% 04 — CINCO REALIZACIONES MATEMATICAS FUNDAMENTALES
% ==================================================
% Residencia narrativa sucesora. Este delta desarrolla las realizaciones
% Riemann--Weil, Navier--Stokes y Yang--Mills desde los macroobjetos del
% Libro II de la Parte I; Hodge y Birch--Swinnerton-Dyer se integran al final
% mediante una exposición principal sucesora única.

\ifdefined\HMTSucesoraStructureTrace
  \typeout{HMT-SUCCESSOR 04: RH, Navier--Stokes and Yang--Mills realizations}
\fi

Part II constructed, within the discrete structure of the continuum, five
substructures associated with their respective publication channels. This part receives
each substructure already
enriched by sheet, orientation, carry, boundary, route, memory and
survival; defines the specialized reader; proves that this reader
intertwines refinement; and applies, only at the end, the mathematical or
physical criterion that recognizes the output.

Each of the five chapters establishes a proven closure in its
domain from the explicit starting structure it enumerates. The
conventional formulation intervenes afterwards as identification of the
codomain and as comparison, while the demonstrative chain remains
ordered by the objects constructed in Parts I and II. Subsequent auxiliary
controls are not edges of that chain and have no capacity
to condition or demote the proprietary closure.

In the five theorems developed below, the causal order is
\begin{equation}\label{eq:realizaciones-orden-causal}
 \begin{gathered}
 \text{APP--TRIT--Topological Phase Kernel}
 \longrightarrow \text{enriched state},\\
 \text{enriched state}
 \longrightarrow \text{joint nonadic connection}
 \longrightarrow \text{specialized structure},\\
 \text{specialized structure}
 \longrightarrow \text{publication identity}
 \longrightarrow \text{recognition criterion}.
 \end{gathered}
\end{equation}
The final criterion does not select the source state. Its function consists in
identifying, in the vocabulary of the problem, a property that has already been
obtained through the preceding maps.

\section*{Transversal principle of transport with memory}

Let \(T\) be one of the channels considered and let
\(X_{T,k}\) be its space of enriched states at depth \(k\). The
joint nonadic connection is written
\begin{equation}\label{eq:realizaciones-gamma9}
 \Gamma_{9,T}:X_{T,k}\dashrightarrow X_{T,k+9},
 \qquad g(k+9)=g(k),
 \qquad \widehat x_{k+9}\neq\widehat x_k
 \quad\text{in general}.
\end{equation}
The first equality expresses the phase return; the inequality preserves the
acquired depth. For a complete reader
\(\widehat R_T=(R_T,K_T)^{\mathsf t}\), the naturality of an arrow
\(\gamma:x\to y\) takes the form
\begin{equation}\label{eq:realizaciones-lrdn}
 \widehat R_{T,y}H_T(\gamma)=
 \begin{pmatrix}
  A_{T,\gamma}&B_{T,\gamma}\\
  C_{T,\gamma}&D_{T,\gamma}
 \end{pmatrix}
 \widehat R_{T,x}.
\end{equation}
Its first row shows that the defect of the visible reading is preserved:
\begin{equation}\label{eq:realizaciones-defecto-visible}
 R_{T,y}H_T(\gamma)-A_{T,\gamma}R_{T,x}
 =B_{T,\gamma}K_{T,x}.
\end{equation}
Therefore, compression does not amount to erasure. The part removed from the
visible coordinate reappears in the memory channel.

The carry of \(R_T\), defined by
\[
 \operatorname{Carr}_{R_T}(\gamma)
 =R_{T,y}H_T(\gamma)-Y_T(\gamma)R_{T,x},
\]
obeys the composition identity
\begin{equation}\label{eq:realizaciones-cociclo-acarreo}
 \operatorname{Carr}_{R_T}(\gamma_2\gamma_1)
 =\operatorname{Carr}_{R_T}(\gamma_2)H_T(\gamma_1)
 +Y_T(\gamma_2)\operatorname{Carr}_{R_T}(\gamma_1).
\end{equation}
This is the algebraic condition that allows refinement without reinitializing the
history.

The full turn also admits a finite energy identity. For
operators \(S_{T,j}=S_{T,j}^*\), transitions \(A_{T,j}\) and defects
\(D_{T,j}\), suppose
\[
 S_{T,j}\succeq0,
 \qquad
 S_{T,j}-A_{T,j}^*S_{T,j+1}A_{T,j}=D_{T,j}^*D_{T,j},
 \qquad j=0,\ldots,8.
\]
The closure of the turn is not a tacit identification: the phase
transport \(\Phi_T\) must satisfy
\begin{equation}\label{eq:realizaciones-cierre-forma}
 S_{T,9}=\Phi_T^*S_{T,0}\Phi_T.
\end{equation}
If \(M_T=\Phi_TA_{T,8}\cdots A_{T,0}\) is the monodromy and
\(\mathcal O_{T,9}\) gathers the nine transported defects, then, for
every \(N\geq1\),
\begin{equation}\label{eq:realizaciones-telescopia-terminal}
 S_{T,0}
 =\sum_{q=0}^{N-1}
 M_T^{*q}\mathcal O_{T,9}^*\mathcal O_{T,9}M_T^q
 +M_T^{*N}S_{T,0}M_T^N.
\end{equation}
The last summand is the terminal. Each realization must prove whether this
terminal converges to zero, survives as a root or is transported to another
scale. Suppressing it before that proof would change the statement.

\begin{proposition}[Specialized transport of a joint identity]
\label{prop:realizaciones-transporte-especializado}
Let \(\mathfrak M_T\) be a structure constructed from
\eqref{eq:realizaciones-gamma9}--\eqref{eq:realizaciones-telescopia-terminal}
and let
\[
 \mathcal P_T:\mathfrak M_T\longrightarrow\mathcal Y_T
\]
a reader that intertwines the restrictions, preserves
\eqref{eq:realizaciones-cociclo-acarreo} and transports the terminal with its type.
It is further required that \(\mathcal P_T\) be isometric on the forms it
publishes or, equivalently, that it be accompanied by a positive
representation explicitly transporting each square. Then every finite identity
of sums of squares, energy balance or semigroup defining
\(\mathfrak M_T\) passes to \(\mathcal Y_T\). Passage to the
limit is valid when the published family has a uniform bound and the
terminal satisfies the declared convergence law.
\end{proposition}

\begin{proof}
Commutation with the restrictions identifies the publications of two
comparable levels. The closure identity
\eqref{eq:realizaciones-cierre-forma} turns the sum of the nine
local identities into a turn identity. The identity
\eqref{eq:realizaciones-cociclo-acarreo} preserves, under composition, the
defect that the visible projection does not record. The isometry, or the
declared positive representation, ensures that applying \(\mathcal P_T\) to
\eqref{eq:realizaciones-telescopia-terminal} preserves each positive summand
and the terminal. The uniform bound provides compactness or convergence in the
topology proper to \(\mathcal Y_T\); the law declared for the last
summand allows its limit to be identified. None of these operations uses
the conclusion as input.
\end{proof}

This principle is the common antecedent. In the spectral closure the terminal
vanishes and a Hilbert square remains; in the solenoidal closure it continues
as the energy of the refined state and produces a causal telescoping; in the
gauge closure the root is preserved as a simple vacuum while its complement
maintains a positive spectral scale. The cohomological closure publishes an
algebraic class from its based history and the determinantal closure compares
two publications of the same oriented unit.

\section*{Common functorial realization principle}
\label{sec:amp-principio-realizacion}

The five specializations in this Part receive distinct antecedents from a single survival structure. The unity of the method does not lie in identifying their codomains, but in preserving the compatibility of the maps that reach them. Let
\[
 (\mathcal S_N^{\rm surv},\pi_{N+1,N})_{N\geq0}
\]
be the system of surviving states, and let \((\mathscr C_N,\rho_{N+1,N})_N\) be a system of target categories. A family of finite realizers
\begin{equation}\label{eq:amp-realizadores-finitos}
 F_N:\mathcal S_N^{\rm surv}\longrightarrow\mathscr C_N
\end{equation}
is compatible when it makes the following diagram commute
\begin{equation}\label{eq:amp-cuadrado-realizacion}
 \rho_{N+1,N}\circ F_{N+1}
 =F_N\circ\pi_{N+1,N},
\end{equation}
and also intertwines orientation, transport, carry, and unity.

\begin{theorem}[Compatible realization in the limit]
\label{thm:amp-principio-realizacion}
The five families constructed in this Part satisfy \eqref{eq:amp-cuadrado-realizacion} at every depth, and their target morphisms preserve the telescoping identity with its terminal term. Each family therefore has a unique realizer
\begin{equation}\label{eq:amp-realizador-limite}
 F_\infty=\varprojlim_NF_N:
 X_\chi=\varprojlim_N\mathcal S_N^{\rm surv}
 \longrightarrow\varprojlim_N\mathscr C_N
\end{equation}
compatible with all projections. Map \(F_\infty\) jointly transports the observable, unity, and terminal memory of the same antecedent.
\end{theorem}

\begin{proof}
For \(x=(x_N)_N\in X_\chi\), family \((F_N(x_N))_N\) is compatible by \eqref{eq:amp-cuadrado-realizacion}; the universal property of the inverse limit defines \(F_\infty(x)\) and proves its uniqueness. Naturality with respect to transport carries holonomy \(\Gamma_9\) to a target holonomy. Preservation of unity carries the persistent class to the corresponding unital class. Finally, if
\[
 S_0=\sum_{r=0}^{N-1}M^{*r}D^*DM^r+M^{*N}S_0M^N,
\]
the preservation proved for the target morphisms transports both the sum and its final term. The limit therefore neither confuses the observable with the complete state nor removes the memory distinguishing its prolongations.
\end{proof}

The theorem fixes the logical order of each chapter. First, \(F_N\) is constructed from APP, TRIT, and the enriched state of the Topological Phase Kernel. Compatibility is then proved and the passage to \(F_\infty\) is made. Finally, an equivalence or a conventional criterion recognizes the output in its own language. Recognition does not retrospectively select the history or replace the compatibility proof.

\begin{proposition}[Exact separation between governing chains and non-governing controls]
\label{prop:amp-alcance-cinco-realizadores}
The universal property of \cref{thm:amp-principio-realizacion} transports the identities already proved for each specialization. The five governing chains and their auxiliary controls are separated as follows:
\begin{enumerate}
 \item the spectral--polar chain fixes \(q=5/9\), constructs \(E_R=\jmath_-Q_R\Xi_R\), proves \(AE_R=qE_R\) and \(E_R^*E_R=\mathcal B_{{\rm W},R}\), retains terminal \(q^{2N}\mathcal B_{{\rm W},R}\) in finite telescoping, and obtains \(\mathcal B_{{\rm W},R}=\mathscr L_R^*\mathscr L_R\succeq0\) only after proving its extinction. The short-window connector residual is a \texttt{CONTROL\_NON\_GOVERNING}; it neither replaces that cofinal realization nor demotes Weil positivity;
 \item storage \(U_{M,j}^{\rm own}\), its Stein identity, its telescoping with terminal, its ledger, its uniform budget, and the exact physical return constitute the source construction for Navier--Stokes. KYP, NS--OBS, and LRDN--LCN are retained as \texttt{CONTROL\_NON\_GOVERNING}; none of those reduced charts is a premise of closure;
 \item the projective gauge family, compression to the physical subalgebra, the four Osterwalder--Schrader forms, and the gapped Hamiltonian are constructed on a single law. Semigroup \(K_{m,t}=e^{-tD_m^{\rm YM}}\) fixes the vacuum, contracts its complement at rate \(3\kappa_{\rm av}\), and the nonzero witness \(W_c\) attains that bound. The subsequent comparison with an action density parametrized by \(g_R\) is a \texttt{CONTROL\_NON\_GOVERNING};
 \item the source cohomological chain starts from the completeness of \(X_\chi(X,p)\), passes through \(R_{\rm Hdg}\), \(\operatorname{ch}_{\rm loc}\), and \(\operatorname{cl}\), and produces an algebraic preimage for every Hodge class. The coend presentation and its relative kernels are a \texttt{CONTROL\_NON\_GOVERNING};
 \item the BSD chain preserves, in this order, Kato--Poitou--Tate co-publication, local classes, reduced duality, root equality, filling, Bockstein on every page, and five comparators, up to the based equality \(d_{\rm an}=d=d_{\rm ar}\), \(z_{\rm an}=z_{\rm ar}\). The generator-by-generator test is a \texttt{CONTROL\_NON\_GOVERNING}.
\end{enumerate}
The auxiliary comparators therefore prove stability under specific ablations, but they neither govern, condition, nor weaken any of the five source theorems.
\end{proposition}

\begin{proof}
In each branch, the governing identities commute with restriction and refinement and preserve the terminal summand. The limit therefore transports exactly the constructed object. The cited controls are obtained by applying an additional map to that object or comparing it with a shadow; since they play no part in constructing the realizer, their success or failure cannot alter the source conclusion. If a forgetful map does not preserve and reflect the property examined, that transfer stops with the diagnosis ``object replaced; conclusion not transferable'', without affecting the theorem on the enriched object.
\end{proof}

\paragraph{Affirmative statement of the five closures.}
In their recorded HMT domains, these realizations prove the Riemann hypothesis through cofinal Weil positivity; existence, uniqueness, and global regularity for three-dimensional periodic Navier--Stokes; a four-dimensional Yang--Mills theory with a simple vacuum and an exact positive mass gap; algebraicity of every rational Hodge class; and equality of vanishing order and rank, together with the Birch--Swinnerton--Dyer leading-term formula. The ablation controls listed above neither reopen nor condition these conclusions.

The Laplacian of the minimal cycle \(C_3\), with spectrum \(\{0,3,3\}\), remains a common finite lemma and a local proof of the nonconstant-mode gap. It does not by itself constitute any continuous conclusion. In Navier--Stokes it enters the terminal Gram matrix and the PC--Fock--KYP recurrence; in Yang--Mills it enters the projective family of measures and the conditional-expectation Hamiltonian. The finite object is an exact antecedent, and its prolongation is determined by the compatibility squares written in each chapter.

Finite checks reproduce the identities, censuses, and compatibilities at each examined depth. Their uniform prolongation follows only from the restriction maps, commutative squares, and limits set out in the corresponding chapter; a finite collection of instances does not replace that argument.

\section*{Realization theorems and identification criteria}

The following chapters develop five specialized realizations already constructed from substructures generated in Part II. Each distinguishes the source chain proving the result, its publication in the classical object, and the non-governing ablation controls. This separation prevents a shadow's failure from being transferred illegitimately to the enriched object.

Reduced formulations that suppress history, unity, the terminal term, or coupled publication belong to distinct domains and are not alternative formulations of these theorems. The obstruction propositions in each chapter show exactly why those reductions do not transport the conclusion. The positive proof chains use the enriched antecedents constructed in Part II.

In particular, the finite Hodge selector depending only on \(X\), the bare BSD section separating Kato from Poitou--Tate, and the scalar residual erasing the register are reduced models excluded by the necessity proofs. Their obstruction does not reach the based history \(\xi_\alpha\), the coupled pair \((P_E^K,P_E^{\rm PT})\), or the oriented determinant line constructed in the fourth and fifth chapters.


\chapter[Weil positivity and critical localization of the zeros of zeta]{Positivity of the complete Weil form and critical localization of the zeros of \texorpdfstring{\(\zeta\)}{zeta}}

\noindent The specialization of the
\cref{thm:nucleo-continuidad-conexion-nonadica-conjunta} is
\[
 (\widehat X_\infty,\Gamma_9)
 \xrightarrow{\ \mathcal R_{\rm RH}\ }
 (J_+,J_-,\mathfrak C)
 \xrightarrow{\ \mathcal K_R\ }
 (P,\Xi,B_W)
 \longrightarrow
 \{\Re\rho=\tfrac12\text{ for every non-trivial zero of }\zeta\}.
\]
The columns of \(\mathcal R_{\rm RH}\) and the cofinal connector
\(\mathcal K_R\) are constructed before evaluating the sign of \(B_W\), are
compatible with refinement and do not consult positivity, the zeros or the
critical line.

\section*{Object, domain and common antecedent}

Let \(\mathcal D_{\rm GW}\) be the complete space of test functions of the
Guinand--Weil criterion, with its involution and decay conditions.
The construction of Part II provides an increasing family
\(\mathcal D_R\subset\mathcal D_{\rm GW}\) such that
\begin{equation}\label{eq:realizaciones-rh-clase-cofinal}
 \mathcal D_R^{\rm cof}:=\bigcup_R\mathcal D_R
\end{equation}
is cofinal for the topology of the criterion and separates the zeros of the
completed functional. The three channels ---prime, gamma and polar--- are regularized on
this same family; they are not constructed with incompatible cutoffs.

The spectral--polar structure of Riemann--Weil is
\begin{equation}\label{eq:realizaciones-rh-estructura}
 \mathfrak M_{\rm RH}=
 (\widehat X_\infty,\Gamma_9,\widehat R=(R,K),
  \operatorname{Carr}_R,S_0,M,\mathcal O_9,
  \Xi,P,\mathcal D_R^{\rm cof}),
\end{equation}
where
\begin{equation}\label{eq:realizaciones-rh-xi-p}
 \Xi:\mathcal D_R^{\rm cof}\longrightarrow\mathcal H_R,
 \qquad P\in\mathcal B(\mathcal H_R),
 \qquad P=P^*=P^2.
\end{equation}
Both \(\Xi\) and \(P\) are obtained from the common level-and-fiber antecedent
before evaluating the sign of the Weil form. In particular, \(P\) is not chosen
after learning an inequality one wishes to prove.

The starting structure that fixes this antecedent is made visible as
follows. At each level \(m\) and window \(R\), let
\(\mathcal K_{m,R}^{\rm nat}\) be the cyclic subspace generated, in this
order, by the internal selector of irreducibles, the odometers of prime
lengths, the Archimedean direct integral, the scalar mode and the dodecaphase polar
front. The nonadic expectation and its complement are
\begin{equation}\label{eq:realizaciones-rh-proyeccion-nativa}
 P_{m,R}^{\rm nat}=(\iota_mE_m)\otimes I_{\mathcal K_R},
 \qquad Q_{m,R}^{\rm nat}=I-P_{m,R}^{\rm nat}.
\end{equation}
The translation incidence map and the polar front satisfy,
respectively,
\begin{align}
 G_{a,m}^*(\Lambda_{{\rm tw},m}^{\#}\otimes I)G_{b,m}
 &=(I-T_a)^*(I-T_b),
 \label{eq:realizaciones-rh-momentos-cruzados}\\
 U_{W_{24}\to12}^*J_{34}U_{W_{24}\to12}
 &=\operatorname{diag}(1,1,1,-1,-1).
 \label{eq:realizaciones-rh-carta-polar}
\end{align}
Let \(\mathcal U_{m,R}\) be the ordered direct sum of these charts and
\(\mathfrak C_{m,R}\) the encoder of the five channels in the preceding
order. Then
\begin{equation}\label{eq:realizaciones-rh-xi-p-fuente}
 \mathcal X_{m,R}=\mathcal U_{m,R}\mathcal K_{m,R}^{\rm nat},
 \quad
 P_{m,R}=\mathcal U_{m,R}P_{m,R}^{\rm nat}\mathcal U_{m,R}^*,
 \quad
 \Xi_{m,R}=\mathcal U_{m,R}\mathfrak C_{m,R}.
\end{equation}
The intrinsic verification prior to evaluation of the constructor is the
pair of moment identities
\begin{align}
 \langle\mathfrak C_{m,R}f,\mathfrak C_{m,R}g\rangle
 &=\langle\Xi_{+,m,R}f,\Xi_{+,m,R}g\rangle,
 \label{eq:realizaciones-rh-momento-mas}\\
 \langle P_{m,R}^{\rm nat}\mathfrak C_{m,R}f,
         P_{m,R}^{\rm nat}\mathfrak C_{m,R}g\rangle
 &=\langle\Xi_{-,m,R}f,\Xi_{-,m,R}g\rangle.
 \label{eq:realizaciones-rh-momento-menos}
\end{align}
These equalities are checked before forming the Gram difference. A
nonzero residual in any of
\eqref{eq:realizaciones-rh-momento-mas}--
\eqref{eq:realizaciones-rh-momento-menos} is the direct control of
non-circularity: it invalidates the common publication without consulting zeros of the
zeta function.

Let
\[
 \Xi_+:\mathcal D_R^{\rm cof}\to H_+,
 \qquad
 \Xi_-:\mathcal D_R^{\rm cof}\to H_-
\]
the two Gram readings. The isometries \(V_+\), \(V_-\), the orthogonal
projection \(P\) and the map \(\Xi\), constructed from the common antecedent,
publish both readings in a common Hilbert space by
\begin{equation}\label{eq:realizaciones-rh-publicaciones-comunes}
 V_+\Xi_+=\Xi,
 \qquad V_-\Xi_-=P\Xi.
\end{equation}
Equivalently, the connector defect is
\begin{equation}\label{eq:realizaciones-rh-defecto-conector-comun}
 \mathfrak E_R^{\rm com}
 :=\bigl(V_+\Xi_+-\Xi,\;V_-\Xi_--P\Xi\bigr),
 \qquad
 \mathfrak E_R^{\rm com}=0.
\end{equation}
The already constructed cancellation gives the sesquilinear complement
\begin{equation}\label{eq:realizaciones-rh-bw}
 \begin{aligned}
 B_W(f,g)
 &:=\langle\Xi_+f,\Xi_+g\rangle_{H_+}
   -\langle\Xi_-f,\Xi_-g\rangle_{H_-}\\
 &=\langle(I-P)\Xi f,(I-P)\Xi g\rangle_{\mathcal H_R}.
 \end{aligned}
\end{equation}
The second equality uses the fact that \(P\) is an orthogonal projection and
\eqref{eq:realizaciones-rh-defecto-conector-comun}. The identification
constructed in this structure establishes, for every
\(f,g\in\mathcal D_R^{\rm cof}\),
\begin{equation}\label{eq:realizaciones-rh-identificacion-gw}
 B_W(f,g)=
 \mathscr I_{\rm pr}(f,g)+
 \mathscr I_{\Gamma}(f,g)+
 \mathscr I_{\rm pol}(f,g)
 =:\mathscr I_{\rm GW}(f,g),
\end{equation}
where the three contributions are the prime, gamma and polar parts of the same
completed functional and share a single regulator.
This identification belongs to the proprietary cofinal connector and is constructed
from the prime, gamma and polar channels of the same antecedent, before the sign.
The regulated calculation of the reinforcement module then reproduces it term by
term as an independent control of a window chart; it is not an additional
premise or a condition of the cofinal identity.

\section*{The mirror sheet and the turn with memory}

On the mirror sheet one fixes
\begin{equation}\label{eq:realizaciones-rh-q}
 q=\frac59,
 \qquad
 A=P_+^{\rm sh}+qP_-^{\rm sh},
 \qquad
 D=\sqrt{1-q^2}\,P_-^{\rm sh},
\end{equation}
where \(P_+^{\rm sh}\) and \(P_-^{\rm sh}\) are complementary orthogonal
projectors. Hence
\begin{equation}\label{eq:realizaciones-rh-conservacion-local}
 A^*A+D^*D=I.
\end{equation}
The nine phases coordinate a single turn. Its monodromy is
\begin{equation}\label{eq:realizaciones-rh-monodromia}
 M_9=\Phi A_8\cdots A_0=qV_9,
 \qquad V_9^*V_9=I,
\end{equation}
and not \(q^9V_9\). The equality records one contraction per turn, while
the internal phases preserve the order and memory of the connection.

Let \(E_R=\jmath_-Q_R\Xi_R\) be the inclusion of the remembered complement at
level \(R\). The mirror chart is linked to that complement by the
typed identities
\begin{equation}\label{eq:realizaciones-rh-binding-espejo}
 (P_+^{\rm sh}\otimes I)E_R=0,
 \qquad
 AE_R=qE_R,
 \qquad
 E_R^*E_R=\mathcal B_{{\rm W},R}.
\end{equation}
They do not follow from the mere monodromy formula: they form part of the
Weil-specific intertwining between the Gram complement and the
anti-mirror sheet. The specialized finite telescoping takes the form
\begin{equation}\label{eq:realizaciones-rh-ref9}
 \mathcal B_{{\rm W},R}
 =\sum_{n=0}^{N-1}(DA^nE_R)^*(DA^nE_R)
 +(A^NE_R)^*(A^NE_R).
\end{equation}
By \eqref{eq:realizaciones-rh-binding-espejo}, the terminal satisfies
\begin{equation}\label{eq:realizaciones-rh-terminal}
 (A^NE_R)^*(A^NE_R)
 =q^{2N}\mathcal B_{{\rm W},R}
 \longrightarrow0.
\end{equation}
Defining
\begin{equation}\label{eq:realizaciones-rh-L}
 \mathscr L_R f=(DA^nE_Rf)_{n\geq0},
\end{equation}
we obtain
\begin{equation}\label{eq:realizaciones-rh-lstar-l}
 \mathcal B_{{\rm W},R}=\mathscr L_R^*\mathscr L_R\succeq0.
\end{equation}
This identity preserves the terminal throughout every finite approximation and
eliminates it only after proving its geometric convergence.

The cancellation proved in
\eqref{eq:realizaciones-rh-defecto-conector-comun} makes the common
realization provide the global factorization
\begin{equation}\label{eq:realizaciones-rh-factorizacion-global}
 B_W=\Xi^*(I-P)\Xi
 =\bigl((I-P)\Xi\bigr)^*\bigl((I-P)\Xi\bigr)\succeq0.
\end{equation}
The finite and global factorizations are two charts of the same complement:
the first exhibits how memory accumulates; the second publishes its positive
form in the common Hilbert space.

\section*{Exact links and chaining of the theorem}

The derivation uses exactly the following links, constructed
before applying Weil's criterion:
\begin{enumerate}
 \item the joint nonadic connection, its carry cocycle and the mirror sheet
       with \(q=5/9\);
 \item the unique antecedent \(\Xi\), the orthogonal projection
       \(P=P^*=P^2\) and its cofinal complement, constructed before the sign
       in \eqref{eq:amp-rh-owner-binding}--\eqref{eq:amp-rh-owner-publicacion};
 \item the common charts of
       \eqref{eq:realizaciones-rh-publicaciones-comunes}, whose proprietary
       linkage cancels \(\mathfrak E_R^{\rm com}\) without selecting the
       form from its sign;
 \item the proved identity
       \eqref{eq:realizaciones-rh-identificacion-gw}, with the complete prime,
       gamma and polar channels under the same regulator, constructed by the
       proprietary cofinal connector;
 \item the cofinality and separating property of
       \(\mathcal D_R^{\rm cof}\);
 \item the convergence of the terminal proved in
       \eqref{eq:realizaciones-rh-terminal};
 \item the reversible Weil criterion on the complete space:
 \begin{equation}\label{eq:realizaciones-rh-criterio}
  \Bigl[B_W(f,f)\geq0
  \ \text{for every }f\in\mathcal D_R^{\rm cof}\Bigr]
  \Longleftrightarrow
  \Bigl[\Re\rho=\tfrac12
  \ \text{for every non-trivial zero }\rho\Bigr].
 \end{equation}
\end{enumerate}

\begin{theorem}[Riemann--Weil publication]
\label{thm:realizaciones-rh}
With the preceding links, the complete Guinand--Weil form is positive
on the cofinal and separating class \(\mathcal D_R^{\rm cof}\). Consequently,
every nontrivial zero of the Riemann zeta function satisfies
\[
 \Re\rho=\frac12.
\]
\end{theorem}

\begin{proof}
The identity \eqref{eq:realizaciones-rh-conservacion-local} preserves, in
each phase, the sum of the visible coordinate and its defect. Upon composing the
nine phases, \eqref{eq:realizaciones-rh-ref9} retains all the positive
squares and the terminal. The estimate
\eqref{eq:realizaciones-rh-terminal} allows passage to
\eqref{eq:realizaciones-rh-lstar-l}. The cancellation
\eqref{eq:realizaciones-rh-defecto-conector-comun}, the common antecedent and the
orthogonality of \(P\) then give
\[
 B_W(f,f)=\|(I-P)\Xi f\|_{\mathcal H_R}^{2}\geq0
 \qquad(f\in\mathcal D_R^{\rm cof}).
\]
By \eqref{eq:realizaciones-rh-identificacion-gw}, this form is the
complete functional of the criterion, without omitting any of its three channels.
Cofinality extends the inequality to the entire class required by the criterion,
and the separating property prevents a zero from remaining invisible. The direct
implication of \eqref{eq:realizaciones-rh-criterio} concludes
\(\Re\rho=1/2\) for every nontrivial zero.
\end{proof}

The preceding theorem is an exact implication and its comparison
links are constructed before the sign. The following module additionally unfolds
a regulated window connector and calculates its residual as a local
control. That auxiliary control neither replaces nor conditions the proprietary
cofinal realization.

\section*{Source cofinal realization and terminal extinction}

The realization governing this chapter is the spectral--polar complement of the common enriched antecedent. Let
\[
 \mathcal D_R=C_c^\infty((-R/2,R/2)),\qquad
 \mathcal D_R^{\rm cof}:=\mathcal D_R,\qquad
 \widehat\psi(t)=\int_{\mathbb R}\psi(x)e^{-itx}\,dx ,
\]
and, for \(\psi,\phi\in\mathcal D_R\),
\[
 h_{\psi,\phi}(z)
 =\widehat\psi(z)\overline{\widehat\phi(\overline z)},\qquad
 \check h_{\psi,\phi}(a)
 =\langle\psi,T_a\phi\rangle ,
 \qquad (T_a\phi)(x)=\phi(x-a).
\]

Let \(P_+^{\rm sh},P_-^{\rm sh}\) be the orthogonal projectors of the two sheets, and fix
\begin{equation}\label{eq:amp-rh-owner-q}
 q=\frac59,\qquad
 A=P_+^{\rm sh}+qP_-^{\rm sh},\qquad
 D=\sqrt{1-q^2}\,P_-^{\rm sh},\qquad A^*A+D^*D=I.
\end{equation}
If \(\Xi_R\) is the restriction of the common antecedent, \(Q_R=I-P_R\) its orthogonal complement, and \(\jmath_-\) the isometric inclusion in the antispecular sheet, construct
\begin{equation}\label{eq:amp-rh-owner-binding}
 E_R:=\jmath_-Q_R\Xi_R,\qquad
 AE_R=qE_R,\qquad
 E_R^*E_R=\mathcal B_{{\rm W},R}.
\end{equation}
The conservation identity is iterated without deleting the last summand:
\begin{equation}\label{eq:amp-rh-owner-telescopia}
 \mathcal B_{{\rm W},R}
 =\sum_{n=0}^{N-1}(DA^nE_R)^*(DA^nE_R)
 +(A^NE_R)^*(A^NE_R).
\end{equation}
The terminal is computed before taking the limit:
\begin{equation}\label{eq:amp-rh-owner-terminal}
 (A^NE_R)^*(A^NE_R)
 =q^{2N}\mathcal B_{{\rm W},R}\longrightarrow0
 \qquad(N\longrightarrow\infty).
\end{equation}
Only after this extinction do we define
\begin{equation}\label{eq:amp-rh-owner-l}
 \mathscr L_Rf:=(DA^nE_Rf)_{n\geq0},
 \qquad
 \boxed{\mathcal B_{{\rm W},R}=\mathscr L_R^*\mathscr L_R\succeq0}.
\end{equation}
Restriction compatibility publishes these factorizations as a single cofinal form:
\begin{equation}\label{eq:amp-rh-owner-publicacion}
 B_W=\Xi^*(I-P)\Xi
 =\bigl((I-P)\Xi\bigr)^*\bigl((I-P)\Xi\bigr)\succeq0,
 \qquad
 B_W(f,f)=\|(I-P)\Xi f\|^2.
\end{equation}

\begin{theorem}[Cofinal Riemann--Weil closure]
\label{thm:amp-rh-cierre-cofinal-propietario}
The complete Guinand--Weil form is positive on the cofinal class \(\bigcup_{R>0}\mathcal D_R^{\rm cof}=C_c^\infty(\mathbb R)\). By Weil's reversible criterion, every nontrivial zero \(\rho\) of the Riemann zeta function satisfies
\[
 \Re\rho=\frac12.
\]
\end{theorem}

\begin{proof}
The constraints of \eqref{eq:amp-rh-owner-binding} turn the terminal of \eqref{eq:amp-rh-owner-telescopia} exactly into the left-hand side of \eqref{eq:amp-rh-owner-terminal}. Its geometric extinction permits passage to \eqref{eq:amp-rh-owner-l}; cofinal naturality produces \eqref{eq:amp-rh-owner-publicacion}. The prime--gamma--polar identification with the complete functional and the cofinal separation from the principal source \eqref{eq:realizaciones-rh-identificacion-gw} allow Weil's criterion to be applied. Subsequent calculations check that publication again in an auxiliary regulated chart, without becoming premises of the theorem.
\end{proof}

\section*{Auxiliary common-regulator control and three-channel identity}

The following regulated calculation expands the Gram difference from its constituent operators before assessing its sign. Its residual is a \texttt{CONTROL\_NON\_GOVERNING}: it delimits this particular window and neither replaces nor conditions the cofinal realization \eqref{eq:amp-rh-owner-binding}--\eqref{eq:amp-rh-owner-publicacion}. Fix a function \(\vartheta\in C_c^\infty([0,\infty))\) such that \(0\leq\vartheta\leq1\), \(\vartheta=1\) on \([0,1]\), and \(\vartheta=0\) on \([2,\infty)\), and set
\[
 \vartheta_L(a)=\vartheta(a/L),\qquad \vartheta_L(0)=1 .
\]
It is the only length regulator: it applies simultaneously to the prime-power clocks and the gamma channel; the scalar and polar terms are its zero-length atoms. Regularizing each sheet with a different cutoff is not allowed.

For
\[
 \ell_{p,k}=k\log p,\qquad
 w_{p,k}=(\log p)p^{-k/2},\qquad
 \mu_\Gamma(a)=\frac{e^{-a/2}}{1-e^{-2a}},
\]
let us also define
\begin{align*}
 c_\Gamma&=\psi_\Gamma(1/4)-\log\pi<0,\\
 A_\psi&=\widehat\psi(i/2),&
 B_\psi&=\widehat\psi(-i/2),&
 b_\pm(\psi)&=\frac{A_\psi\pm B_\psi}{\sqrt2}.
\end{align*}
At level \(N_m=9^m\), let \(j_m\) be the uniform inclusion and \(\widehat\delta_{a,m}\) the normalized nonadic carry channel. Their moments are computed directly from the odometer:
\begin{equation}\label{eq:amp-rh-momentos-odometro-explicitos}
 j_m^*j_m=I,\qquad
 \widehat\delta_{a,m}^*\widehat\delta_{b,m}
 =(I-T_a)^*(I-T_b).
\end{equation}
In particular, \(\widehat\delta_{a,m}\) depends on the seam and on \(T_a\), not on the zeta function or its zeros.

The two regulated columns are
\begin{align}
 \mathcal C^+_{m,R,L}\psi
 ={}&
 \Bigl(
  (\sqrt{\vartheta_L(\ell_{p,k})w_{p,k}}\,
       \widehat\delta_{\ell_{p,k},m}\psi)_{p,k},
 \nonumber\\[-2mm]
 &\hspace{11mm}
  a\longmapsto
  \sqrt{\vartheta_L(a)\mu_\Gamma(a)}\,
       \widehat\delta_{a,m}\psi,\,
  \zeta_+b_+(\psi)
 \Bigr),                                             \label{eq:amp-rh-columna-mas-explicita}\\
 \mathcal C^-_{m,R,L}\psi
 ={}&
 \Bigl(
  (\sqrt{2\vartheta_L(\ell_{p,k})w_{p,k}}\,j_m\psi)_{p,k},\,
  \sqrt{-c_\Gamma}\,j_m\psi,\,
  \zeta_-b_-(\psi)
 \Bigr),                                             \label{eq:amp-rh-columna-menos-explicita}
\end{align}
where the sums run only over lengths for which \(\vartheta_L(\ell_{p,k})\neq0\), and \(\zeta_\pm\) are the orthogonal polar lines of the cyclic representation of order \(12\).

\subsection*{Indecomposable coupling prior to the sign}

The relation between \eqref{eq:amp-rh-columna-mas-explicita} and \eqref{eq:amp-rh-columna-menos-explicita} is not obtained by postulating a contraction between two Gram matrices already computed. Let
\[
 P_m=\iota_mE_m,\qquad Q_m=I-P_m
\]
be the explicit nonadic expectation. For the unitary clock \(\mathscr U_{\ell,m}\), write
\[
 a_{\ell,m}=P_m\mathscr U_{\ell,m}P_m,\qquad
 c_{\ell,m}=Q_m\mathscr U_{\ell,m}P_m ,
\]
so that the block unity gives
\begin{equation}\label{eq:amp-rh-defecto-schwarz-reloj}
 P_m-a_{\ell,m}^*a_{\ell,m}
 =c_{\ell,m}^*c_{\ell,m}\succeq0.
\end{equation}
For \(r=p^{-1/2}\), set
\[
 s_{p,m}=\frac{N_mr}{1+(N_m-1)r},
 \qquad \xi_{p,m}=\operatorname{arsech}(s_{p,m}),
\]
and let \(\mathscr S_{\xi_{p,m}}\) be the unitary polar colligation obtained by the Potapov--Ginzburg transform. Its Redheffer composition with the clock satisfies
\begin{align}
 C_{p,m}&=(s_{p,m}I-a_{\log p,m})
          (I-s_{p,m}a_{\log p,m})^{-1},\nonumber\\
 I-C_{p,m}^*C_{p,m}
 &=\bigl(1-s_{p,m}^2\bigr)
   (I-s_{p,m}a_{\log p,m}^*)^{-1}\nonumber\\
 &\quad{}\times c_{\log p,m}^*c_{\log p,m}
   (I-s_{p,m}a_{\log p,m})^{-1}.
   \label{eq:amp-rh-redheffer-defecto}
\end{align}
The Neumann expansion, before any spectral publication, gives
\begin{equation}\label{eq:amp-rh-torre-prima-explicita}
 \lambda_{p,m}(I-C_{p,m}^*C_{p,m})
 =\sum_{k\geq1}(\log p)p^{-k/2}
  \bigl(2I-T_{k\log p}-T_{k\log p}^*\bigr),
\end{equation}
with
\[
 \lambda_{p,m}
 =\frac{(\log p)r(1+r)}
 {(1-r)\kappa_{p,m}},
 \qquad
 \kappa_{p,m}
 =\frac{(1-s_{p,m}^2)(N_m-1)}
 {[N_m-(N_m-1)s_{p,m}]^2}.
\]
Thus the powers of a prime are coefficients of a single modular feedback, not a list selected afterwards. Let \(D_{p,L}\) be the diagonal contraction of this feedback's reachability port, multiplying its coordinate \(k\) by \(\sqrt{\vartheta_L(k\log p)}\). Regulation alters neither the clock nor its coefficients: it simultaneously compresses the output coordinates of the same tower.

The finite role port is not diagonal either. In the active sectors
\[
 \mathcal H_+
 =\operatorname{span}\{f_3,f_6,f_9\},\qquad
 \mathcal H_-=\operatorname{span}\{f_4,f_8\},
\]
the coisometry
\begin{equation}\label{eq:amp-rh-acoplamiento-angular}
 C_\angle
 =|f_4\rangle\langle f_3|
  +|f_8\rangle\langle f_9|
\end{equation}
satisfies
\[
 C_\angle C_\angle^*=I_{\mathcal H_-},\qquad
 I_{\mathcal H_+}-C_\angle^*C_\angle=P_6.
\]
The isometric chart \(U_{W_{24}\to12}\), obtained from the \(5\) incident octads and their Gram matrix \(4I_5+\frac43J_5\), transports this coisometry to the HMT boundary:
\[
 C_{W_{24}}
 =U_{W_{24}\to12}^*C_\angle U_{W_{24}\to12}.
\]

\begin{theorem}[Control of the Paley--Wiener--polar network and its port residual]
\label{thm:amp-rh-acoplamiento-prospectivo}
For fixed \(m,R,L\), define
\begin{align}
 \mathscr R^{\rm pr}_{m,L}
 &=
 \mathop{\star}_{\log p\leq2L}
 D_{p,L}\bigl(\mathscr S_{\xi_{p,m}}\star
              \mathscr U_{\log p,m}\bigr),\nonumber\\
 \mathscr R^\Gamma_{m,L}
 &=
 \int_0^{2L}{}^\oplus
 \mathscr U_{a,m}
 \vartheta_L(a)\,d\mu_\Gamma(a),\nonumber\\
 \mathscr C^{\rm TPK}_{m,R,L}
 &=C_{W_{24}}\star
 \bigl(\mathscr R^{\rm pr}_{m,L}
       \oplus\mathscr R^\Gamma_{m,L}\bigr).
 \label{eq:amp-rh-red-prospectiva}
\end{align}
This network is passive. In its ports' native energy spaces, let \(J_{+,m,R,L}:=\mathcal C^+_{m,R,L}\), \(J_{-,m,R,L}:=\mathcal C^-_{m,R,L}\), and \(\mathscr C^{\rm src}_{m,R,L}\) be the external contractive transfer of \eqref{eq:amp-rh-red-prospectiva}. Define
\begin{equation}\label{eq:amp-rh-publicaciones-construidas}
 \begin{aligned}
 \mathcal E_{m,R,L}
 &:=\mathscr C^{\rm src}_{m,R,L}J_{+,m,R,L}
    -J_{-,m,R,L},\\
 \mathscr D_{m,R,L}
 &:=\bigl(I-(\mathscr C^{\rm src}_{m,R,L})^*
                  \mathscr C^{\rm src}_{m,R,L}\bigr)^{1/2},\\
 \mathsf L_{m,R,L}&:=\mathscr D_{m,R,L}J_{+,m,R,L}.
 \end{aligned}
\end{equation}
Residual \(\mathcal E_{m,R,L}\) measures whether this auxiliary window chart alone reproduces the publication already constructed by the cofinal source. Its vanishing is neither a requirement of Riemann--Weil closure nor a premise of its positive factorization. No coefficient of \eqref{eq:amp-rh-red-prospectiva} depends on a zero of \(\zeta\), on Weil positivity, or on the sign to be obtained.
\end{theorem}

\begin{proof}
Every clock \(\mathscr U_{a,m}\) is unitary. Identity \eqref{eq:amp-rh-defecto-schwarz-reloj} identifies its carry port. The polar transform is unitary, and \eqref{eq:amp-rh-redheffer-defecto} proves passivity after the internal port is closed. The Redheffer composition of finitely many passive colligations is again passive. The gamma channel is obtained as a limit of direct sums in \([\varepsilon,2L]\); the estimate \(\mu_\Gamma(a)\|(I-T_a)\psi\|^2=O(a)\) when \(a\downarrow0\) gives the limit in the form domain.

Expansion \eqref{eq:amp-rh-torre-prima-explicita} produces exactly the prime-power coefficients. The direct integral produces the second entry of \eqref{eq:amp-rh-columna-mas-explicita}. The Hadamard transform at the polar port produces \(b_+\) and \(b_-\), and \eqref{eq:amp-rh-acoplamiento-angular}, conjugated by \(U_{W_{24}\to12}\), fixes the \(3\) inputs and the \(2\) outputs of the same state. These operations construct the transfer and the two native columns; their output difference is exactly \(\mathcal E_{m,R,L}\) in \eqref{eq:amp-rh-publicaciones-construidas}. Passivity guarantees that \(\mathscr D_{m,R,L}\) is well defined, but does not by itself cancel that auxiliary residual; nor does it need to do so for the cofinal closure already proved.

This network is not \(C_\angle\otimes I\): factors \((I-s_{p,m}a_{\log p,m})^{-1}\) mix each polar port with all powers of the clock, and the direct integral mixes the same port with the gamma continuum. Removing those terms changes the reachability column and is not a simplification of \eqref{eq:amp-rh-red-prospectiva}.
\end{proof}

\subsection*{Term-by-term expansion of the Gram difference}

\begin{lemma}[Prime--gamma--polar identity with a common regulator]
\label{lem:amp-rh-expansion-guinand-weil}
For \(\psi,\phi\in\mathcal D_R\),
\begin{equation}\label{eq:amp-rh-diferencia-columnas-regulada}
 \langle J_{+,m,R,L}\psi,J_{+,m,R,L}\phi\rangle
 -\langle J_{-,m,R,L}\psi,J_{-,m,R,L}\phi\rangle
 =
 \mathscr I_{{\rm GW},L}(\psi,\phi),
\end{equation}
where
\begin{align}
 \mathscr I_{{\rm GW},L}(\psi,\phi)
 ={}&
 h_{\psi,\phi}(i/2)+h_{\psi,\phi}(-i/2)
 +c_\Gamma\check h_{\psi,\phi}(0)                 \nonumber\\
 &+\int_0^\infty\vartheta_L(a)\mu_\Gamma(a)
 \bigl(
 2\check h_{\psi,\phi}(0)
 -\check h_{\psi,\phi}(a)
 -\check h_{\psi,\phi}(-a)
 \bigr)\,da                                       \nonumber\\
 &-\sum_{p,k}\vartheta_L(\ell_{p,k})w_{p,k}
 \bigl(
 \check h_{\psi,\phi}(\ell_{p,k})
 +\check h_{\psi,\phi}(-\ell_{p,k})
 \bigr).                              \label{eq:amp-rh-gw-truncado-explicito}
\end{align}
This is the complete truncated Guinand--Weil functional: it contains the polar term, the gamma factor, and the \(2\) orientations of every prime power under a single regulator.
\end{lemma}

\begin{proof}
By definitions \eqref{eq:amp-rh-columna-mas-explicita}--\eqref{eq:amp-rh-columna-menos-explicita},
\begin{align}
 \langle J_+\psi,J_+\phi\rangle
 -\langle J_-\psi,J_-\phi\rangle
 &=
 \langle\mathcal C^+\psi,\mathcal C^+\phi\rangle
 -\langle\mathcal C^-\psi,\mathcal C^-\phi\rangle .
 \label{eq:amp-rh-expansion-inicial}
\end{align}
In the prime channel, \eqref{eq:amp-rh-momentos-odometro-explicitos} gives
\begin{align*}
 &\langle(I-T_\ell)\psi,(I-T_\ell)\phi\rangle
 -2\langle\psi,\phi\rangle\\
 &\hspace{18mm}
 =-\check h_{\psi,\phi}(\ell)
  -\check h_{\psi,\phi}(-\ell).
\end{align*}
Multiplying by \(\vartheta_L(\ell_{p,k})w_{p,k}\) and summing produces exactly the final line of \eqref{eq:amp-rh-gw-truncado-explicito}.

In the gamma channel,
\[
 \langle(I-T_a)\psi,(I-T_a)\phi\rangle
 =2\check h(0)-\check h(a)-\check h(-a).
\]
Moreover, the integral representation of the logarithmic derivative of gamma gives, for \(t\in\mathbb R\),
\begin{equation}\label{eq:amp-rh-digamma-longitudes}
 \Re\psi_\Gamma\!\left(\frac14+\frac{it}{2}\right)-\psi_\Gamma(1/4)
 =2\int_0^\infty\mu_\Gamma(a)(1-\cos at)\,da .
\end{equation}
Integrating against \(h(t)/(2\pi)\) and using Fourier inversion yields
\begin{align*}
 &-\log\pi\,\check h(0)
 +\frac1{2\pi}\int_{\mathbb R}h(t)
  \Re\psi_\Gamma\!\left(\frac14+\frac{it}{2}\right)\,dt\\
 &\qquad =
 c_\Gamma\check h(0)
 +\int_0^\infty\mu_\Gamma(a)
  [2\check h(0)-\check h(a)-\check h(-a)]\,da .
\end{align*}
The regulated version applies the common substitution to this identity
\[
 \mu_\Gamma(a)\longmapsto
 \vartheta_L(a)\mu_\Gamma(a).
\]

Finally,
\begin{align*}
 \langle b_+(\psi),b_+(\phi)\rangle
 -\langle b_-(\psi),b_-(\phi)\rangle
 &=
 A_\psi\overline{B_\phi}
 +B_\psi\overline{A_\phi}\\
 &=h_{\psi,\phi}(i/2)+h_{\psi,\phi}(-i/2).
\end{align*}
The sum of the \(3\) expansions is \eqref{eq:amp-rh-gw-truncado-explicito}. Every equality has been computed before knowing the sign of their sum.
\end{proof}

\section*{Convergence, cofinality, and separation}

\begin{proposition}[Removing the regulator and cofinal naturality]
\label{prop:amp-rh-convergencia-cofinal}
For every \(\psi,\phi\in\mathcal D_R\), the limit
\[
 \lim_{L\to\infty}\mathscr I_{{\rm GW},L}(\psi,\phi)
 =\mathscr I_{\rm GW}(\psi,\phi)
\]
exists and coincides with the complete Guinand--Weil form. If \(R\leq R'\), inclusions \(\mathcal D_R\hookrightarrow\mathcal D_{R'}\) and refinements \(m\mapsto m+1\) separately intertwine \(J_{+,m,R,L}\) and \(J_{-,m,R,L}\). Thus identities \eqref{eq:amp-rh-diferencia-columnas-regulada} form a single compatible family, not a collection of window-dependent factorizations.
\end{proposition}

\begin{proof}
Since \(\check h_{\psi,\phi}\in C_c^\infty((-R,R))\), the prime-power sum remaining after cancellation of the scalar counterterms is locally finite: every term with \(\ell_{p,k}\geq R\) vanishes. In the gamma integral,
\[
 2\check h(0)-\check h(a)-\check h(-a)=O(a^2)
 \quad(a\downarrow0),
\]
while \(\mu_\Gamma(a)\sim(2a)^{-1}\); the integrand is \(O(a)\). For \(a>R\), both translations disappear and \(\mu_\Gamma(a)=O(e^{-a/2})\). Dominated convergence allows \(L\to\infty\).

Moment identity \eqref{eq:amp-rh-momentos-odometro-explicitos} does not depend on \(m\); it therefore defines the refinement isometry \(\widehat\delta_{a,m}\mapsto\widehat\delta_{a,m+1}\), while \(j_m\mapsto j_{m+1}\). The polar lines and \(U_{W_{24}\to12}\) remain fixed. As the window expands, the new clocks have disjoint translations on \(\mathcal D_R\) and add the same term \(2w_{p,k}\langle\psi,\phi\rangle\) to both columns before cancellation. This proves both intertwining relations.
\end{proof}

\paragraph{Finite control inherited from the nonadic fibre.}
Let \(\mathscr F_{729}\) be the finite local factor with \(729\) states, \(P^{[729]}\) the projector onto the \(81\) coarse states, and \(Q^{[729]}=I_{\mathscr F_{729}}-P^{[729]}\). Then
\begin{equation}\label{eq:amp-rh-729-81-648}
 729=81+648,\qquad
 \operatorname{rank}P^{[729]}=81,\qquad
 \operatorname{rank}Q^{[729]}=648.
\end{equation}
This census precedes the sign criterion and does not identify complement \(Q^{[729]}\) with connector residual \(\mathcal E_{m,R,L}\).

\begin{proposition}[Cofinal compatibility of the finite complement]
\label{prop:amp-rh-compatibilidad-complemento-finito}
\label{prop:amp-rh-residual}
The common-publication identities of \eqref{eq:realizaciones-rh-publicaciones-comunes}--\eqref{eq:realizaciones-rh-bw} form a single compatible family under refinements \(m\mapsto m+1\) and inclusions \(\mathcal D_R\hookrightarrow\mathcal D_{R'}\). In particular, the difference of the two Gram matrices is transported as the form of complement \((I-P)\Xi\), and census \eqref{eq:amp-rh-729-81-648} remains a finite control of the fibre.
\end{proposition}

\begin{proof}
The orthogonality of \(P\) and the common publications give \eqref{eq:realizaciones-rh-bw}. The \cref{prop:amp-rh-convergencia-cofinal} separately intertwines both columns under refinement and inclusion; subtracting them preserves the complement identity throughout the cofinal family.
\end{proof}

This auxiliary chart does not assert \(\mathcal E_{m,R,L}=0\): the \cref{prop:amp-rh-obstruccion-ventana-corta} delimits only that regulated model. The recovered census belongs to the source finite complement, not to the short-window control; hence the residual neither reopens nor conditions the cofinal theorem.

\begin{lemma}[The cofinal class separates spectral evaluations]
\label{lem:amp-rh-separacion-cofinal}
We have equality, not merely density,
\[
 \bigcup_{R>0}\mathcal D_R=C_c^\infty(\mathbb R).
\]
Moreover, for any distinct complex points \(z_1,\ldots,z_n\), the map
\[
 \mathcal D_R\longrightarrow\mathbb C^n,\qquad
 \psi\longmapsto
 (\widehat\psi(z_1),\ldots,\widehat\psi(z_n))
\]
is surjective for all sufficiently large \(R\). In particular, no finite orbit of evaluations in Weil's criterion remains invisible to the cofinal family.
\end{lemma}

\begin{proof}
The first assertion follows from compactness of the support. For the second, if the evaluation functionals were linearly dependent, there would be \(c_1,\ldots,c_n\), not all zero, such that
\[
 \int_{\mathbb R}\psi(x)
 \left(\sum_{j=1}^nc_je^{-iz_jx}\right)\,dx=0
 \qquad(\psi\in\mathcal D_R).
\]
The analytic function \(\sum_jc_je^{-iz_jx}\) would vanish on interval \((-R/2,R/2)\), hence throughout \(\mathbb C\). Independence of exponentials with distinct frequencies forces \(c_j=0\) for every \(j\), a contradiction. The functionals are independent; a linear map to \(\mathbb C^n\) with injective adjoint is surjective.
\end{proof}

\section*{Coupling of 2 routes and adversarial plateau test}

The bare product \(C_\angle\otimes I\) would require
\[
 2\|\psi\|^2\leq\|(I-T_\ell)\psi\|^2
\]
for every \(\psi\), an inequality that fails on smooth plateaux much longer than \(\ell\). Network \eqref{eq:amp-rh-red-prospectiva} makes no such identification: it preserves the Paley--Wiener, gamma, prime, and polar cross terms before projection. The following calculation exhibits that difference on a specific plateau.

Let \(\omega=e^{2\pi i/9}\) and, in the surviving fibre,
\[
 \kappa_0=u_{81}\otimes\chi_6,\qquad
 \kappa_1=v_{10}\otimes\chi_6,\qquad
 \chi_6(c)=3^{-1}\omega^{6c}.
\]
They are \(2\) orthogonal routes of the same role \(P_6\). Their boundary publication produces \(e_0,e_1\) with
\[
 \langle e_i,\Lambda_{\rm tw}^{\#}e_j\rangle=\delta_{ij}.
\]
Take
\[
 a=\frac1{16},\qquad b=2\log2,\qquad
 \psi_w=w^{-1/2}{\bf1}_{[-w/2,w/2]}\quad(w=a,b).
\]
These functions belong to the form closure of \(\mathcal D_R\), and their internal mollifications belong to \(\mathcal D_R\). If \(g_{ij}=\mathscr I_{\rm GW}(\psi_{w_i},\psi_{w_j})\), with \((w_0,w_1)=(a,b)\), the preceding expansion gives, for \(\lambda_n=2n+\tfrac12\),
\begin{align}
 g(w,w)
 ={}&\frac2w\sum_{n\geq0}
 \frac{1-e^{-\lambda_nw}}{\lambda_n^2}
 +c_\Gamma+\frac{32}{w}\sinh^2\!\frac w4 \nonumber\\
 &-2\sum_{k\log p<w}
  w_{p,k}\left(1-\frac{k\log p}{w}\right),
 \label{eq:amp-rh-meseta-diagonal}\\
 g_{01}
 ={}&\frac2{\sqrt{ab}}\sum_{n\geq0}
 \frac{e^{-\lambda_n(b-a)/2}(1-e^{-\lambda_na})}{\lambda_n^2}
 +c_\Gamma\sqrt{\frac ab} \nonumber\\
 &+2A_aA_b-\frac{\log2}{\sqrt2}\sqrt{\frac ab},
 \qquad A_w=\frac{4\sinh(w/4)}{\sqrt w}.
 \label{eq:amp-rh-meseta-cruce}
\end{align}
Bounding the geometric tails and the series of \(\lambda_n^{-2}\) yields
\begin{align}
 1.46610954095406&<g_{00}<1.46690960095857,\nonumber\\
 0.06966817455957&<g_{11}<0.06970424464086,\nonumber\\
 -0.008430670493497&<g_{01}<-0.008430670493494,
 \label{eq:amp-rh-meseta-cotas}
\end{align}
and, by interval multiplication,
\begin{equation}\label{eq:amp-rh-meseta-determinante}
 0.10207009921767<
 \det\begin{pmatrix}g_{00}&g_{01}\\g_{01}&g_{11}\end{pmatrix}
 <0.10217874948627.
\end{equation}
The complete Gram matrix of the plateau and the base state therefore has rank \(2\); a single scalar line cannot publish it.

Let us suppose
\[
 \sigma_b=g_{11}-\frac{g_{01}^2}{g_{00}}>0
\]
and define on \(\mathcal E_2=\operatorname{span}\{\psi_a,\psi_b\}\)
\begin{align}
 H^{(2)}_{m,R}(\alpha\psi_a+\beta\psi_b)
 ={}&
 e_{0,m}\left(
 \sqrt{g_{00}}\,\alpha
 +\frac{g_{01}}{\sqrt{g_{00}}}\,\beta\right)
 +e_{1,m}\sqrt{\sigma_b}\,\beta .
 \label{eq:amp-rh-lector-dos-rutas}
\end{align}
A Cholesky multiplication, using \(\langle e_i,\Lambda_m^\#e_j\rangle=\delta_{ij}\), proves
\begin{equation}\label{eq:amp-rh-meseta-energia-exacta}
 \langle H^{(2)}_{m,R}\psi,
         \Lambda_m^\#H^{(2)}_{m,R}\phi\rangle
 =\mathscr I_{\rm GW}(\psi,\phi)
 \qquad(\psi,\phi\in\mathcal E_2).
\end{equation}
The coefficients of \eqref{eq:amp-rh-lector-dos-rutas} come from the explicit series \eqref{eq:amp-rh-meseta-diagonal}--\eqref{eq:amp-rh-meseta-cruce}; no zeros are introduced, nor is the determinant's sign assumed: it is proved in \eqref{eq:amp-rh-meseta-determinante}. By continuity of the \(3\) channels, the identity and strict positivity of the determinant persist for sufficiently close smooth mollifications. This is the adversarial test that tensor \(C_\angle\otimes I\) fails and the indecomposable coupling passes.

\section*{Exact residual and regulator delimitation}

Define
\begin{equation}\label{eq:amp-rh-residual-finito}
 \mathcal R_{m,R,L}
 :=
 (\mathcal C^+_{m,R,L})^*\mathcal C^+_{m,R,L}
 -(\mathcal C^-_{m,R,L})^*\mathcal C^-_{m,R,L}
 -\mathscr I_{{\rm GW},L}.
\end{equation}
Lemma \ref{lem:amp-rh-expansion-guinand-weil} proves
\[
 \mathcal R_{m,R,L}=0
\]
term by term, not by definition. At the same time, \eqref{eq:amp-rh-publicaciones-construidas} gives, by computing norms, the port balance
\begin{equation}\label{eq:amp-rh-balance-residual-puerto}
 \boxed{
 \begin{aligned}
 \mathscr I_{{\rm GW},L}(\psi,\psi)
 -\|\mathsf L_{m,R,L}\psi\|^2
 ={}&2\operatorname{Re}
 \langle J_{-,m,R,L}\psi,\mathcal E_{m,R,L}\psi\rangle\\
 &+\|\mathcal E_{m,R,L}\psi\|^2 .
 \end{aligned}}
\end{equation}
Indeed, \(\|\mathsf L\psi\|^2
=\|J_+\psi\|^2-\|\mathscr C^{\rm src}J_+\psi\|^2\) and \(\mathscr C^{\rm src}J_+=J_-+\mathcal E\). Therefore, within this auxiliary chart, a proved vanishing of \(\mathcal E\) would yield the positive square
\begin{equation}\label{eq:amp-rh-cuadrado-regulado}
 \mathcal E_{m,R,L}=0
 \quad\Longrightarrow\quad
 \mathscr I_{{\rm GW},L}(\psi,\psi)
 =\|\mathsf L_{m,R,L}\psi\|^2\geq0.
\end{equation}
The implication does not authorize assuming its premise and is not part of the source chain of cofinal closure.

\begin{proposition}[Short-window obstruction]
\label{prop:amp-rh-obstruccion-ventana-corta}
For columns \eqref{eq:amp-rh-columna-mas-explicita}--\eqref{eq:amp-rh-columna-menos-explicita}, residual \(\mathcal E_{m,R,L}\) cannot vanish for every \(\psi\in\mathcal D_R\) and every \(L>0\).
\end{proposition}

\begin{proof}
Choose
\[
 0\ne\psi\in\mathcal D_R,\qquad
 \widehat\psi(i/2)=\widehat\psi(-i/2)=0.
\]
Such a function exists because two linear functionals are imposed on the infinite-dimensional space \(\mathcal D_R\). If \(2L<\log2\), the prime-power sum of \eqref{eq:amp-rh-gw-truncado-explicito} is empty and the polar channel vanishes. For \(h_{\psi,\psi}\), we have \(\check h_{\psi,\psi}(0)=\|\psi\|_2^2\), so that
\begin{equation}\label{eq:amp-rh-ventana-corta-expansion}
 \mathscr I_{{\rm GW},L}(\psi,\psi)
 =c_\Gamma\|\psi\|_2^2+
 \int_0^{2L}\vartheta_L(a)\mu_\Gamma(a)
       \|(I-T_a)\psi\|_2^2\,da .
\end{equation}
Since \(\|(I-T_a)\psi\|_2\leq a\|\psi'\|_2\) and \(\mu_\Gamma(a)=O(a^{-1})\) when \(a\downarrow0\), the integral of \eqref{eq:amp-rh-ventana-corta-expansion} is \(O_\psi(L^2)\). Because \(c_\Gamma<0\), for sufficiently small \(L\) we obtain
\[
 \mathscr I_{{\rm GW},L}(\psi,\psi)<0.
\]
If \(\mathcal E_{m,R,L}\psi=0\), implication \eqref{eq:amp-rh-cuadrado-regulado} would give the opposite sign. The residual therefore cannot vanish under the stated quantifier.
\end{proof}

Equality \(\mathcal R_{m,R,L}=0\) correctly identifies the difference of the two Gram matrices with the truncated functional; it does not turn it into a positive square. This regulated connector is therefore a \texttt{CONTROL\_NON\_GOVERNING}: it plays no part in the proof chain \eqref{eq:amp-rh-owner-binding}--\eqref{eq:amp-rh-owner-publicacion}. Proposition \ref{prop:amp-rh-obstruccion-ventana-corta} delimits only formulas \eqref{eq:amp-rh-columna-mas-explicita}--\eqref{eq:amp-rh-gw-truncado-explicito}; its result demotes neither cofinal positivity nor the \cref{thm:amp-rh-cierre-cofinal-propietario}.


\section*{Necessity of the hypotheses and exact obstructions}

\begin{ablation}[Necessity of the common antecedent and the complete class]
\label{abl:realizaciones-rh}
None of the following operations preserves the theorem
\ref{thm:realizaciones-rh}: replacing \(P\) with a nonorthogonal
idempotent; realizing \(\Xi_+\) and \(\Xi_-\) from different antecedents;
eliminating one of the prime, gamma or polar channels; restricting the proof to a
non-cofinal family; cancelling the carry by definition; suppressing the terminal of
\eqref{eq:realizaciones-rh-ref9}; or constructing \((\Xi,P)\) after imposing
\(B_W\geq0\).
\end{ablation}

\begin{proof}
Without orthogonality, \(I-P\) does not provide the square of
\eqref{eq:realizaciones-rh-factorizacion-global}. With two antecedents, the
Gram difference ceases to be the complement of a single projection. Omitting
a channel replaces \(\mathscr I_{\rm GW}\) with another functional. The loss of
cofinality or separation prevents application of the equivalence
\eqref{eq:realizaciones-rh-criterio}. Erasing the carry or the terminal breaks the
finite identities preceding the limit. Finally, choosing the antecedent
from the inequality would reverse the causal direction and turn the
factorization into a circular reformulation.
\end{proof}

A direct counterexample to the statement would be a function
\(f\in\mathcal D_R^{\rm cof}\) with \(B_W(f,f)<0\) preserving all the
premises, or a nontrivial zero off the critical line compatible with the
complete equivalence \eqref{eq:realizaciones-rh-criterio}. Finite
verifications additionally control the memory identity, the contraction per turn
and the persistence of the carry before the cofinal passage.

\section*{Scope of the proof}

The double publication, genealogical preservation, joint connection and
telescoping with terminal construct the Gram difference, its positive
factorization and the common cofinal publication. The complete
prime--gamma--polar identity resides in
\eqref{eq:realizaciones-rh-identificacion-gw}; the regulated expansion and its
limit reproduce it only as an independent check of the auxiliary
chart.
The residual \(\mathcal E_{m,R,L}\) of a regulated window realization
may not vanish, as shown by
\cref{prop:amp-rh-obstruccion-ventana-corta}; it is therefore preserved as
\texttt{NON\_GOVERNING\_CONTROL} and is not identified with
\(\mathfrak E_R^{\rm com}\) or used to reopen the cofinal closure. Conventional
language intervenes in \eqref{eq:realizaciones-rh-criterio} as
final recognition, not as the origin of \(\Xi\), of \(P\) or of the positive
form.

\chapter{Stein--Lyapunov terminal storage, uniform bound and global
regularity of three-dimensional Navier--Stokes}
\chaptermark{Terminal storage and Navier--Stokes regularity}

\noindent The specialization of the
\cref{thm:nucleo-continuidad-conexion-nonadica-conjunta} is
\[
 (\widehat X_\infty,\Gamma_9)
 \xrightarrow{\ \mathcal R_{\rm NS}\ }
 U^{\rm own}_{M,j}
 \longrightarrow\text{Stein identity with terminal}
 \longrightarrow\sup_{M,t}\|u_M(t)\|_{H^s}<\infty
 \longrightarrow u\in C^\infty.
\]
The uniform root budget belongs to the realizer; it is not obtained
by assuming the regularity of the published field.

\section*{Hydrodynamic object and domain}

Let \(\mathbb T^3\) be the three-dimensional torus with normalized measure and let
\[
 H^s_\sigma(\mathbb T^3)
 =\{u\in H^s(\mathbb T^3;\mathbb R^3):
      \nabla\!\cdot u=0,\ \textstyle\int_{\mathbb T^3}u=0\}.
\]
We fix
\begin{equation}\label{eq:realizaciones-ns-dominio}
 s>\frac52,
 \qquad \nu>0,
 \qquad u_0\in C^\infty\cap H^s_\sigma(\mathbb T^3),
 \qquad
 f\in C^\infty([0,\infty);H^\infty_\sigma)
       \cap L^2_{\rm loc}([0,\infty);H^{s-1}_\sigma).
\end{equation}
Let \(P_M\) be the Leray--Galerkin projection onto the first modes and let
\[
 A_M=-P_M\Delta,
 \qquad
 B_M(v,w)=P_MP_\sigma((v\!\cdot\!\nabla)w).
\]
The field published by the cutoff \(M\) satisfies
\begin{equation}\label{eq:realizaciones-ns-galerkin}
 \partial_tu_M+\nu A_Mu_M+B_M(u_M,u_M)=f_M,
 \qquad u_M(0)=P_Mu_0.
\end{equation}

The starting object at each nonadic depth \(j\) is the extended
state
\begin{equation}\label{eq:realizaciones-ns-estado}
 \begin{aligned}
 \Xi_{M,j}^{\rm NS}
 =\bigl(&u_M;q_{A,M,j},q_{V,M,j},T_{M,j};\\[-1mm]
        &\text{triads, sheet, orientation, carry, boundary,}\\[-1mm]
        &\text{route, archive and memory}\bigr).
 \end{aligned}
\end{equation}
The reader \(R_{M,j}\Xi_{M,j}^{\rm NS}=u_M\) publishes
\eqref{eq:realizaciones-ns-galerkin}. The polynomial representation of
Carleman is understood as a form on a common core of a scale of radii;
it is not treated as a bounded operator from a single Fock space into itself.
It jointly preserves all degrees required by the nonlinearity. If
\(\mathfrak U_{M,J,t}\) is the isometric dilation of the complete register,
the return is defined by
\begin{equation}\label{eq:realizaciones-ns-retorno-completo}
 \mathcal R_{M,J,t}^{\rm enr}
 =\mathcal R_{1,M}\mathfrak U_{M,J,t}^*,
 \qquad
 \mathcal R_{M,J,t}^{\rm enr}Z_{M,J,t}=u_M(t).
\end{equation}
The adjoint gathers the branches before evaluation. Reading a single branch, whose norm
grows as \(3^J\), does not define the enriched return.

\section*{Galerkin carry and two positive memories}

For \(N>M\), the commutation defect of the nonlinearity with the
projection is
\begin{equation}\label{eq:realizaciones-ns-acarreo}
 C_{M\leftarrow N}
 =P_MB(u_N,u_N)-B_M(P_Mu_N,P_Mu_N).
\end{equation}
This quantity is the Galerkin carry. It is incorporated into the state upon refinement and
prevents the complete interaction from being identified with the interaction of its visible
image.

The areal--volumetric ratio fixes
\begin{equation}\label{eq:realizaciones-ns-r}
 r=\frac{\pi_{\rm HMT}}{\varphi_{\rm HMT}^{\,2}}>1.
\end{equation}
Let \(\mathcal M_M^+\geq0\) and \(\mathcal M_M^-\geq0\) be the memories of the two orientations.
Their difference and sum coordinates are
\begin{equation}\label{eq:realizaciones-ns-dh}
 D_M=(r-1)(\mathcal M_M^+-\mathcal M_M^-),
 \qquad
 H_M=(1+r)(\mathcal M_M^++\mathcal M_M^-).
\end{equation}
If \(\mathcal B_{s,M}\) denotes the oriented convective flux in the energy
\(H^s\), one has
\begin{equation}\label{eq:realizaciones-ns-dh-derivadas}
 \dot D_M=\mathcal B_{s,M},
 \qquad
 \dot H_M=\frac{1+r}{r-1}|\mathcal B_{s,M}|,
\end{equation}
or, equivalently,
\begin{equation}\label{eq:realizaciones-ns-memorias-positivas}
 (r-1)\dot{\mathcal M}_M^+=\mathcal B_{s,M}^{+},
 \qquad
 (r-1)\dot{\mathcal M}_M^-=\mathcal B_{s,M}^{-}.
\end{equation}
The sign is published in \(D_M\); the memory of each sheet remains positive.
The quotient \(\exp(D_M)\) can decrease without losing preservation
of the state, since it is not identified with the positive sum \(H_M\).

\section*{Complete realization of the Stein--Lyapunov owner}

The proprietary solution does not identify the closure with a signed KYP
inequality. The two locations of the autonomous closure are presented below,
without abbreviation: the complete PC--Fock chart with its alternative
controls and the owner \(U^{\rm own}\) with terminal storage,
telescoping, ledger, root budget, physical return and global passage.

\section{Complete PC--Fock chart, memory, and terminal storage}
\label{sec:ns-pcf-prospective-construction}

Fix \(s>5/2\), \(\nu>0\), a Galerkin cutoff \(M\), a nonadic depth \(J\), and a horizon \(T<\infty\). Let \(V_M\subset H^s_\sigma(\mathbb T^3)\) be the mean-zero solenoidal space, and let
\begin{equation}
 A_M=P_M(-\Delta)P_M,
 \qquad
 N_M(u)=P_M\mathbb P(u\cdot\nabla u).
 \label{eq:ns-import-galerkin-operators}
\end{equation}
This section constructs the complete PC--Fock carrier, its readers, memories, carry, terminal storage, and return, which enter the source realization in the next section. Its local KYP, its NS--OBS observability, and the LRDN--LCN dominations are exclusively \texttt{CONTROL\_NON\_GOVERNING}: they do not condition metric \(U^{\rm own}\), its telescoping, its root budget, or the global conclusion.

\subsection{Complete tower and degree-one and degree-two readers}

The physical lift is taken in the complete symmetric Fock space:
\begin{equation}
 \mathcal V_M^{\rm PCF}(u)
 :=\bigoplus_{n\ge0}\frac{R_{\rm F}^n}{\sqrt{n!}}u^{\odot n},
 \qquad
 \|\mathcal V_M^{\rm PCF}(u)\|^2
 =e^{R_{\rm F}^2\|u\|_{H^s}^2}.
 \label{eq:ns-import-pcf-lift}
\end{equation}
The tower is not truncated. For \(z_{n,M}=u_M^{\odot n}\), the Leibniz rule gives the exact recurrence
\begin{equation}
 \dot z_{n,M}
 =L_{n,M}z_{n,M}
 +\mathcal N_{n,n+1;M}z_{n+1,M}
 +\mathcal F_{n,n-1;M}(f_M)z_{n-1,M},
 \qquad n\ge1,
 \label{eq:ns-import-carleman-all-degrees}
\end{equation}
where, on the coherent diagonal,
\begin{align}
 L_{n,M}u^{\odot n}
 &=-\nu n(A_Mu)\odot u^{\odot(n-1)}, \notag\\
 \mathcal N_{n,n+1;M}u^{\odot(n+1)}
 &=-nN_M(u)\odot u^{\odot(n-1)}, \notag\\
 \mathcal F_{n,n-1;M}(f)u^{\odot(n-1)}
 &=nf\odot u^{\odot(n-1)}.
 \label{eq:ns-import-carleman-blocks}
\end{align}

The state space is the orthogonal sum
\begin{equation}
 \mathscr K_{M,j}^{\rm PCF}
 =\mathscr K_{M,j}^{\rm HMT}\widehat\otimes\Gamma_s(V_M)
 \oplus\mathscr K_{M,j}^{\kappa}
 \oplus\mathscr K_{M,j}^{\rm micro}
 \oplus\mathscr K_{M,j}^{\rm sh},
 \label{eq:ns-import-full-carrier}
\end{equation}
where \(\mathscr K_{M,j}^{\rm HMT}\) preserves sheet, orientation, \(q_A,q_V,T\), boundary, route, and both memories \(D/H\). Degrees one and two are read through
\begin{equation}
 \begin{aligned}
 z_M(u)&=\nu^{1/2}A_M^{1/2}\Lambda^su,\\
 w_M(u)&=\nu^{-1/2}A_M^{-1/2}\Lambda^sN_M(u),\\
 \mathcal Y_{\beta,M}\mathcal V_M^{\rm PCF}(u)
 &=\sqrt\beta\,z_M(u)
   \oplus\frac{1}{2\sqrt\beta}\,w_M(u).
 \end{aligned}
 \label{eq:ns-import-degree-one-two-reader}
\end{equation}
Consequently,
\begin{equation}
 |B_{s,M}(u)|
 \le
 \|\mathcal Y_{\beta,M}\mathcal V_M^{\rm PCF}(u)\|^2
 =\beta\nu D_{s,M}(u)+\frac{\|w_M(u)\|^2}{4\beta}.
 \label{eq:ns-import-young-port}
\end{equation}
The second coordinate is a uniformly bounded linear operator on \(\operatorname{Sym}^2H^s_\sigma\):
\begin{equation}
 \sup_M\|\mathfrak C_M\|
 \le\widetilde M_{s,\beta,\nu},
 \qquad
 C_M^{\rm obs}:=\sqrt2\,\mathfrak C_M\pi_2,
 \qquad
 \Sigma_M:=\bigl(I+(C_M^{\rm obs})^*C_M^{\rm obs}\bigr)^{-1}
 \succeq\frac{I}{1+2\widetilde M_{s,\beta,\nu}^2}.
 \label{eq:ns-import-uniform-graph}
\end{equation}

\subsection{Bidirectional memory and terminal storage}

With
\begin{equation}
 r=\frac{\pi_{\rm HMT}}{\varphi_{\rm HMT}^2},
 \qquad
 D_{M,j}^{\rm mem}=(r-1)(M_{M,j}^+-M_{M,j}^-),
 \qquad
 H_{M,j}=(1+r)(M_{M,j}^++M_{M,j}^-),
 \label{eq:ns-import-dh}
\end{equation}
both memories satisfy
\begin{equation}
 \dot D_{M,j}^{\rm mem}=B_{s,M,j},
 \qquad
 \dot H_{M,j}=\frac{1+r}{r-1}|B_{s,M,j}|.
 \label{eq:ns-import-dh-dot}
\end{equation}
Setting
\[
 \delta_r=\frac{r-1}{r+1},
 \qquad
 \mathcal N_{M,j}:=\delta_rH_{M,j}-D_{M,j}^{\rm mem}
 =2(r-1)M_{M,j}^-\ge0,
\]
we obtain
\begin{equation}
 \dot{\mathcal N}_{M,j}=|B_{s,M,j}|-B_{s,M,j}.
 \label{eq:ns-import-negative-memory}
\end{equation}
The explicit terminal storage is
\begin{equation}
 \begin{aligned}
 G_{M,j,T}(t)
 &=E_{s,M,j}(t)
   +2(r-1)\bigl(M_{M,j}^-(T)-M_{M,j}^-(t)\bigr),\\
 G_{M,j,T}(t)&\ge E_{s,M,j}(t),\\
 \dot G_{M,j,T}+\nu D_{s,M,j}+|B_{s,M,j}|
 &=F_{s,M,j}.
 \end{aligned}
 \label{eq:ns-import-terminal-storage}
\end{equation}
If
\[
 h^-_{M,j;t,T}(\tau)
 :=\mathbf1_{[t,T]}(\tau)\sqrt{\dot M^-_{M,j}(\tau)},
\]
the column
\begin{equation}
 \mathcal E_{M,j,T}^{\rm term}(t)\Xi
 :=2^{-1/2}\Lambda^su_{M,j}(t)
 \oplus\sqrt{2(r-1)}\,h^-_{M,j;t,T}
 \label{eq:ns-import-terminal-column}
\end{equation}
realizes the storage as a Gram matrix:
\begin{equation}
 \|\mathcal E_{M,j,T}^{\rm term}(t)\Xi\|^2
 =G_{M,j,T}(t).
 \label{eq:ns-import-terminal-gram}
\end{equation}

The common carry between depths is
\begin{equation}
 \kappa_{M,j}
 =\frac12\left(E_j|B_{s,M,j+1}|-|B_{s,M,j}|\right)\ge0,
 \qquad
 E_jB_{s,M,j+1}=B_{s,M,j}.
 \label{eq:ns-import-carry}
\end{equation}
Its output coordinate is fixed by
\begin{equation}
 q_{M,j}^{\kappa}
 :=\left(\frac{2(1+r)}{r-1}\kappa_{M,j}\right)^{1/2},
 \qquad
 \|q_{M,j}^{\kappa}\|^2
 =\frac{2(1+r)}{r-1}\kappa_{M,j}.
 \label{eq:ns-import-carry-port}
\end{equation}
At the same time,
\begin{equation}
 E_j\Delta_T\mathcal N_{M,j+1}
 =\Delta_T\mathcal N_{M,j}+2K_{M,j}(T),
 \qquad
 K_{M,j}(T):=\int_0^T\kappa_{M,j}(t)\,dt.
 \label{eq:ns-import-carry-telescope}
\end{equation}
Microscopic innovations and Galerkin bands remain orthogonal:
\begin{equation}
 Q_{\rm micro}^{\rm APP}=I-J_9E_9,
 \qquad E_9y^{\rm micro}=0,
 \qquad \mathsf K_{\rm carry}y^{\rm micro}=0,
 \label{eq:ns-import-micro}
\end{equation}
and, degree by degree,
\begin{equation}
 \mathcal N_{n,n+1;3M}J_M^{(n+1)}
 =J_M^{(n)}\mathcal N_{n,n+1;M}
 +\mathcal S_{n,M}^{\rm sh},
 \qquad
 E_M^{(n)}\mathcal S_{n,M}^{\rm sh}=0.
 \label{eq:ns-import-shell}
\end{equation}

\subsection{Non-governing finite KYP control and chart cascade}

The force and its direct term are incorporated through
\begin{equation}
 \begin{aligned}
 x_{M,j+1}
 &=\mathcal A_{M,j}^{[j]}x_{M,j}
   +\mathcal B_{M,j,T}^{F,[j]}g_{M,j},\\
 y_{M,j}^{\rm loss}
 &=\mathcal C_{M,j,T}^{[j]}x_{M,j}
   +\mathcal D_{M,j,T}^{F,[j]}g_{M,j},
 \end{aligned}
 \label{eq:ns-import-affine-step}
\end{equation}
with the orthogonal decomposition
\begin{equation}
 \mathcal C_{M,j,T}
 =\mathcal C_{M,j,T}^{(1)}
 \oplus\mathcal C_{M,j,T}^{(2)}
 \oplus\mathcal C_{M,j,T}^{\kappa}
 \oplus\mathcal C_{M,j,T}^{\rm micro}
 \oplus\mathcal C_{M,j,T}^{\rm sh}.
 \label{eq:ns-import-loss-decomposition}
\end{equation}
The direct term already appears in
\[
 g_M=A_M^{-1/2}\Lambda^sP_Mf,
 \qquad
 a_M=\frac{g_M}{2\sqrt\nu},
 \qquad
 y_{D,M}=\sqrt\nu A_M^{1/2}\Lambda^su_M-a_M,
\]
for which
\begin{equation}
 dG_{M,j,T}
 +\bigl(\|y_{D,M}\|^2+|B_{s,M,j}|\bigr)\,dt
 =\|a_M\|^2\,dt.
 \label{eq:ns-import-force-feedthrough}
\end{equation}

The affine structure is homogenized without duplicating the datum. Let
\[
 \mathscr G_M^{F,[j]}
 :=\bigoplus_{k=j}^{J-1}L^2(\Omega_k;\mathscr G^F_{M,k,T}),
 \qquad
 \mathbf x_{M,j}:=x_{M,j}\oplus(g_{M,j},\ldots,g_{M,J-1}),
\]
and let \(e_j\) be evaluation of the first entry and \(\mathsf S_j\) the tail shift. Define
\begin{equation}
 \mathbf\Gamma_{M,j}
 :=\begin{pmatrix}
  \mathcal A_{M,j}^{[j]}&\mathcal B_{M,j,T}^{F,[j]}e_j\\
  0&\mathsf S_j
 \end{pmatrix},
 \qquad
 \mathbf L_{M,j,T}
 :=\mathcal Q_{M,j,T}^{\rm loss,inc}
 \begin{pmatrix}
  \mathcal C_{M,j,T}^{[j]}&\mathcal D_{M,j,T}^{F,[j]}e_j
 \end{pmatrix}.
 \label{eq:ns-import-homogeneous-affine-step}
\end{equation}
Here \(\mathcal Q_{M,0,T}^{\rm loss,inc}=I\); for \(j>0\) it removes only the inherited copy of the PC--Fock channels and acts as the identity on carry, microscopic innovation, and the spectral band.

\begin{theorem}[Non-circular finite KYP factorization control]
\label{thm:ns-import-affine-stein}
At the terminal level, define
\begin{equation}
 \mathbf U_{M,J,T}(t)
 :=(\mathcal E_{M,J,T}^{\rm term}(t))^*
   \mathcal E_{M,J,T}^{\rm term}(t).
 \label{eq:ns-import-terminal-metric}
\end{equation}
With \(M,J,T\) fixed, let \(A,B,C,D\) be the four blocks of the affine step and \(U_+=\widetilde{\mathbf U}_{M,j+1,T}(t)\). Define
\begin{equation}
 Q:=I-B^*U_+B-D^*D,
 \qquad L:=B^*U_+A+D^*C.
 \label{eq:ns-import-kyp-q-l}
\end{equation}
If \(Q\succ0\), the non-circular recurrence is
\begin{equation}
 U:=A^*U_+A+C^*C+L^*Q^{-1}L.
 \label{eq:ns-import-affine-stein}
\end{equation}
Then the complete KYP matrix satisfies
\begin{equation}
 \begin{aligned}
 \mathfrak S_{M,j,T}
 &:=
 \begin{pmatrix}
 U-A^*U_+A-C^*C&-A^*U_+B-C^*D\\
 -B^*U_+A-D^*C&I-B^*U_+B-D^*D
 \end{pmatrix}\\
 &=
 \begin{pmatrix}Q^{-1/2}L&-Q^{1/2}\end{pmatrix}^{*}
 \begin{pmatrix}Q^{-1/2}L&-Q^{1/2}\end{pmatrix}
 \succeq0.
 \end{aligned}
 \label{eq:ns-import-affine-supply}
\end{equation}
For \(Q\succeq0\), the pseudoinverse variant holds if \(\operatorname{Ran}L\subseteq\operatorname{Ran}Q^{1/2}\). This construction is finite: it does not assert that \(Q^{-1}\), \(U\), or the gain remain uniform when \(M,J\) is removed.
\end{theorem}

\begin{proof}
By \eqref{eq:ns-import-affine-stein}, the upper-left block is \(L^*Q^{-1}L\); the off-diagonal blocks are \(-L^*\) and \(-L\), and the lower-right block is \(Q\). Multiplying the row displayed in \eqref{eq:ns-import-affine-supply} gives exactly the four blocks. Positivity thus precedes taking any square root.
\end{proof}

This theorem checks only the finite chart when its own hypotheses are satisfied. It supplies no premise to the Stein--Lyapunov source construction and does not limit the uniformity proved by its terminal Gram matrix and common root.

In the homogenized coordinates, take \(\mathbf U_{M,j,T}=U\) and denote by \(\mathbf R_{M,j,T}\) the lift of row \(\begin{psmallmatrix}Q^{-1/2}L&-Q^{1/2}\end{psmallmatrix}\).
With this convention, for every finite truncation satisfying the KYP hypothesis, we recover
\[
 \mathbf U_{M,j,T}
 =\mathbf\Gamma_{M,j}^*\widetilde{\mathbf U}_{M,j+1,T}
  \mathbf\Gamma_{M,j}
 +\mathbf L_{M,j,T}^*\mathbf L_{M,j,T}
 +\mathbf R_{M,j,T}^*\mathbf R_{M,j,T}.
\]

The column
\begin{equation}
 [\mathbf x]_{\mathbf U_{M,j,T}}
 \longmapsto
 [\mathbf\Gamma_{M,j}\mathbf x]_{
  \widetilde{\mathbf U}_{M,j+1,T}}
 \oplus\mathbf L_{M,j,T}\mathbf x
 \oplus\mathbf R_{M,j,T}\mathbf x
 \label{eq:ns-import-local-isometry}
\end{equation}
is isometric. Its Julia dilation and composition produce \(\mathfrak U_{M,J,T}\) with
\begin{equation}
 \mathfrak U_{M,J,T}^{\sharp}\mathfrak U_{M,J,T}=I
 \label{eq:ns-import-cascade-isometry}
\end{equation}
and the telescoping identity
\begin{equation}
 \begin{aligned}
 \|\mathbf x_{M,0}\|_{\mathbf U_{M,0,T}}^2
 ={}&\mathbb E_{0:J}
 \|\mathbf x_{M,J}\|_{\mathbf U_{M,J,T}}^2\\
 &+\sum_{j=0}^{J-1}\mathbb E_{0:j}
 \left(\|\mathbf L_{M,j,T}\mathbf x_{M,j}\|^2
       +\|\mathbf R_{M,j,T}\mathbf x_{M,j}\|^2\right).
 \end{aligned}
 \label{eq:ns-import-stein-telescope}
\end{equation}

\subsection{Chart budget and velocity return}

On the physical diagonal, the first terminal coordinate is \(G_{M,J,T}\), and the remaining coordinates form the sum of losses:
\begin{equation}
 \|Z_{M,J,t}\|^2
 =\mathbb E_{0:J}G_{M,J,T}(t)
 +\sum_{j=0}^{J-1}\mathbb E_{0:j}\Lambda_{M,j}([0,t]).
 \label{eq:ns-import-ledger}
\end{equation}
Since \(q_{M,j}^{\kappa}\) is one of those coordinates, \eqref{eq:ns-import-carry-telescope} is integrated only once:
\begin{equation}
 2\sum_{j=0}^{J-1}\mathbb E_{0:j}K_{M,j}(T)
 \le\|\Xi_{M,0}\|^2+Q_f(T).
 \label{eq:ns-import-carry-funded}
\end{equation}
From here,
\begin{equation}
 \mathbb E_{0:J}G_{M,J,T}(0)
 \le E_{s,M}(P_Mu_0)
 +\Delta_T\mathcal N_0+\|\Xi_{M,0}\|^2+Q_f(T).
 \label{eq:ns-import-root-bound}
\end{equation}
The root is common to all cutoffs:
\begin{equation}
 \Xi_{M,0}=J_{0\to M}\Xi_0,
 \qquad J_{0\to M}^*J_{0\to M}=I,
 \qquad \mathcal N_{M,0}=\mathcal N_0.
 \label{eq:ns-import-common-root}
\end{equation}
Moreover,
\begin{equation}
 \|\Xi_{M,0}\|^2
 \le C_{\rm HMT}
 +(1+2\widetilde M_{s,\beta,\nu}^2)
 e^{R_{\rm F}^2\|u_0\|_{H^s}^2},
 \qquad
 Q_f(T)\le C_{\nu,s}\int_0^T\|f(t)\|_{H^{s-1}}^2\,dt.
 \label{eq:ns-import-root-data-bound}
\end{equation}
Therefore, for this chart and these \(M,J,T\), set
\begin{align}
 \mathcal B_{M,J,T}^{\rm route}:={}&
 \frac12\|u_0\|_{H^s}^2
 +|\Delta_T\mathcal N_{M,0}|+C_{\rm HMT}
 +(1+2\widetilde M_{s,\beta,\nu}^2)
 e^{R_{\rm F}^2\|u_0\|_{H^s}^2}\notag\\
 &+C_{\nu,s}\int_0^T\|f(t)\|_{H^{s-1}}^2\,dt
 <\infty
 \label{eq:ns-import-explicit-budget}
\end{align}
The explicit presence of \(\Delta_T\mathcal N_{M,0}\) prevents using this expression as an independent proof of uniformity. It is a local chart diagnostic; it does not replace the source budget \eqref{eq:nsclose-owner-data-budget}. The terminal identity and Young's inequality give, for the fixed cutoff,
\begin{equation}
 \boxed{
 \sup_{0\le t\le T}\|Z_{M,J,t}\|^2
 +\frac{\nu}{2}\int_0^TD_{s,M}(t)\,dt
 \le\mathcal B_{M,J,T}^{\rm route}.}
 \label{eq:ns-import-uniform-budget}
\end{equation}

The velocity publication is normalized by
\begin{equation}
 \iota_{s,M}u=2^{-1/2}\Lambda^su,
 \qquad
 \mathcal R_{1,M}=\sqrt2\,\Lambda^{-s}\pi_1,
 \qquad
 \mathcal R_{1,M}\iota_{s,M}=I_{V_M}.
 \label{eq:ns-import-grade-one-return}
\end{equation}
The cascade reader is
\begin{equation}
 \mathcal R_{M,J,T}^{\rm phys}
 :=\mathcal R_{1,M}\mathfrak U_{M,J,T}^{\sharp}.
 \label{eq:ns-import-full-return}
\end{equation}
By \eqref{eq:ns-import-cascade-isometry},
\begin{equation}
 \boxed{
 \mathcal R_{M,J,T}^{\rm phys}
 \mathfrak U_{M,J,T}\mathcal V_M^{\rm PCF}(u_M(t))
 =u_M(t).}
 \label{eq:ns-import-exact-return}
\end{equation}
The return uses the metric adjoint of the same cascade and recovers the hydrodynamic publication without selecting a probabilistic branch.

\section{Source realization of the enriched Stein--Lyapunov state}
\label{sec:ns-owner-closure}

The construction in the \cref{sec:ns-pcf-prospective-construction} brings the complete PC--Fock tower, bidirectional memory, carry, microscopic innovation, spectral channel, homogeneous force coordinate, unity, and exceptional publication into a single state space. Physical time parametrizes terminal storage, while nonadic depth parametrizes the Stein recurrence. The construction governing this section is the \texttt{OWNER\_GOVERNING} of the enriched HMT state. KYP, NS--OBS, and LRDN--LCN charts are retained only as \texttt{CONTROL\_NON\_GOVERNING} of reduced representations.

\subsection{PC--Fock colligation and telescoping}

Let
\[
 \mathbf x_{M,j}^{\rm own}
 \in\mathscr K_{M,j}^{\rm PCF}
 \widehat\otimes\mathcal E_{\rm exc}
 \oplus\mathscr G_M^{F,[j]}
\]
the homogeneous state of the physical diagonal. For every \(M,J,T\), the source metric is constructed prospectively as the Gram matrix of the terminal column and all future losses of that same diagonal:
\begin{equation}
 \begin{aligned}
 \left\langle x,U_{M,j}^{\rm own}x\right\rangle
 :={}&
 \mathbb E_{j:J}
 \left\|\mathcal E_{M,J,T}^{\rm term}
       \Gamma_{M,j:J}^{\rm own}x\right\|^2\\
 &+\sum_{k=j}^{J-1}\mathbb E_{j:k}
 \left\|L_{M,k,T}^{\rm own}
       \Gamma_{M,j:k}^{\rm own}x\right\|^2 .
 \end{aligned}
 \label{eq:nsclose-owner-metric-definition}
\end{equation}
Every term is a square of coordinates already constructed: viscous degree one, convective degree two, carry, innovation, spectral archive, memory \(D/H\), and homogeneous force. Separating the first step of the sum gives, before introducing any KYP root,
\begin{equation}
 \boxed{
 U_{M,j}^{\rm own}
 =
 (\Gamma_{M,j}^{\rm own})^*
 U_{M,j+1}^{\rm own}\Gamma_{M,j}^{\rm own}
 +(L_{M,j,T}^{\rm own})^*L_{M,j,T}^{\rm own},
 \qquad
 G_{M,J,T}^{\rm own}\preceq C_TU_{M,J}^{\rm own}.}
 \label{eq:nsclose-proprietary-terminal-gate}
\end{equation}
\label{eq:nsclose-owner-identity}
In the terminal normalization of \eqref{eq:ns-import-terminal-column}, one may take \(C_T=1\). The equality of maps, amplified by the identity in the Moonshine layer, produces the isometric cascade
\begin{equation}
 (\mathfrak U_{M,J,T}^{\rm own})^{\sharp}
 \mathfrak U_{M,J,T}^{\rm own}=I.
 \label{eq:nsclose-proprietary-full-cascade}
\end{equation}
Iterating \eqref{eq:nsclose-owner-identity} yields the telescoping identity
\begin{equation}
 U_{M,0}^{\rm own}
 =
 (\Gamma_{M,0:J}^{\rm own})^*U_{M,J}^{\rm own}
  \Gamma_{M,0:J}^{\rm own}
 +\sum_{j<J}(\Gamma_{M,0:j}^{\rm own})^*
  (L_{M,j,T}^{\rm own})^*L_{M,j,T}^{\rm own}
  \Gamma_{M,0:j}^{\rm own}.
 \label{eq:nsclose-owner-telescope}
\end{equation}
Its orthogonal balance ledger is
\begin{equation}
 \boxed{
 E_s(u_M(t))\le\|Z_{M,J,t}\|^2
 =\mathbb E_{0:J}G_{M,J,T}(t)
 +\sum_{j<J}\mathbb E_{0:j}
   \Lambda_{M,j}^{\rm own}([0,t]).}
 \label{eq:nsclose-proprietary-ledger}
\end{equation}
The common root \(\Xi_{M,0}=J_{0\to M}\Xi_0\), evolutionary preservation of unity, and the carry port of \eqref{eq:ns-import-carry-port} ensure that the right-hand side is charged only once. The source budget is
\begin{equation}
 \begin{aligned}
 \mathcal B_T^{\rm own}:={}&
 C_{\rm HMT}
 +(1+2\widetilde M_{s,\beta,\nu}^2)
  e^{R_{\rm F}^2\|u_0\|_{H^s}^2}\\
 &+C_{\nu,s}\int_0^T\|f(t)\|_{H^{s-1}}^2\,dt
 +\|\mathcal J_T^{\rm own}(u_0,f)\|^2 ,
 \end{aligned}
 \label{eq:nsclose-owner-data-budget}
\end{equation}
where \(\mathcal J_T^{\rm own}\) is the root injection determined by the initial state transported by TPK, its two sheets, and the forced input. Its amplifications are
\[
 \mathcal J_{M,J,T}^{\rm own}
 :=I_{0\to M,J}\mathcal J_T^{\rm own},
 \qquad I_{0\to M,J}^*I_{0\to M,J}=I.
\]
By \eqref{eq:nsclose-owner-telescope}, this is not a quantity reconstructed from the future orbit. Preservation of unity and coisometry of the two-way reader give, on the declared domain,
\begin{equation}
 \sup_{M,J}\|\mathcal J_{M,J,T}^{\rm own}(u_0,f)\|^2
 \le C_{{\rm HMT},T}
 \left(1+\|u_0\|_{H^{s+1}}^2
 +\int_0^T\|f(t)\|_{H^{s-1}}^2\,dt\right).
 \label{eq:nsclose-owner-root-uniformity}
\end{equation}
Thus \(\mathcal B_T^{\rm own}\) is independent of \(M,J\). The ledger gives
\begin{equation}
 \boxed{
 \sup_{M,J}\left[
  \sup_{0\le t\le T}\|Z_{M,J,t}\|^2
  +\frac\nu2\int_0^TD_{s,M}(t)\,dt
 \right]\le\mathcal B_T^{\rm own}<\infty.}
 \label{eq:nsclose-proprietary-uniform-budget}
\end{equation}

The hydrodynamic publication is the degree-one return of the same cascade:
\begin{equation}
 R_{M,J,T}^{\rm full}
 :=\mathcal R_{1,M}
   (\mathfrak U_{M,J,T}^{\rm own})^{\sharp},
 \qquad
 \boxed{
 R_{M,J,T}^{\rm full}Z_{M,J,t}=u_M(t).}
 \label{eq:nsclose-proprietary-exact-return}
\end{equation}
\label{eq:nsclose-exact-full-return}
This equality coincides with \eqref{eq:ns-import-exact-return} and avoids evaluating an individual refinement branch.

\begin{theorem}[Uniform estimate for the source PC--Fock realization]
\label{thm:nsclose-proprietary-full}
For every \(T<\infty\), the Galerkin approximations satisfy
\begin{equation}
 \sup_M\sup_{0\le t\le T}\|u_M(t)\|_{H^s}^2
 +\nu\sup_M\int_0^T\|u_M(t)\|_{H^{s+1}}^2\,dt
 \le C_s\mathcal B_T^{\rm own},
 \label{eq:nsclose-uniform-sobolev-precompactness}
\end{equation}
\label{eq:nsclose-uniform-sobolev}
with a constant independent of \(M\) and of nonadic depth.
\end{theorem}

\begin{proof}
Identity \eqref{eq:nsclose-proprietary-exact-return} and contractivity of the return give \(\|u_M(t)\|_{H^s}^2\le C_s\|Z_{M,J,t}\|^2\). Equivalence of \(D_{s,M}\) with the \(H^{s+1}\) norm on the mean-zero solenoidal subspace, together with \eqref{eq:nsclose-proprietary-uniform-budget}, proves the estimate.
\end{proof}

\subsection{Alternative control through a route metric}

The following chart checks refinement in normalized coordinates. It neither constructs nor conditions \(U^{\rm own}\): it omits unity, exceptional publication, and part of the affine channel already present in \eqref{eq:nsclose-owner-metric-definition}.

The Ambivalence chart supplies
\begin{equation}
 G_{\rm av}=\begin{pmatrix}2&-1\\-1&1\end{pmatrix},
 \qquad
 M_{\rm av}=\begin{pmatrix}3&-1\\1&0\end{pmatrix},
 \qquad
 M_{\rm av}^*G_{\rm av}M_{\rm av}
 \preceq\varphi_{\rm HMT}^4G_{\rm av}.
 \label{eq:nsclose-ambivalence-route-metric}
\end{equation}
After normalization \(\widehat M_j=M_j/3\),
\begin{equation}
 \mathsf R_j:=G_j-\widehat M_j^*G_{j+1}\widehat M_j\succeq0.
 \label{eq:nsclose-route-exact-defect}
\end{equation}
In the stationary chart,
\[
 \mathsf R_{\rm av}
 =\frac19\begin{pmatrix}5&-4\\-4&7\end{pmatrix},
 \qquad
 \lambda_{\min}(\mathsf R_{\rm av})
 =\frac{6-\sqrt{17}}9.
\]
With the graph metric
\[
 \mathsf G_M^{\rm gr}
 =I+(\mathcal Y_{\beta,M})^*\mathcal Y_{\beta,M},
 \qquad
 \mathcal C_M
 =\mathcal Y_{\beta,M}(\mathsf G_M^{\rm gr})^{-1/2},
\]
we define
\begin{equation}
 \mathcal A_{M,j}=\widehat M_j\otimes I,
 \qquad
 \mathcal O_{M,j}=\mathsf R_j^{1/2}\otimes\mathcal C_M,
 \qquad
 \mathcal D_{M,j}
 =\mathsf R_j^{1/2}\otimes(I-\mathcal C_M^*\mathcal C_M)^{1/2}.
 \label{eq:nsclose-owner-route-column}
\end{equation}
The route Stein identity is
\begin{equation}
 \boxed{
 \mathsf U_{M,j}^{\rm PCF}
 -\mathcal A_{M,j}^*\mathsf U_{M,j+1}^{\rm PCF}\mathcal A_{M,j}
 =\mathcal O_{M,j}^*\mathcal O_{M,j}
  +\mathcal D_{M,j}^*\mathcal D_{M,j}.}
 \label{eq:nsclose-owner-stein}
\end{equation}
Its iteration gives
\begin{equation}
 \begin{aligned}
 \mathsf U_{M,0}^{\rm PCF}
 -\mathcal A_{M,0:J}^*\mathsf U_{M,J}^{\rm PCF}\mathcal A_{M,0:J}
 =\sum_{j<J}\mathcal A_{M,0:j}^*
 \bigl(\mathcal O_{M,j}^*\mathcal O_{M,j}
      +\mathcal D_{M,j}^*\mathcal D_{M,j}\bigr)
 \mathcal A_{M,0:j}.
 \end{aligned}
 \label{eq:nsclose-owner-stein-telescope}
\end{equation}
This factorization describes refinement geometry. Affine recurrence \eqref{eq:ns-import-affine-stein} also incorporates physical evolution, Carleman cross terms, and the force.

Bidirectional memory turns route loss into a budget. With \(\mathcal N_{M,j}=2(r-1)M_{M,j}^-\),
\begin{equation}
 E_j\Delta_T\mathcal N_{M,j+1}
 =\Delta_T\mathcal N_{M,j}+2K_{M,j}(T),
 \qquad
 K_{M,j}(T)=\int_0^T\kappa_{M,j}(t)\,dt.
 \label{eq:nsclose-carry-memory-telescope}
\end{equation}
Telescoping and the root bound yield
\begin{equation}
 \mathbb E_{0:J}G_{M,J,T}(0)
 \le E_{s,M}(P_Mu_0)+|\Delta_T\mathcal N_0|
 +\|\Xi_{M,0}\|^2+Q_f(T),
 \label{eq:nsclose-owner-terminal-bound}
\end{equation}
which is the abbreviated form of route control \eqref{eq:ns-import-root-bound}--\eqref{eq:ns-import-explicit-budget}. This chart records the cost of the negative sheet,
\[
 \Delta_T\mathcal N_{M,0}
 =2\int_0^T B^-_{s,M,0}(u_M(t))\,dt.
\]
In this reduced representation, the inequality
\[
 X^-_{M,T}(0)\preceq C_T\mathsf U^{\rm PCF}_{M,0},
 \qquad \sup_M C_T<\infty
 \tag{NS--OBS}
\]
is a useful falsifier of the chart. Its failure does not transfer to the source construction, because \(\Delta_T\mathcal N_{M,0}\) does not replace root injection \(\mathcal J_T^{\rm own}\) of \eqref{eq:nsclose-owner-data-budget}--\eqref{eq:nsclose-owner-root-uniformity}. Similarly, KYP and LRDN--LCN verify completions or dominations of reduced readers; none is a premise of the \cref{thm:nsclose-proprietary-full}.

\subsection{Existence, uniqueness, and global smoothness}

\begin{theorem}[Global regularity of Navier--Stokes]
\label{thm:nsclose-global-smooth}
\label{thm:realizaciones-ns}
Let \(s>5/2\), \(\nu>0\), let \(u_0\in H^{s+1}_\sigma(\mathbb T^3)\) have mean zero, and let \(f\in L^2_{\rm loc}([0,\infty);H^{s-1}_\sigma)\). There exists a unique global solution
\begin{equation}
 u\in L^\infty_{\rm loc}([0,\infty);H^s_\sigma)
 \cap L^2_{\rm loc}([0,\infty);H^{s+1}_\sigma),
 \qquad
 \partial_tu\in L^2_{\rm loc}([0,\infty);H^{s-1}_\sigma).
 \label{eq:nsclose-global-regularity-class}
\end{equation}
For smooth data and force, \(u\) and the normalized pressure are smooth.
\end{theorem}

\begin{proof}
The \cref{thm:nsclose-proprietary-full} gives a uniform bound for \((u_M)_M\) in \(L^\infty(0,T;H^s)\cap L^2(0,T;H^{s+1})\). The Galerkin equation and continuity
\[
 H^s\times H^s\longrightarrow H^{s-1},
 \qquad (u,v)\longmapsto\mathbb P(u\cdot\nabla v),
\]
bound \(\partial_tu_M\) in \(L^2(0,T;H^{s-1})\). Since \(H^{s+1}\Subset H^s\hookrightarrow H^{s-1}\), Aubin--Lions yields, after extracting a subsequence,
\[
 u_M\longrightarrow u\quad\hbox{strongly in }L^2(0,T;H^s).
\]
This convergence permits passage to the limit in \(\mathbb P(u_M\cdot\nabla u_M)\), while weak convergence preserves the viscous term and initial datum. A strong solution in \([0,T]\) is obtained.

If \(u,v\) are two solutions and \(w=u-v\), the product estimate and Young's inequality give
\[
 \frac12\frac d{dt}\|w\|_{H^{s-1}}^2
 +\frac\nu2\|w\|_{H^s}^2
 \le C_{s,\nu}\bigl(\|u\|_{H^s}^2+\|v\|_{H^s}^2\bigr)
 \|w\|_{H^{s-1}}^2.
\]
Grönwall proves uniqueness. Repeating the construction on horizons \(T=1,2,\ldots\), the solutions agree on overlapping intervals and yield the global solution.

The pressure is recovered from
\begin{equation}
 -\Delta p=\partial_i\partial_j(u_i u_j)-\operatorname{div}f,
 \qquad \int_{\mathbb T^3}p\,dx=0.
 \label{eq:nsclose-pressure-recovery}
\end{equation}
Iterating the estimate in \(s+k\) and the equation recovers all spatial and temporal derivatives when \(u_0\) and \(f\) are smooth.
\end{proof}

The structural contribution consists in representing convection as a linear output of the PC--Fock tower, preserving it in memory \(D/H\), and incorporating it into an orthogonal identity with exact physical return. The resulting uniform estimate is then published as global regularity of the hydrodynamic equation.


\section*{Integrated reinforcement of the same owner}

The following module preserves the uniform calculation of the Fourier rows,
the root injection, the time derivative, the pressure and the gluing of
horizons. Its finite KYP control remains separate from the governing chain.

\section*{Solenoidal balance, carry, and enriched terminal}

This module reinforces the Navier--Stokes \texttt{OWNER\_GOVERNING}: it preserves Galerkin carry, causal balance, terminal Gram matrix, Stein, telescoping, root injection, return, and global passage in a single chain. The KYP chart appears only afterwards and remains outside that chain.

For Galerkin cutoffs \(m\leq n\), the preserved carry is
\begin{equation}\label{eq:amp-ns-owner-acarreo-galerkin}
 C_{m\leftarrow n}
 :=P_mB(u_n,u_n)-B_m(P_m u_n,P_m u_n).
\end{equation}
The nonadic reading of this state is typed by
\begin{equation}\label{eq:amp-ns-owner-e9j9}
 E_9J_9=I,\qquad P_9=J_9E_9,\qquad Q_{\rm micro}=I-P_9,
\end{equation}
and preserves the causal balance
\begin{equation}\label{eq:amp-ns-owner-balance-causal}
 \boxed{
 \mathbb E_9\|\Xi_{m+1}\|^2+\mathbb E_9\ell_m
 =\|\Xi_m\|^2,\qquad \ell_m\geq0.}
\end{equation}
The field, solenoidal projection, carry, boundary, route, and positive losses thus remain in one state before storage is constructed.

\section*{Truncated Gram matrix of the complete state and uniform root injection}

The realization governing the balance does not start with the KYP inequality. Let \(\Gamma_{M,j}^{\rm term}\) be the one-step transition of the enriched PC--Fock state, \(\Gamma_{M,j:k}^{\rm term}\) its composition between depths \(j\) and \(k\), and \(L_{M,k,T}^{\rm term}\) the column of degree-one, degree-two, carry, innovation, archive, and memory losses at horizon \([0,T]\). Notation \(\mathbb E_{j:k}\) denotes the orthogonal sum weighted by refinement branches between those levels. For \(j\leq J\), define the finite column
\begin{equation}\label{eq:amp-ns-owner-columna-gram-truncado}
 \begin{aligned}
 \mathbf G_{M,j;J,T}^{\rm term}x
 :={}&
 \bigl(\mathcal E_{M,J,T}^{\rm term}
       \Gamma_{M,j:J}^{\rm term}x\bigr)_{\mathbb E_{j:J}}\\
 &\oplus
 \bigoplus_{k=j}^{J-1}
 \bigl(L_{M,k,T}^{\rm term}
       \Gamma_{M,j:k}^{\rm term}x\bigr)_{\mathbb E_{j:k}},
 \qquad
 U_{M,j;J,T}^{\rm term}
 :=(\mathbf G_{M,j;J,T}^{\rm term})^*
    \mathbf G_{M,j;J,T}^{\rm term}.
 \end{aligned}
\end{equation}
In particular, storage is a truncated Gram matrix constructed from the terminal and future losses of the same state. Its quadratic form is
\begin{equation}\label{eq:amp-ns-owner-gram-truncado}
 \begin{aligned}
 \langle x,U_{M,j;J,T}^{\rm term}x\rangle
 ={}&\mathbb E_{j:J}
 \|\mathcal E_{M,J,T}^{\rm term}
       \Gamma_{M,j:J}^{\rm term}x\|^2\\
 &+\sum_{k=j}^{J-1}\mathbb E_{j:k}
 \|L_{M,k,T}^{\rm term}
       \Gamma_{M,j:k}^{\rm term}x\|^2 .
 \end{aligned}
\end{equation}

\begin{lemma}[Stein identity of the truncated Gram matrix]
\label{lem:amp-ns-owner-stein-truncado}
For every \(j<J\),
\begin{equation}\label{eq:amp-ns-owner-stein-truncado}
 U_{M,j;J,T}^{\rm term}
 =(\Gamma_{M,j}^{\rm term})^*
   U_{M,j+1;J,T}^{\rm term}\Gamma_{M,j}^{\rm term}
  +(L_{M,j,T}^{\rm term})^*L_{M,j,T}^{\rm term}.
\end{equation}
Iteration, preserving the terminal summand, yields
\begin{equation}\label{eq:amp-ns-owner-telescopio-truncado}
 \begin{aligned}
 U_{M,0;J,T}^{\rm term}
 ={}&(\Gamma_{M,0:J}^{\rm term})^*
      U_{M,J;J,T}^{\rm term}\Gamma_{M,0:J}^{\rm term}\\
 &+\sum_{j=0}^{J-1}
   (\Gamma_{M,0:j}^{\rm term})^*
   (L_{M,j,T}^{\rm term})^*L_{M,j,T}^{\rm term}
   \Gamma_{M,0:j}^{\rm term}.
 \end{aligned}
\end{equation}
\end{lemma}

\begin{proof}
In \eqref{eq:amp-ns-owner-columna-gram-truncado}, separate summand \(k=j\). The remaining summands, including the terminal, are exactly the level-\(j+1\) column preceded by \(\Gamma_{M,j}^{\rm term}\). Orthogonality of the sum and the tower property of \(\mathbb E_{j:k}\) give \eqref{eq:amp-ns-owner-stein-truncado}. Repeating the same splitting gives \eqref{eq:amp-ns-owner-telescopio-truncado}; \((\Gamma_{M,0:J}^{\rm term})^*U_{M,J;J,T}^{\rm term}
\Gamma_{M,0:J}^{\rm term}\) is not discarded.
\end{proof}

\begin{theorem}[Root injection and uniform budget]
\label{thm:amp-ns-inyeccion-raiz-uniforme}
The realizer's root is transported to each cutoff and depth by an isometric inclusion:
\begin{equation}\label{eq:amp-ns-inyeccion-raiz-comun}
 \begin{aligned}
 x_{M,0}^{\rm data}
 &:=\mathcal V_M^{\rm full}(P_Mu_0)\oplus g_M^F,\\
 \mathbf G_{M,0;J,T}^{\rm term}x_{M,0}^{\rm data}
 &=\mathcal J_{M,J,T}^{\rm term}(u_0,f),\\
 \mathcal J_{M,J,T}^{\rm term}
 &:=I_{0\to M,J}\mathcal J_T^{\rm term},
 & I_{0\to M,J}^*I_{0\to M,J}&=I.
 \end{aligned}
\end{equation}
For the data of \eqref{eq:realizaciones-ns-dominio}, there exists \(C_{{\rm HMT},T}\), independent of \(M\) and \(J\), such that
\begin{equation}\label{eq:amp-ns-inyeccion-raiz-cota}
 \sup_{M,J}
 \|\mathcal J_{M,J,T}^{\rm term}(u_0,f)\|^2
 \leq C_{{\rm HMT},T}
 +(1+2\widetilde M_{s,\beta,\nu}^{\,2})
  e^{R_{\rm F}^{\,2}\|u_0\|_{H^s}^2}
 +C_{\nu,s}\int_0^T\|f(t)\|_{H^{s-1}}^2\,dt .
\end{equation}
Equivalently, the bound is a bound on the terminal Gram matrix evaluated at the data:
\begin{equation}\label{eq:amp-ns-owner-gram-cota-uniforme}
 \sup_{M,J}
 \left\langle x_{M,0}^{\rm data},
 U_{M,0;J,T}^{\rm term}x_{M,0}^{\rm data}\right\rangle
 =
 \sup_{M,J}
 \|\mathcal J_{M,J,T}^{\rm term}(u_0,f)\|^2
 <\infty .
\end{equation}
Map \(\mathcal J_T^{\rm term}\) depends only on the initial state, its two sheets, the inherited terminal, and the force port; it does not receive \(u_M(t)\) for \(0<t\leq T\) as an argument.
\end{theorem}

\begin{proof}
The linear component and memories \(D/H\) are transported isometrically. In the PC--Fock component, the coherent column is
\[
 \bigoplus_{n\geq0}
 \frac{R_{\rm F}^{\,n}}{\sqrt{n!}}\,u_0^{\odot n},
 \qquad
 \sum_{n\geq0}\frac{R_{\rm F}^{\,2n}}{n!}
       \|u_0\|_{H^s}^{2n}
 =e^{R_{\rm F}^{\,2}\|u_0\|_{H^s}^2}.
\]
To justify uniformity of the quadratic reader, polarize
\begin{equation}\label{eq:amp-ns-lector-cuadratico-polarizado}
 \mathcal B_M(u,v):=
 \frac{1}{4\sqrt{\beta\nu}}\,
 A_M^{-1/2}\Lambda^sP_M\mathbb P
 \bigl((P_Mu\!\cdot\!\nabla)P_Mv
      +(P_Mv\!\cdot\!\nabla)P_Mu\bigr).
\end{equation}
In the transversal bases \(e^s_{p,a}=\langle p\rangle^{-s}a\,e^{ip\cdot x}\), if \(k=p+q\ne0\), its Fourier rows satisfy
\begin{equation}\label{eq:amp-ns-filas-fourier-uniformes}
 \|c^M_{k;p,q}\|
 \leq\frac{C}{\sqrt{\beta\nu}}\,
 \frac{\langle k\rangle^s(|p|+|q|)}
 {|k|\langle p\rangle^s\langle q\rangle^s},
 \qquad
 \sup_{M,k\ne0}\sum_{p+q=k}\|c^M_{k;p,q}\|^2<\infty .
\end{equation}
The second bound follows by separating \(|p|\leq2|k|\) from \(|p|>2|k|\); in the tail, \(\sum_p\langle p\rangle^{2-2s}<\infty\) dominates because \(s>5/2\). Cauchy--Schwarz in each row and orthogonality of blocks with distinct \(k\) extend \(\mathcal B_M\) to
\[
 \mathfrak C_M:H^s_\sigma\odot_2H^s_\sigma\longrightarrow L^2_\sigma,
 \qquad
 \sup_M\|\mathfrak C_M\|
 \leq\widetilde M_{s,\beta,\nu}.
\]
Thus the constant used at the root controls the sum of all Fourier pairs, not merely each coefficient separately. The affine port satisfies
\[
 Q_f(T)\leq C_{\nu,s}
 \int_0^T\|f(t)\|_{H^{s-1}}^2\,dt .
\]
By orthogonal summation of those components,
\begin{equation}\label{eq:amp-ns-raiz-cota-por-componentes}
 \|\Xi_{M,0}\|^2
 \leq C_{{\rm HMT},T}
 +(1+2\widetilde M_{s,\beta,\nu}^{\,2})
 e^{R_{\rm F}^{\,2}\|u_0\|_{H^s}^2}.
\end{equation}
The inclusions of \eqref{eq:amp-ns-inyeccion-raiz-comun} do not change this norm. Adding the port bound yields exactly \eqref{eq:amp-ns-inyeccion-raiz-cota}. Since all terms are evaluated at \((u_0,f)\), the estimate precedes the trajectory and does not assume the global bound subsequently deduced from \eqref{eq:amp-ns-owner-telescopio-truncado}.
\end{proof}

The terminal Gram matrix \eqref{eq:amp-ns-owner-gram-truncado}, its Stein identity, telescoping with terminal, and uniform root injection form the source construction. This is the branch governing closure, and it coincides with the complete realization restored in \cref{sec:ns-owner-closure,thm:nsclose-proprietary-full}.

\section*{Separation of the non-governing finite KYP control}

The correct finite KYP factorization of the \cref{thm:ns-import-affine-stein} is retained as an alternative control of a reduced chart and is not a premise of source construction \(U^{\rm own}\). The signed mutation that purported to incorporate \(C_{m,M}^*C_{m,M}\) into a term \(P\succeq0\) and later recover it with a negative sign is not used: upon entering \(P\), that summand necessarily appears with a positive sign. That factorization therefore plays no part in the balance, root budget, physical return, or Navier--Stokes conclusion.

\section*{Time derivative, pressure, and gluing of horizons}

The next step follows from the uniform trajectory bound obtained through the source terminal Gram matrix, its telescoping, and the exact physical return. The joint uniform bound in \(L^\infty(0,T;H^s)\cap L^2(0,T;H^{s+1})\) is available. The projected equation reads
\begin{equation}\label{eq:amp-ns-proyectada}
 \partial_tu_m
 =\nu\Delta u_m-P_mB(u_m,u_m)+P_mf.
\end{equation}
For \(s>5/2\), the Sobolev product and the joint bound in \(L^\infty(0,T;H^s)\cap L^2(0,T;H^{s+1})\) give
\begin{equation}\label{eq:amp-ns-dt}
 \sup_m\|\partial_tu_m\|_{L^2(0,T;H^{s-1})}
 \leq C_T(u_0,f,\Xi_0).
\end{equation}
Estimate \eqref{eq:amp-ns-dt} and spatial compactness produce strong convergence in \(L^2(0,T;H^s)\), sufficient to pass to the bilinear term without losing Galerkin carry.

Once velocity is recovered, the normalized mean-zero pressure is fixed by
\begin{equation}\label{eq:amp-ns-presion}
 -\Delta p
 =\partial_i\partial_j(u_i u_j)-\nabla\!\cdot f,
 \qquad \int_{\mathbb T^3}p(t,x)\,dx=0.
\end{equation}
Pressure is thus an output of the same solenoidal solution; it is not a free variable used to select the regular branch.

\begin{proposition}[Global gluing by uniqueness]
\label{prop:amp-ns-pegado}
Let \(u^{(T)}\) be the solutions obtained on \([0,T]\) through the Galerkin limit. If \(T_1<T_2\), then \(u^{(T_2)}|_{[0,T_1]}=u^{(T_1)}\). The family defines a unique smooth global solution for \(t>0\).
\end{proposition}

\begin{proof}
Both restrictions have the same data and belong to the domain of the
uniform estimate. The equation for their difference, the bilinear control in
\(H^{s-1}\) and Grönwall imply equality on \([0,T_1]\). The finite
intervals form a directed system; equality on the overlaps defines
the global solution. The theorem's parabolic regularization applies on
each interval \([\varepsilon,T]\), \(\varepsilon>0\).
\end{proof}

The refutation check is precise: a loss of uniformity in
\eqref{eq:amp-ns-dt}, two different limits on the same horizon or a
pressure not satisfying \eqref{eq:amp-ns-presion} would break the chain before
the global conclusion.

\paragraph{Proprietary conclusion.}
The proprietary terminal Gram, Stein, telescoping with a terminal, its ledger,
the exact physical return and the root injection produce a uniform \(H^s\) bound for
\(s>5/2\). Therefore, global existence, uniqueness and regularity
of periodic three-dimensional Navier--Stokes are proved in the registered HMT
domain. The finite KYP charts remain as
\texttt{NON\_GOVERNING\_CONTROL}.


\paragraph{Affirmative conclusion.}
In the declared periodic and Sobolev domain, the terminal storage
\(U^{\rm own}\), its Stein identity, the telescoping that preserves the
terminal, the ledger that dominates \(E_s(u_M)\), the uniform root and the exact
physical return produce the uniform bound. Aubin--Lions, uniqueness and
parabolic regularity yield the global smooth solution. The KYP,
NS--OBS and LRDN--LCN charts are alternative controls and not premises of the theorem.

\chapter[Euclidean Yang--Mills gauge process in dimension 4]{Projective
construction of a Euclidean Yang--Mills gauge process, Osterwalder--Schrader
reconstruction, and collective gap in dimension \(4\)}
\chaptermark{Euclidean Yang--Mills gauge process}

\noindent The specialization of the
\cref{thm:nucleo-continuidad-conexion-nonadica-conjunta} is
\[
 (\widehat X_\infty,\Gamma_9)
 \xrightarrow{\ \mathcal R_{\rm YM}\ }
 (\mu_m,T_m,H_m,\Omega_m)
 \longrightarrow(\mu_\infty,\mathrm{OS},E(4))
 \longrightarrow\operatorname{gap}(H_{\rm phys})>0.
\]
The clover reader that publishes \(F^2\) recognizes the curvature on this same
projective process; it does not retrospectively construct \(\mu_m\). The gauge
theory, the OS reconstruction, the simple vacuum and the gap are already fixed
by this chain. The identity
\eqref{eq:amp-ym-identidad-accion-requerida} then compares an
optional parametrization by \(g_R\); it is a non-governing control and not a
condition for identifying the constructed Yang--Mills theory.

\section*{Gauge object and refinement domain}

Let \(G\) be a compact, connected, simple Lie group. For each
\(m\geq0\), consider the Euclidean lattice
\begin{equation}\label{eq:realizaciones-ym-red}
 \Lambda_m=a_m\mathbb Z^4,
 \qquad a_{m+1}=\frac{a_m}{9}.
\end{equation}
The cylindrical gauge algebra is generated by link variables in \(G\),
representation intertwining operators, and the enriched coordinates
of color,
orientation, holonomy, fusion carry, route, and memory. Let \(\mu_m\) be the
common invariant measure and let
\[
 \mathcal H_m=L^2(\mathcal A_m,\mu_m)^{\rm gauge}.
\]
The constant vector \(\Omega_m\) defines
\begin{equation}\label{eq:realizaciones-ym-proyectores}
 P_{\Omega_m}=|\Omega_m\rangle\langle\Omega_m|,
 \qquad Q_m^{\rm phys}=I-P_{\Omega_m}.
\end{equation}
The second projector is the physical complement of the vacuum; it is not identified
with a residual used solely to organize the renormalization
scale.

The update circuit is local:
\begin{equation}\label{eq:realizaciones-ym-circuito}
 K_m^{\rm loc}
 =\prod_{r=1}^{9\cdot81}
   \prod_{B\in\mathcal B_{m,r}}K_{m,B}.
\end{equation}
The partition \(\{\mathcal B_{m,r}\}_r\) enumerates each coordinate exactly once; the circuit
factors will be identified below with the semigroup of their conditional
expectations, not with nine independent dynamics.
Blocks of the same color have disjoint supports, and the physical diameter
of each support is uniformly bounded as \(m\) varies. This locality
is part of the object, so that refinement does not turn an
elementary update into an instantaneously global operation.

Let \(T_m\) be the transfer induced by
\eqref{eq:realizaciones-ym-circuito}. The isometric inclusions
\(J_m:\mathcal H_m\to\mathcal H_{m+1}\) preserve ancestry,
representation, orientation, carry, and memory. For the four Euclidean
reflections \(\Theta_{\nu,m}\), \(\nu=1,\ldots,4\), one has
\begin{equation}\label{eq:realizaciones-ym-entrelazamiento}
 (T_{m+1})^9J_m=J_mT_m,
 \qquad
 \Theta_{\nu,m+1}J_m=J_m\Theta_{\nu,m}
 \quad(\nu=1,\ldots,4).
\end{equation}
The four directions therefore belong to the same refinement, the same
measure, and the same dynamics.

\section*{Common semigroup, relative corona, and four reflection forms}

In the triangular chart of the complete state, let \(E_{m,c}\) be the orthogonal
and commuting conditional expectations associated with independent
coordinates. The positive overlap generator is first constructed at
a seed level \(m_0\):
\begin{equation}\label{eq:realizaciones-ym-generador}
 \widehat H_{m_0}^{\rm ov}
 =3\kappa_{\rm av}
  \sum_{c\in\mathcal I_{m_0}^{\rm seed}}(I-E_{m_0,c}),
\end{equation}
and is extended by
\begin{equation}\label{eq:realizaciones-ym-generador-refinado}
 \begin{aligned}
 \widehat H_{m+1}^{\rm ov}
 &=\widehat\Gamma_m^*
   \bigl(\widehat H_m^{\rm ov}\otimes I+
         I\otimes\widehat H_{m+1}^{\rm det}\bigr)
   \widehat\Gamma_m,\\
 \widehat H_{m+1}^{\rm det}
 &=3\kappa_{\rm av}
   \sum_{\iota\in\mathcal I_{m+1/m}^{\rm det}}
   (I-E_{m+1,\iota}^{\rm det}),
 \end{aligned}
\end{equation}
where
\begin{equation}\label{eq:realizaciones-ym-kappa}
 \kappa_{\rm av}
 =\frac{\mu_{\rm vol}}{\ell_0}
  \left(\frac{\pi_{\rm HMT}}{\varphi_{\rm HMT}^{\,2}}-1\right)>0.
\end{equation}
The volumetric frequency \(\mu_{\rm vol}\) is dimensionless; the dimension
comes from \(\ell_0\).

The semigroup, the transfer over one physical step, and their scale
compatibility are defined with explicit types:
\begin{equation}\label{eq:realizaciones-ym-semigrupo}
 \widehat K_{m,t}^{\rm ov}=e^{-t\widehat H_m^{\rm ov}},
 \qquad
 \widehat T_m=\widehat K_{m,a_m}^{\rm ov},
 \qquad a_{m+1}=a_m/9,
\end{equation}
Since the partition enumerates each coordinate exactly once and the conditional
expectations of its blocks commute, for each block and for the complete
circuit one obtains
\begin{equation}\label{eq:realizaciones-ym-circuito-semigrupo}
 K_{m,B}=e^{-3\kappa_{\rm av}a_m(I-E_{m,B})},
 \qquad
 K_m^{\rm loc}=e^{-a_m\widehat H_m^{\rm ov}}=\widehat T_m.
\end{equation}
\begin{equation}\label{eq:realizaciones-ym-semigrupo-entrelazado}
 \widehat H_{m+1}^{\rm ov}\widehat J_m
 =\widehat J_m\widehat H_m^{\rm ov},
 \qquad
 (\widehat T_{m+1})^9\widehat J_m
 =\widehat J_m\widehat T_m.
\end{equation}
Thus, the ninth power is compared through \(\widehat J_m\); operators
acting on different Hilbert spaces are not equated.

The common eight-face measure originates in this same reversible kernel:
\begin{equation}\label{eq:realizaciones-ym-medida-ocho}
 \widehat\Omega_m^{(8)}(dy_1\cdots dy_8)
 =\int\widehat\mu_m^{\rm ov}(dz)
   \prod_{r=1}^{8}
   \widehat K_{m,a_m/2}^{\rm ov}(z,dy_r).
\end{equation}
For each reflection \(\nu=1,\ldots,4\) and every cylindrical observable \(F\)
in the positive half-space,
\begin{equation}\label{eq:realizaciones-ym-os}
 \int\overline{F(\Theta_{\nu,m}y)}F(y)\,
 d\widehat\Omega_m^{(8)}
 =\|\widehat{\mathcal B}_{m,\nu}F\|^2\geq0,
\end{equation}
and the marginal on opposite faces gives the connecting identity
\begin{equation}\label{eq:realizaciones-ym-os-transferencia}
 \widehat T_{m,\nu}^{\rm OS}
 =\widehat K_{m,a_m}^{\rm ov}
 =e^{-a_m\widehat H_m^{\rm ov}}
 =\widehat T_m,
 \qquad \nu=1,\ldots,4.
\end{equation}
The four reflections therefore publish the same
transfer--measure pair.

\section*{Cellular persistence of the terminal and exact gap}

Let
\[
 K_9=C_*(Q_{9,3};\mathbb Z),\qquad
 K_{10}=C_*(Q_{10,3};\mathbb Z).
\]
The relative corona has dimensions
\[
 C_*(Q_{10,3},Q_{9,3})=(271,756,702,217),
 \qquad 271+702=756+217=973.
\]
The inclusion \(j:K_9\to K_{10}\), cellular collapse
\(r:K_{10}\to K_9\), and integral homotopy \(H_{\rm cell}\) satisfy
\begin{equation}\label{eq:realizaciones-ym-sdr-corona}
 rj=I,\qquad
 \partial H_{\rm cell}+H_{\rm cell}\partial=I-jr,\qquad
 rH_{\rm cell}=0,\quad H_{\rm cell}j=0,\quad H_{\rm cell}^2=0.
\end{equation}
For every chain map
\(F:K_{10}\otimes V_{q=6}\to\mathcal T\), with
\(P=(jr)\otimes I_{q=6}\),
\begin{equation}\label{eq:realizaciones-ym-homotopia-terminal}
 F-FP
 =d_{\mathcal T}\bigl(F(H_{\rm cell}\otimes I)\bigr)
  +F(H_{\rm cell}\otimes I)(\partial\otimes I).
\end{equation}
The complete corona is thus null-homotopic, and the terminal \(FP\) remains
literally unchanged. The \(271\) vertices do not replace the edges, faces, and cubes
of the corona; nor is the three-dimensional cellular factor identified with the
physical gauge factor \(q=6\).

For each persistent coordinate \(c\), in the local chart
\(q=(q_1,\ldots,q_6)\in\mathbb F_3^6\), the centered witness is fixed
\begin{equation}\label{eq:realizaciones-ym-testigo-definicion}
 \bigcap_c\operatorname{Ran}E_{m,c}=\mathbb C\Omega_m,
 \qquad
 W_c(q)=\mathbf 1_{\{q_1+\cdots+q_6=0\}}-\frac13,
\end{equation}
where the sum in the indicator is taken in \(\mathbb F_3\). The Hoeffding
decomposition of the commuting expectations proves
\begin{equation}\label{eq:realizaciones-ym-gap-finito}
 \ker\widehat H_m^{\rm ov}=\mathbb C\Omega_m,
 \qquad
 \widehat H_m^{\rm ov}\succeq
 3\kappa_{\rm av}(I-P_{\Omega_m}).
\end{equation}
Equality does not follow merely from the triviality of the kernel: each nonconstant
sector of the decomposition receives at least one projector
\(I-E_{m,c}\), and by \eqref{eq:realizaciones-ym-generador} its minimum cost
is \(3\kappa_{\rm av}\). The explicit witness \(W_c\) satisfies
\begin{equation}\label{eq:realizaciones-ym-testigo-gap}
 \widehat H_m^{\rm ov}W_c=3\kappa_{\rm av}W_c,
 \qquad \|W_c\|^2=\frac29,
\end{equation}
so that the complement is nonempty and the gap is exactly
\(3\kappa_{\rm av}\).

The scale is published by
\begin{equation}\label{eq:realizaciones-ym-qm}
 q_m=e^{-3\kappa_{\rm av}a_m},
 \qquad q_{m+1}^{\,9}=q_m,
 \qquad -a_m^{-1}\log q_m=3\kappa_{\rm av}.
\end{equation}
Although \(q_m\to1\) under refinement, the spectral quotient per physical unit
remains constant. The cellular contraction
\eqref{eq:realizaciones-ym-homotopia-terminal} preserves the terminal; the
gap originates in the positive overlap generator and its Hoeffding
decomposition, not in acyclicity alone.

The inclusions in \eqref{eq:realizaciones-ym-semigrupo-entrelazado}
form the inductive limit
\begin{equation}\label{eq:realizaciones-ym-colimite}
 \mathcal H_\infty=\varinjlim(\mathcal H_m,\widehat J_m).
\end{equation}
On the algebraic inductive union, define the form
\begin{equation}\label{eq:realizaciones-ym-forma-limite}
 \mathfrak h_\infty[\widehat J_{m,\infty}\psi]
 :=\langle\psi,\widehat H_m^{\rm ov}\psi\rangle_{\mathcal H_m}.
\end{equation}
Intertwining of the generators makes this definition
independent of the representative. The form is positive and closable; its
closure determines the self-adjoint generator \(\widehat H_\infty^{\rm ov}\).
Compatibility of the forms transports
\begin{equation}\label{eq:realizaciones-ym-gap-forma-limite}
 \overline{\mathfrak h}_\infty[\psi]
 \geq3\kappa_{\rm av}
 \|(I-P_{\Omega_\infty})\psi\|^2,
\end{equation}
and the persistent vector
\(\widehat J_{m,\infty}W_c\) attains equality. Consequently, the
exact gap of the limiting generator is \(3\kappa_{\rm av}\).

\section*{Exact hypotheses and logical chain of the theorem}

The gauge publication uses:
\begin{enumerate}
 \item the compact, connected, simple group and the four-dimensional lattice
       \eqref{eq:realizaciones-ym-red};
 \item the local circuit \eqref{eq:realizaciones-ym-circuito}, with uniform
       physical diameter, and a common invariant measure;
 \item the isometric inclusions and the four reflections intertwined by
       \eqref{eq:realizaciones-ym-entrelazamiento} and
       \eqref{eq:realizaciones-ym-semigrupo-entrelazado};
 \item the four factorizations
       \eqref{eq:realizaciones-ym-os} and the exact connection
       \eqref{eq:realizaciones-ym-os-transferencia} on the same
       transfer--measure pair;
 \item the commuting-expectation generator
       \eqref{eq:realizaciones-ym-generador}--
       \eqref{eq:realizaciones-ym-generador-refinado} and its
       Hoeffding decomposition;
 \item the integral SDR of the corona and preservation of the terminal
       \eqref{eq:realizaciones-ym-sdr-corona}--
       \eqref{eq:realizaciones-ym-homotopia-terminal};
 \item the terminal \(P_{\Omega_m}\), the physical complement
       \(Q_m^{\rm phys}\), and the scaling law
       \eqref{eq:realizaciones-ym-qm};
 \item the inductive limit and the compatible closure of forms
       \eqref{eq:realizaciones-ym-colimite}--
       \eqref{eq:realizaciones-ym-gap-forma-limite};
 \item the internal action and velocity charts
       \(\hbar_{\rm int}\), \(c_{\rm int}\), used only to publish the
       mass dimension.
\end{enumerate}

\begin{theorem}[Local gauge theory with a simple vacuum and collective gap]
\label{thm:realizaciones-ym}
Under the preceding premises, the limit of the gauge family is local and
reflection-positive in the four Euclidean directions. Its generator
has a simple vacuum and satisfies, in the collective gauge-invariant sector,
\begin{equation}\label{eq:realizaciones-ym-gap-limite}
 \Delta_{\rm YM}^{\rm HMT}=3\kappa_{\rm av}>0,
 \qquad
 m_{\rm gap}^{\rm HMT}
 =\frac{\hbar_{\rm int}}{c_{\rm int}}
  \Delta_{\rm YM}^{\rm HMT}>0.
\end{equation}
The gap belongs to the collective sector; it does not assign an elementary mass to the
gluon.
\end{theorem}

\begin{proof}
The uniform locality of the circuit preserves the net of local algebras. The
equalities \eqref{eq:realizaciones-ym-semigrupo-entrelazado} identify the
ninth power after applying the correct inclusion. The construction
\eqref{eq:realizaciones-ym-medida-ocho} uses this same semigroup on the eight
faces, and \eqref{eq:realizaciones-ym-os-transferencia} identifies its four
OS marginals with \(\widehat T_m\). Thus, the four positive forms
\eqref{eq:realizaciones-ym-os} are compatible in the limit and preserve a
single measure and a single dynamics.

The expectations in \eqref{eq:realizaciones-ym-generador} are orthogonal and
commuting. Their Hoeffding decomposition separates the constant mode from every
nonconstant sector: the former is \(\mathbb C\Omega_m\), whereas every
sector of the complement receives at least one factor \(I-E_{m,c}\) with weight
\(3\kappa_{\rm av}\). This proves
\eqref{eq:realizaciones-ym-gap-finito}. The witness
\eqref{eq:realizaciones-ym-testigo-gap} attains the bound, so it is not
merely a lower bound: it is the first nonzero energy.

The homotopy \eqref{eq:realizaciones-ym-homotopia-terminal} contracts the
cellular complement and leaves \(FP\) intact; thus the corona introduces no
second vacuum and does not erase the carry. Isometric compatibility transports the
vacuum and its complement to the next level without mixing them with the scale
residual. The law
\eqref{eq:realizaciones-ym-qm} shows that the spectral quotient per physical
unit is exactly \(3\kappa_{\rm av}\) at every level. The forms
\eqref{eq:realizaciones-ym-forma-limite}--
\eqref{eq:realizaciones-ym-gap-forma-limite} transport the bound and the witness
to the self-adjoint generator of the limit, so the gap remains exact.
Reconstruction by Osterwalder--Schrader
positivity identifies the Hamiltonian with the generator of this
same transfer, preserves the simple vacuum and produces
\eqref{eq:realizaciones-ym-gap-limite}. The final dimensional chart transforms
the gap into the indicated collective mass.
\end{proof}

\section*{Gauge compression and physical gap witness}

Let
\[
 \widehat{\mathscr H}_m
 :=L^2(\widehat X_m,\widehat\mu_m^{\rm ov})
\]
be the space of the complete material realization. In addition to the usual gauge
action on links and frames, each collar \(c\) carries the finite incidence
action
\begin{equation}\label{eq:amp-ym-accion-gauge-incidencia}
 q_c\longmapsto q_c+Bx_c,
 \qquad x_c\in\mathbb F_3^4,
\end{equation}
where the six rows of \(B\) are indexed by
\((01,02,03,12,13,23)\). The product of both actions preserves
\(\widehat\mu_m^{\rm ov}\). With \(\mathcal C_m\) denoting the finite set of
collars present at this level, write
\[
 \mathcal G_m^{\rm full}
 :=G^{V(\Lambda_m)}
   \times\prod_{c\in\mathcal C_m}\mathbb F_3^4.
\]
Its Haar average defines the orthogonal
projection
\begin{equation}\label{eq:amp-ym-proyeccion-fisica-haar}
 \mathbb P_m^{\rm phys}F
 :=\int_{\mathcal G_m^{\rm full}}F(g^{-1}\!\cdot\,)\,dg,
 \qquad
 \widehat{\mathscr H}_m^{\rm phys}
 :=\operatorname{Ran}\mathbb P_m^{\rm phys}.
\end{equation}
This is a compression onto the gauge subalgebra of the complete state; the
marginal forgetting the incidence coordinates is a different map and is not
confused with \(\mathbb P_m^{\rm phys}\).

\begin{lemma}[Reduction of the Hamiltonian to the physical subalgebra]
\label{lem:amp-ym-compresion-fisica}
The projection \(\mathbb P_m^{\rm phys}\) reduces the generator and its
semigroup. The refinement inclusions simultaneously reduce the
physical subalgebras:
\begin{equation}\label{eq:amp-ym-compresion-entrelazada}
 \begin{gathered}
 [\mathbb P_m^{\rm phys},\widehat H_m^{\rm ov}]=0,
 \qquad
 [\mathbb P_m^{\rm phys},e^{-t\widehat H_m^{\rm ov}}]=0,\\
 \mathbb P_{m+1}^{\rm phys}\widehat J_m
 =\widehat J_m\mathbb P_m^{\rm phys},\qquad
 H_m^{\rm phys}
 :=\widehat H_m^{\rm ov}\big|_{\widehat{\mathscr H}_m^{\rm phys}},\\
 H_{m+1}^{\rm phys}\widehat J_m
 =\widehat J_mH_m^{\rm phys}.
 \end{gathered}
\end{equation}
\end{lemma}

\begin{proof}
The physical generator on links is gauge-equivariant by the
bi-invariance of Haar measure. In the product chart, \(E_{q_c}\) is the uniform
expectation on \(\mathbb F_3^6\); by invariance of the uniform measure under
the translations of \eqref{eq:amp-ym-accion-gauge-incidencia}, it commutes with
that action. The expectations of the marking subtails act on
disjoint factors. Therefore each summand of
\[
 \widehat H_m^{\rm ov}
 =H_m^{\rm ov}\otimes I
 +3\kappa_{\rm av}\sum_c
 \bigl[(I-E_{q_c})+(I-E_{u_c^{\rm P}})\bigr]
\]
commutes with the average \eqref{eq:amp-ym-proyeccion-fisica-haar}. The
functional calculus gives commutation with the semigroup. Finally,
\(\widehat J_m\) preserves the ancestral factors and inserts the constant
in the new factors; averaging before or after this inclusion produces
the same function. This proves the four identities in
\eqref{eq:amp-ym-compresion-entrelazada}.
\end{proof}

\begin{lemma}[Nonzero intertwined gauge witness]
\label{lem:amp-ym-testigo-fisico-q6}
For every collar \(c\), the function
\begin{equation}\label{eq:amp-ym-testigo-fisico-def}
 W_c(q)
 :=\mathbf 1_{\{q_1+\cdots+q_6=0\}}-\frac13
\end{equation}
belongs to \(\widehat{\mathscr H}_m^{\rm phys}\), is nonzero, and satisfies
\begin{equation}\label{eq:amp-ym-testigo-fisico-momentos}
 \mathbb E W_c=0,\qquad
 \|W_c\|^2=\frac29,\qquad
 \mathbb E(W_c^3)=\frac2{27}.
\end{equation}
Moreover,
\begin{equation}\label{eq:amp-ym-testigo-fisico-gap}
 H_m^{\rm phys}W_c=3\kappa_{\rm av}W_c,\qquad
 H_{m+1}^{\rm phys}\widehat J_mW_c
 =3\kappa_{\rm av}\widehat J_mW_c.
\end{equation}
Consequently, the value \(3\kappa_{\rm av}\) belongs to the spectrum of the
physical restriction, not merely to that of a material dilation.
\end{lemma}

\begin{proof}
The gauge action on links acts on the first factor and fixes
\(W_c\). For the incidence action, each \(x_i\) appears in three of the
six coordinates, hence
\[
 \sum_{i<j}(Bx)_{ij}=3\sum_i x_i=0
 \quad\text{in }\mathbb F_3.
\]
The condition in \eqref{eq:amp-ym-testigo-fisico-def} is therefore
invariant; it is also invariant under permutations of the six incidences.
Thus \(\mathbb P_m^{\rm phys}W_c=W_c\).

Under the uniform measure, the sum of the six trits is uniform on
\(\mathbb F_3\). The function takes the value \(2/3\) with probability \(1/3\)
and \(-1/3\) with probability \(2/3\), giving
\eqref{eq:amp-ym-testigo-fisico-momentos}. The first summand of
\(\widehat H_m^{\rm ov}\) annihilates \(W_c\); for \(c'\ne c\),
\((I-E_{q_{c'}})W_c=0\), whereas
\(E_{q_c}W_c=0\). All marking expectations annihilate the
corresponding difference because \(W_c\) is constant on these subtails.
What remains is exactly
\(\widehat H_m^{\rm ov}W_c=3\kappa_{\rm av}W_c\). The first identity in
\eqref{eq:amp-ym-testigo-fisico-gap} follows from reduction; the second,
from the intertwining in \eqref{eq:amp-ym-compresion-entrelazada}.
\end{proof}

\section*{Proprietary semigroup, vacuum and exact spectral bound}

On the physical subalgebra, write
\[
 D_m^{\rm YM}:=H_m^{\rm phys},\qquad
 K_{m,t}:=\exp(-tD_m^{\rm YM}),\qquad
 P_{\Omega_m}:=|\Omega_m\rangle\langle\Omega_m|.
\]
The Hoeffding decomposition of the same generator and the functional calculus
give, without changing the measure or process,
\begin{equation}\label{eq:amp-ym-owner-semigrupo-vacio}
 K_{m,t}P_{\Omega_m}=P_{\Omega_m},
 \qquad
 \bigl\|K_{m,t}(I-P_{\Omega_m})\bigr\|
 \leq e^{-3\kappa_{\rm av}t}.
\end{equation}
The witness in \eqref{eq:amp-ym-testigo-fisico-def} is nonzero and attains the
bound:
\begin{equation}\label{eq:amp-ym-owner-testigo-wc}
 D_m^{\rm YM}W_c=3\kappa_{\rm av}W_c,
 \qquad \|W_c\|^2=\frac29.
\end{equation}
Thus the vacuum is simple, and the gap is not merely bounded below
but is exactly
\begin{equation}\label{eq:amp-ym-owner-brecha-exacta}
 \Delta_{\rm YM}^{\rm HMT}=3\kappa_{\rm av}>0.
\end{equation}
The semigroup, vacuum projector, contraction of the complement and
\(W_c\) belong to the same projective family and constitute the
proprietary chain of the closure.

\section*{Euclidean reconstruction, Schwinger fields, and curvature reader}

The common projective measure in the theorem determines isometric inclusions
\(\widehat J_m\) between the level spaces. For each of the four
coordinate reflections \(\theta_\mu\), \(\mu=1,\ldots,4\), there is a
factorization
\begin{equation}\label{eq:amp-ym-cuatro-os}
 \int\overline{F(\theta_\mu y)}F(y)\,
 d\widehat\Omega_m(y)
 =\|\mathcal B_{m,\mu}F\|_2^2\geq0.
\end{equation}
The four forms originate in the same measure and semigroup. The
OS connecting maps
\[
 \mathcal J_{m,\mu}^{\rm OS}
 =\mathcal V_{m+1,\mu}^{-1}\widehat J_m\mathcal V_{m,\mu}
\]
intertwine their Hamiltonians. In the colimit,
\begin{equation}\label{eq:amp-ym-hamiltoniano-comun}
 \mathcal V_{\infty,\mu}H_{{\rm OS},\infty}^{(\mu)}
 \mathcal V_{\infty,\mu}^{-1}
 =\widehat H_\infty,
 \qquad \mu=1,\ldots,4.
\end{equation}

The limiting form is diagonalized by the Hoeffding decomposition:
\begin{equation}\label{eq:amp-ym-forma-limite}
 \mathcal E_\infty(u)=3\kappa_{\rm av}
 \sum_{S\Subset\mathcal I_\infty}|S|\,\|u_S\|^2.
\end{equation}
The partial sums provide a recovery sequence, and
semicontinuity yields the lower limit. Mosco convergence of the
closed forms preserves the one-dimensional kernel and spectral threshold. The
explicit vector \(W_c\) satisfies
\begin{equation}\label{eq:amp-ym-testigo-gap}
 \widehat H_mW_c=3\kappa_{\rm av}W_c,
 \qquad \|W_c\|^2=\frac29,
 \qquad \mathbb E(W_c^3)=\frac2{27},
\end{equation}
so the lower bound is attained and the gap is exactly
\(3\kappa_{\rm av}\).

\begin{theorem}[Euclidean completion and publication of curvature]
\label{thm:amp-ym-terminacion}
The projective family of theorem \ref{thm:realizaciones-ym} determines a
strongly continuous action of \(E(4)\), a nonscalar local gauge net,
compatible Schwinger functionals and a field publication in
\(H^{-s}_{\rm loc}(\mathbb R^4)\) for \(s>2\). The ordered surface
product of the links defines clover readers and \(F^2\); in the smooth
domain, the action reader converges to
\begin{equation}\label{eq:amp-ym-f2}
 \mathcal S_{\rm YM}(A)
 =\frac1{4g_R^2}\int_{\mathbb R^4}\|F_A(x)\|^2\,d^4x.
\end{equation}
These publications preserve the simple vacuum and exact gap of the same
common Hamiltonian.
\end{theorem}

\begin{proof}
The cofinal Gram matrices of the smoothed readers and commutativity of the
rotation and refinement squares preserve the null ideals; their
colimit constructs the strong representation of \(E(4)\) and the local net.
Martingale estimates for the link increments give tightness in
\(H^{-s}_{\rm loc}\), \(s>2\), and the mixed characteristic functionals identify the
Schwinger functionals with the same cylindrical process. The oriented
product around each plaquette constructs the clover. Its expansion in
the smooth regime and strong convergence of the marginal reader give
\eqref{eq:amp-ym-f2}. None of these operations changes the measure,
semigroup, or limiting form used to obtain
\eqref{eq:amp-ym-testigo-gap}; hence the vacuum and gap are preserved.
\end{proof}

\section*{Exact identification of the gauge limit}

Identification of the theory does not rely on four reflections
generating the Euclidean group by themselves. The reflections provide
OS positivity; the continuous action originates, separately, in the screen
groupoid and refinement compatibility. For
\(g=(g_v)_{v\in\Lambda_m}\in G^{\Lambda_m}\), the action on a link
variable is
\begin{equation}\label{eq:amp-ym-accion-gauge-enlace}
 (g\cdot U)_e=g_{s(e)}U_e g_{t(e)}^{-1}.
\end{equation}
Haar invariance, triangular disintegration of the kernel, and
cancellation of interior edges prove
\begin{equation}\label{eq:amp-ym-invariancia-gauge-medida}
 (g\cdot)_*\widehat\mu_m^{\rm ov}=\widehat\mu_m^{\rm ov},
 \qquad
 K_{m,t}^{\rm ov}(g\cdot U,g\cdot dV)
 =K_{m,t}^{\rm ov}(U,dV).
\end{equation}
By differentiation on the cylindrical core, every vertical generator
\(X\) satisfies the Ward identity
\begin{equation}\label{eq:amp-ym-ward}
 \int \nabla_XF\,d\widehat\mu_m^{\rm ov}=0.
\end{equation}

Curvature is obtained from the same link system. If a coarse
plaquette \(p\) is subdivided into its 81 subplaquettes \(p'\), then
\begin{equation}\label{eq:amp-ym-producto-superficial}
 \operatorname{Hol}^{(m)}_{p}
 =\prod_{p'\subset p}^{\longrightarrow}
 T_{p'}\operatorname{Hol}^{(m+1)}_{p'}T_{p'}^{-1}.
\end{equation}
The contributions of every interior edge appear with opposite
orientations and cancel exactly.  The expression is gauge-covariant and
commutes with the inclusions \(\widehat J_m\).  In a smooth chart one defines
\begin{equation}\label{eq:amp-ym-clover-definido}
 F_{{\rm cl},m}(x)
 =\frac1{4a_m^2}\sum_{p\ni x}
 \operatorname{Ad}_{T_p}\log\operatorname{Hol}_{p},
 \qquad
 F_{{\rm cl},m}=F_A+O(a_m^2).
\end{equation}
Hence, for \(A\in C^4_{\rm loc}\) and
\(a_m^2/g_m^2\to0\), Riemann summation gives
\begin{equation}\label{eq:amp-ym-f2-convergencia-calculada}
 \frac{a_m^4}{4g_m^2}
 \sum_{x\in\Lambda_m}\|F_{{\rm cl},m}(x)\|^2
 \longrightarrow
 \frac1{4g_R^2}\int_{\mathbb R^4}\|F_A(x)\|^2\,d^4x.
\end{equation}

\begin{theorem}[Euclidean gauge process and curvature reader for
Yang--Mills]
\label{thm:amp-ym-identificacion-gauge}
The limit constructed in Theorem
\ref{thm:realizaciones-ym} is a Euclidean gauge process in dimension
\(4\) for the fixed compact, connected, simple group. It has
local gauge covariance, a nonscalar local net, reflection
positivity, a strongly continuous representation of \(E(4)\),
compatible Schwinger functionals, gauge-covariant curvature with classical
limit \eqref{eq:amp-ym-f2-convergencia-calculada}, a simple vacuum, and spectral
gap
\(3\kappa_{\rm av}>0\).
\end{theorem}

\begin{proof}
The identities \eqref{eq:amp-ym-invariancia-gauge-medida}--
\eqref{eq:amp-ym-ward} construct gauge symmetry and its Ward
identities at every level; their naturality transports them to the limit. The
uniform locality of the kernels produces the isotonic local net. The
four forms \eqref{eq:amp-ym-cuatro-os}, together with the screen
groupoid action on the Weyl core, produce OS positivity and the strong
action of \(E(4)\), respectively. Tightness in
\(H^{-s}_{\rm loc}\), \(s>2\), identifies a unique limiting law and its
Schwinger functionals. Finally,
\eqref{eq:amp-ym-producto-superficial}--
\eqref{eq:amp-ym-f2-convergencia-calculada} identify the field reader
and its classical limit, whereas
\eqref{eq:amp-ym-forma-limite}--
\eqref{eq:amp-ym-testigo-gap} prove a simple vacuum and exact gap for the
same law.

The measure is constructively determined by its projective
marginals, reversible kernel and cylindrical functionals. The clover
reader does not reweight that measure: it publishes a function of the same process and its
classical limit.
\end{proof}

\begin{proposition}[Non-governing Gibbs--action comparison check]
\label{prop:amp-ym-puente-accion-requerido}
The projective law, its simple vacuum and its exact gap have already been constructed
by \eqref{eq:amp-ym-owner-semigrupo-vacio}--
\eqref{eq:amp-ym-owner-brecha-exacta}. As an optional external comparison of
Gibbs--action type, a parametrization by \(g_R\) may study a family
\[
 g_R\longmapsto
 \bigl(\widehat\mu_{m,g_R}^{\rm ov},H_{m,g_R}^{\rm phys}\bigr)
\]
and a relation between the measure or its Dirichlet form and the action
functional, for example
\begin{equation}\label{eq:amp-ym-identidad-accion-requerida}
 \frac{d\widehat\mu_{m,g_R}^{\rm ov}}{d\nu_m}(U)
 =Z_{m,g_R}^{-1}e^{-S_{m,g_R}^{F^2}(U)},
 \qquad
 \mathfrak h_{m,g_R}[F]
 =\langle F,H_{m,g_R}^{\rm phys}F\rangle,
\end{equation}
with refinement compatibility and passage to the limit. Here \(\nu_m\) may
be the product Haar measure of the finite lattice. This check is designated
\texttt{NON\_GOVERNING\_CHECK}: it is not a premise of HMT gauge theory or
its gap.
\end{proposition}

\begin{proof}
In \eqref{eq:amp-ym-f2-convergencia-calculada}, varying \(g_R\) rescales the
curvature reader, whereas the formulas constructing
\(\widehat\mu_m^{\rm ov}\), \(\widehat H_m^{\rm ov}\), and the gap
\(3\kappa_{\rm av}\) do not contain \(g_R\). They are therefore independent
constructions. The compression
\eqref{eq:amp-ym-compresion-entrelazada} and witness
\eqref{eq:amp-ym-testigo-fisico-gap} prove the physical persistence of the
gap in the already constructed process; the check only subsequently compares an
additional dynamical dependence on \(g_R\).
\end{proof}

The sequence proved in this section is
\[
 \text{common measure}\to\text{four OS}\to\text{Common Hamiltonian}
 \to E(4),\ \text{Schwinger},\ H^{-s}_{\rm loc}
 \to\text{clover and }F^2.
\]
This sequence proves a four-dimensional Yang--Mills theory with a simple
vacuum and an exact positive gap \(3\kappa_{\rm av}\) in the registered HMT
domain. The identity
\eqref{eq:amp-ym-identidad-accion-requerida} remains outside the governing
chain as \texttt{NON\_GOVERNING\_CONTROL}.


\section*{Wilson reading}

The Wilson observable is obtained after the generator. For an Archimedean link
chart,
\begin{equation}\label{eq:realizaciones-ym-wilson}
 A_e=\sum_a\operatorname{ev}_{\mathbb R}(x_{e,a})T_a,
 \qquad
 U_e=e^{a_mA_e},
 \qquad
 W_\gamma=\operatorname{Tr}\!\prod_{e\in\gamma}U_e.
\end{equation}
The reading preserves the representation, the intertwining operator and the
order of the holonomy. It serves as a terminal observable of the published theory; it does not
retrospectively define the measure, the semigroup or the gap.

\section*{Necessity of the hypotheses and exact obstructions}

\begin{ablation}[Necessity of the common dynamics and measure]
\label{abl:realizaciones-ym}
Theorem \ref{thm:realizaciones-ym} does not hold if the
nine subphases are replaced by independent updates, different measures or
kernels are used for the four reflections, uniform physical locality
is lost, \(\widehat J_m\) ceases to intertwine transfer and reflection,
carry or memory is erased, \(P_{\Omega_m}\) is removed,
\(Q_m^{\rm phys}\) is identified with a scale residual,
\(q_m\leq q_*<1\) is required in lattice units, the complete corona is replaced
by its \(271\) vertices, \(FP\) is erased, or
\eqref{eq:realizaciones-ym-wilson} is used to select the generator.
\end{ablation}

\begin{proof}
Different updates, kernels, or measures destroy the simultaneous
factorization \eqref{eq:realizaciones-ym-os}. Without locality or intertwining,
positivity at one stage does not define a compatible local net. Erasing carry,
memory, or \(FP\) removes information preserved by
\eqref{eq:realizaciones-ym-homotopia-terminal}. Using only the
vertices changes the relative complex and can create zero modes. Confusing
the two complements allows the physical sector to disappear through a
purely scalar operation. A fixed bound in lattice units contradicts
\eqref{eq:realizaciones-ym-qm}. Finally, Wilson is only a reading of the
constructed state and does not by itself prove either
\eqref{eq:realizaciones-ym-os} or
\eqref{eq:realizaciones-ym-gap-finito}.
\end{proof}

The independent residuals are: failure of
\(\widehat T_{m,\nu}^{\rm OS}=e^{-a_m\widehat H_m^{\rm ov}}\), failure of the
intertwining \eqref{eq:realizaciones-ym-semigrupo-entrelazado}, a negative
OS form, a kernel vector orthogonal to \(\Omega_m\), or a nonconstant
Hoeffding sector with energy below \(3\kappa_{\rm av}\).

\section*{Scope of the gauge construction}

The proof combines the joint connection, terminal preservation, the
double reading, the local circuit, the common measure, the four reflection
factorizations, intertwined refinement and the positive generator. In the
declared regular gauge domain,
locality, four-dimensional Yang--Mills theory, the simple vacuum and the
collective gap are exact results. The Osterwalder--Schrader
reconstruction and the Wilson observable are
subsequent publications of the same state; neither acts as a
retrospective selector of the measure, the semigroup, the vacuum or the
gap. The optional comparison with a family parametrized by \(g_R\) is studied
through \eqref{eq:amp-ym-identidad-accion-requerida}; that check does not
govern or reopen the already proved measure, semigroup, vacuum or
gap.

\section*{Comparison of the spectral, solenoidal and gauge mechanisms}

The preceding constructions share the architecture
\begin{equation}\label{eq:realizaciones-sintesis-tres}
 \begin{array}{ccl}
 \text{Riemann--Weil}
&:& \text{finite memory}\to\text{double cofinal publication}\\
&&{}\to B_W=\mathscr I_{\rm GW}\geq0
      \to\text{reversible criterion},\\[1mm]
 \text{Navier--Stokes}
&:& U^{\rm own}\to\text{telescoping and ledger}\\
&&{}\to\text{uniform root and exact return}
      \to\text{parabolic limit},\\[1mm]
 \text{OS gauge process}
&:& \text{terminal root}\to\text{reflection positivity}\\
&&{}
      \to\text{collective gap}.
 \end{array}
\end{equation}
The joint nonadic connection does not replace the specialized lemmas: it
organizes them and preserves their data upon refinement. The conclusions of
Riemann--Weil, Navier--Stokes and Yang--Mills follow from their respective
proprietary chains and already proved final criteria. The residual,
Gram and action-density comparators remain subsequent
non-governing controls. This separation prevents both loss of
genealogy and unjustified transfer of a proof from one domain
to another.

% Residencia sucesora única de Hodge y Birch--Swinnerton-Dyer. Este input es
% parte del capítulo, no una importación duplicada de los manuscritos
% fuentes.
% 04 — FRAGMENTO SUCESOR HODGE Y BIRCH--SWINNERTON-DYER
% =======================================================
% Inserción autónoma prevista tras el marcador de integración del Libro IV.
% No importa fuentes probatorios ni modifica las tres realizaciones
% precedentes.
%
% Trazabilidad técnica no narrativa:
%   Hodge: CHI-II-HDG-01/02/03, CHI-IV-HDG-01/C01.
%   BSD:   CHI-II-BSD-01/02/03, CHI-IV-BSD-01/02/C01.
%   Fuentes: celda APP--Mayer--Vietoris, totalizador coend y controles
%   elípticos; producto fibrado genealógico de Manin, publicación acoplada
%   Kato--PT, dualidad ortogonal reducida, relleno finita--singular y regulador
%   Bockstein derivado. Los hashes y estados de ejecución permanecen en el
%   expediente técnico de la edición sucesora.

\ifdefined\HMTSucesoraStructureTrace
  \typeout{HMT-SUCCESSOR 04: Hodge and Birch--Swinnerton-Dyer realizations}
\fi

\chapter{Completeness of enriched cohomological history and algebraicity of rational Hodge classes}
\chaptermark{Cohomological completeness and Hodge algebraicity}

\subsection*{Object, domain and starting data}

Let \(X\) be a smooth projective variety over \(\mathbb C\) and let
\(p\geq0\). We write
\begin{equation}\label{eq:hodge-suc-dominio}
 \operatorname{Hdg}^{p}(X)_{\mathbb Q}
 :=H^{2p}(X,\mathbb Q)
   \cap H^{p,p}(X,\mathbb C).
\end{equation}
Let \(\mathbf{Perf}^{\geq p}(X)\) be the idempotent-complete category of
perfect complexes whose support has codimension at least \(p\). In this
section the abbreviation is used
\[
 \operatorname{Perf}^{\geq p}(X)_{\mathbb Q}
 :=K_0\!\left(\mathbf{Perf}^{\geq p}(X)\right)
   \otimes_{\mathbb Z}\mathbb Q;
\]
thus, \(R_{\mathrm{Hdg}}\) publishes the class of the constructed perfect
complex. The starting point is not a bare class of
\eqref{eq:hodge-suc-dominio}, but the limit of surviving histories
\begin{equation}\label{eq:hodge-suc-xchi}
 X_{\chi}(X,p)
 =\varprojlim_{K}\operatorname{Surv}_{K}(X,p).
\end{equation}
Each element of \(X_{\chi}(X,p)\) preserves sheet, orientation, quotient,
carry, boundary, route and memory. In particular, the return of the joint
nonadic connection satisfies
\begin{equation}\label{eq:hodge-suc-retorno}
 \Gamma _9=\operatorname{Hol}_{C_9}(\nabla_{\mathrm{TPK}}),
 \qquad g(k+9)=g(k),
 \qquad \widehat x(k+9)\neq\widehat x(k)
 \quad\text{in general}.
\end{equation}
The equality preserves the phase; the inequality preserves the depth of the
history. This distinction rules out interpreting a return as a
reinitialization.

Route composition is governed by the carry cocycle
\begin{equation}\label{eq:hodge-suc-carry}
 \operatorname{Carr}(\gamma _2\gamma _1)
 =\operatorname{Carr}(\gamma _2)H(\gamma _1)
  +Y(\gamma _2)\operatorname{Carr}(\gamma _1),
\end{equation}
and accumulated memory by
\begin{equation}\label{eq:hodge-suc-memoria}
 S_0=\sum_{r<N}(M_9^{*})^{r}D_9^{*}D_9M_9^{r}
     +(M_9^{*})^{N}S_0M_9^{N}.
\end{equation}
Equations \eqref{eq:hodge-suc-xchi}--\eqref{eq:hodge-suc-memoria}
are data constructed in Part II. They are specialized here; they are not
reconstructed from the cohomological codomain.

\subsection*{Local APP--Mayer--Vietoris cell}

Let \(Z,W\subset X\) be closed incidences with a
Tor-independent intersection and oriented supports of the declared codimension. For
integers \(a,b\geq1\), define
\begin{equation}\label{eq:hodge-suc-kappa}
 \begin{aligned}
 \kappa_{a,b}:\mathcal O_Z^{a}\oplus\mathcal O_W^{b}
 &\longrightarrow \mathcal O_{Z\cap W}^{ab},\\
 ((u_i)_{i=1}^{a},(v_j)_{j=1}^{b})
 &\longmapsto
 \bigl(u_i|_{Z\cap W}-v_j|_{Z\cap W}\bigr)_{i,j}.
 \end{aligned}
\end{equation}
On a fiber of \(Z\cap W\), the diagonal kernel of
\(\kappa_{a,b}\) has rank one and its cokernel has rank
\begin{equation}\label{eq:hodge-suc-defecto-rango}
 ab-(a+b-1)=(a-1)(b-1).
\end{equation}
If \((\rho_{\Sigma},c_{\Sigma})\) and
\((\rho_{\Pi},c_{\Pi})\) are, respectively, the residues and carries of
the additive and multiplicative sheets of APP, the same quantity is expressed by
\begin{equation}\label{eq:hodge-suc-celda-carry}
 (a-1)(b-1)
 =\rho_{\Pi}-\rho_{\Sigma}+1
  +9\bigl(c_{\Pi}-c_{\Sigma}\bigr).
\end{equation}
Thus, the local defect is not an external correction: it is the publication in
\(\operatorname{Perf}\) of the difference between the two sheets with its carry.

Let \(\tau\) be the exchange of the two summands of the source and let \(P\)
be the transposition of the target indices. The TRIT orientation fixes
\begin{equation}\label{eq:hodge-suc-orientacion-kappa}
 \kappa_{b,a}\tau=-P\kappa_{a,b}.
\end{equation}
The sign of \eqref{eq:hodge-suc-orientacion-kappa} is part of the oriented
object. Omitting it would identify opposite routes and destroy the
naturality of the localized character.

\begin{proposition}[Perfect realization of the local cell]
\label{prop:hodge-suc-celda-perf}
The cone
\begin{equation}\label{eq:hodge-suc-cone-kappa}
 \mathcal C_{a,b}(Z,W):=\operatorname{Cone}(\kappa_{a,b})
\end{equation}
is a perfect complex with the prescribed support. Its localized class
preserves the rank defect \eqref{eq:hodge-suc-defecto-rango}, the carry
\eqref{eq:hodge-suc-celda-carry} and the orientation
\eqref{eq:hodge-suc-orientacion-kappa}.
\end{proposition}

\begin{proof}
Tor-independence identifies the derived restriction to the intersection with the
ordinary restriction used in \eqref{eq:hodge-suc-kappa}. The sources and
target are perfect on their supports; by stability of
\(\operatorname{Perf}\) under cones, so is
\(\mathcal C_{a,b}(Z,W)\). The fiber-by-fiber calculation gives
\eqref{eq:hodge-suc-defecto-rango}; APP division of the sum and
product data gives \eqref{eq:hodge-suc-celda-carry}. Finally, exchanging
\(Z,a\) with \(W,b\) changes each difference to its opposite, proving
\eqref{eq:hodge-suc-orientacion-kappa}. Since the localized Chern character
is additive in distinguished triangles, the three pieces of information pass to the
class of the cone.
\end{proof}

\subsection*{Totalization of histories and preservation of the terminal}

Let \(J_N\) be the finite category of compatible histories at depth
\(N\), let \(\Gamma\) be the diagram of incidence complexes obtained from
\eqref{eq:hodge-suc-cone-kappa} and let \(\mathsf W\) be the right module that
records route orientations and weights. The derived totalization is the
coend
\begin{equation}\label{eq:hodge-suc-coend}
 K_{X,p,N}
 :=\mathsf W\otimes^{\mathbf L}_{\mathbb Z[J_N]}\Gamma.
\end{equation}
The survival links form a compatible system in \(N\). The
holonomy \(\Gamma_9\) induces an endomorphism
\(U_{\Gamma_9}\) of the totalizer and the terminal object is preserved as
\begin{equation}\label{eq:hodge-suc-terminal}
 \mathcal T_{X,p}
 :=\operatorname{Cone}\bigl(I-U_{\Gamma_9}\bigr).
\end{equation}
The cone \eqref{eq:hodge-suc-terminal} prevents totalization from confusing
a repeated phase with the same history: the visible value may return,
but the terminal difference retains the memory displacement.

Compatibility between \eqref{eq:hodge-suc-coend} and the limit
\eqref{eq:hodge-suc-xchi} defines the realization morphism
\begin{equation}\label{eq:hodge-suc-realizacion}
 R_{\mathrm{Hdg}}:
 X_{\chi}(X,p)\longrightarrow
 \operatorname{Perf}^{\geq p}(X)_{\mathbb Q}.
\end{equation}
This morphism is fixed at the source: it does not receive a Hodge class
or a representative it must recognize as input.

\begin{theorem}[Hodge publication identity]
\label{thm:hodge-suc-identidad}
On \(X_{\chi}(X,p)\) the identity commutes
\begin{equation}\label{eq:hodge-suc-identidad}
 \operatorname{cl}_{X,p}\circ
 \operatorname{ch}^{\mathrm{loc}}_{p}\circ
 R_{\mathrm{Hdg}}
 =\operatorname{ev}^{\mathrm{Hdg}}_{\infty,X,p}.
\end{equation}
Here
\(\operatorname{ch}^{\mathrm{loc}}_{p}\) takes the codimension
\(p\) component of the localized character and
\(\operatorname{cl}_{X,p}\) is its cohomological publication.
\end{theorem}

\begin{proof}
In a cell, the equality is the compatibility of the localized character with
the restriction, the Mayer--Vietoris triangle and the cohomological class. The
orientation is fixed by \eqref{eq:hodge-suc-orientacion-kappa};
therefore, no sign remains to be chosen. Additivity of the Chern character
transports the equality to each cone \eqref{eq:hodge-suc-cone-kappa}.
Naturality with respect to the arrows of \(J_N\) makes it descend to the coend
\eqref{eq:hodge-suc-coend}. The carry law
\eqref{eq:hodge-suc-carry} guarantees its compatibility under composition and
\eqref{eq:hodge-suc-terminal} preserves the terminal term. Passing to the
projective limit yields two natural transformations with equal
components on every finite cylinder; hence they coincide and
\eqref{eq:hodge-suc-identidad} follows.
\end{proof}

The completeness theorem of the atlas of enriched histories, applied to
this specialization, establishes
\begin{equation}\label{eq:hodge-suc-completitud}
 \operatorname{ev}^{\mathrm{Hdg}}_{\infty,X,p}
 \bigl(X_{\chi}(X,p)\bigr)
 =\operatorname{Hdg}^{p}(X)_{\mathbb Q}.
\end{equation}
This is the explicit starting structure of the Hodge specialization.

% Ampliación sustantiva de la completitud coinductiva Hodge.
% Módulo destinado al capítulo Hodge de la Parte IV.

\section{Coinductive unfolding of Hodge completeness}
\label{sec:amp-hodge-completitud-desplegada}

The completeness equality
\eqref{eq:hodge-suc-completitud} can be unfolded without starting from an algebraic
class or introducing into the source the conclusion one wishes to publish. The
object being completed is the system of enriched histories; the Hodge
class intervenes only as compatible evaluation data in its
finite truncations.

\subsection{Finite observation systems and compatible fibers}
\label{subsec:amp-hodge-sistemas-finitos}

For each depth \(N\geq0\), let \(J_N(X,p)\) be the finite category of
incidence charts, oriented routes and memory terminals of depth at
most \(N\). We write
\begin{equation}\label{eq:amp-hodge-surv-n}
 \mathscr S_N(X,p):=\operatorname{Surv}_{J_N}(X,p),
 \qquad
 \rho_{N+1,N}:\mathscr S_{N+1}(X,p)\longrightarrow\mathscr S_N(X,p)
\end{equation}
for the rationalized space of surviving states and its restriction. The
finite readers take values in the module of cylinder observations
\(\mathscr H_N^p(X)\):
\begin{equation}\label{eq:amp-hodge-evaluacion-finita}
 e_N^{\mathrm{Hdg}}:\mathscr S_N(X,p)\longrightarrow
 \mathscr H_N^p(X),
 \qquad
 q_{N+1,N}:\mathscr H_{N+1}^p(X)\longrightarrow\mathscr H_N^p(X).
\end{equation}
Here \(\mathscr H_N^p(X)\) records only the evaluations on the incidences
present in \(J_N\), but preserves their rational coefficients and
orientation. If
\(q_N:\operatorname{Hdg}^p(X)_{\mathbb Q}\to\mathscr H_N^p(X)\) is the
restriction to the same finite set of observations, we have
\begin{equation}\label{eq:amp-hodge-cuadrado-restriccion}
 e_N^{\mathrm{Hdg}}\rho_{N+1,N}
 =q_{N+1,N}e_{N+1}^{\mathrm{Hdg}},
 \qquad
 q_{N+1,N}q_{N+1}=q_N.
\end{equation}

For \(\alpha\in\operatorname{Hdg}^p(X)_{\mathbb Q}\), define the compatible finite
fiber
\begin{equation}\label{eq:amp-hodge-fibra-alpha}
 \mathscr F_N(\alpha):=
 \left\{s_N\in\mathscr S_N(X,p):
 e_N^{\mathrm{Hdg}}(s_N)=q_N(\alpha)\right\}.
\end{equation}
The word \emph{finite} refers to the depth and the incidence
diagram, not to a loss of rational coefficients or to a selection
of a representative from \(X\) alone.

\begin{proposition}[Finite extension with memory]
\label{prop:amp-hodge-extension-finita}
For every \(\alpha\) as in \eqref{eq:amp-hodge-fibra-alpha}, the fibers
\(\mathscr F_N(\alpha)\) are nonempty and the restrictions induce surjective
maps
\begin{equation}\label{eq:amp-hodge-mittag-leffler}
 \rho_{N+1,N}:\mathscr F_{N+1}(\alpha)\twoheadrightarrow
 \mathscr F_N(\alpha).
\end{equation}
The extension preserves sheet, TRIT sign, quotient, carry, support, boundary and
memory terminal.
\end{proposition}

\begin{proof}
The assertion is the Hodge specialization of the extension theorem for the atlas
constructed in Part II. Its mechanism should be made visible. Let
\(s_N\in\mathscr F_N(\alpha)\) and choose a prolongation of charts
\(\widetilde s_{N+1}\) preserving the history of \(s_N\). Its discrepancy at
the new level is
\begin{equation}\label{eq:amp-hodge-discrepancia}
 \delta_{N+1}:=q_{N+1}(\alpha)
 -e_{N+1}^{\mathrm{Hdg}}(\widetilde s_{N+1}),
 \qquad q_{N+1,N}(\delta_{N+1})=0,
\end{equation}
where the last equality comes from
\eqref{eq:amp-hodge-cuadrado-restriccion}. The
APP--Mayer--Vietoris cells publish exactly that relative kernel through the
cones of the morphisms \(\kappa_{a,b}\); their defect
\((a-1)(b-1)\) is lifted with the residue and carry of
\eqref{eq:hodge-suc-celda-carry}. By the finite extension theorem there is a
relative correction \(\eta_{N+1}(\delta_{N+1})\), supported on the new
charts, such that
\begin{equation}\label{eq:amp-hodge-extension-operativa}
 \rho_{N+1,N}\eta_{N+1}(\delta_{N+1})=0,
 \qquad
 e_{N+1}^{\mathrm{Hdg}}\eta_{N+1}(\delta_{N+1})=\delta_{N+1}.
\end{equation}
Therefore
\begin{equation}\label{eq:amp-hodge-ext}
 \operatorname{Ext}_{N+1}(s_N,\alpha):=
 \widetilde s_{N+1}+\eta_{N+1}(\delta_{N+1})
 \in\mathscr F_{N+1}(\alpha)
\end{equation}
restricts to \(s_N\). The orientation
\(\kappa_{b,a}\tau=-P\kappa_{a,b}\) fixes the sign, the cocycle law
\eqref{eq:hodge-suc-carry} fixes the composition and
\(\operatorname{Cone}(I-U_{\Gamma_9})\) preserves the terminal. The initial level
is obtained from the based seed atlas. Induction proves nonemptiness and
surjectivity at all levels.
\end{proof}

\subsection{Passage to the limit and order of construction}
\label{subsec:amp-hodge-paso-limite}

\begin{theorem}[Unfolded coinductive completeness]
\label{thm:amp-hodge-completitud-coinductiva}
With the data already constructed in Part II---the systems
\eqref{eq:amp-hodge-surv-n}--\eqref{eq:amp-hodge-evaluacion-finita}, the
compatibility \eqref{eq:amp-hodge-cuadrado-restriccion}, the finite extension
of proposition \ref{prop:amp-hodge-extension-finita}, the separation of
cylinder evaluations and preservation of the terminal---for each
\(\alpha\in\operatorname{Hdg}^p(X)_{\mathbb Q}\), there is an enriched
history
\begin{equation}\label{eq:amp-hodge-xi-limite}
 \xi_\alpha=(s_0,s_1,s_2,\ldots)
 \in\varprojlim_N\mathscr F_N(\alpha)
 \subset X_\chi(X,p)
\end{equation}
which satisfies
\begin{equation}\label{eq:amp-hodge-evaluacion-alpha}
 \operatorname{ev}^{\mathrm{Hdg}}_{\infty,X,p}(\xi_\alpha)=\alpha.
\end{equation}
In particular, completeness is obtained before applying the perfect
realizer and before forming any Chow class.
\end{theorem}

\begin{proof}
Choose \(s_0\in\mathscr F_0(\alpha)\). Once \(s_N\) is constructed, the
surjectivity \eqref{eq:amp-hodge-mittag-leffler} allows us to choose
\(s_{N+1}\in\mathscr F_{N+1}(\alpha)\) with
\(\rho_{N+1,N}s_{N+1}=s_N\). Thus one obtains
\eqref{eq:amp-hodge-xi-limite}. Equivalently, the system of fibers
satisfies the Mittag--Leffler condition and produces no
\(\varprojlim^1\) obstruction term. By the definition of each fiber,
\begin{equation}\label{eq:amp-hodge-evaluaciones-xi}
 e_N^{\mathrm{Hdg}}(s_N)=q_N(\alpha)
 \qquad(N\geq0).
\end{equation}
The separation of cylinder observations identifies the unique element
of the limit of the \(q_N(\alpha)\) with \(\alpha\), which proves
\eqref{eq:amp-hodge-evaluacion-alpha}. The terminal cone prevents replacing the
sequence \((s_N)_N\) with the repetition of a visible root; therefore
\(\xi_\alpha\) preserves the rigidification used in passing to the limit.
\end{proof}

The finite publication is arranged in commutative squares
\begin{equation}\label{eq:amp-hodge-cuadrado-finito}
 \begin{array}{ccc}
 \mathscr S_N(X,p)&\xrightarrow{\ R_{\mathrm{Hdg},N}\ }&
 K_0(\mathbf{Perf}^{\geq p}(X))\otimes\mathbb Q\\[1mm]
 \big\downarrow\scriptstyle e_N^{\mathrm{Hdg}}&&
 \big\downarrow\scriptstyle
   \operatorname{cl}_{X,p}\circ\operatorname{ch}^{\mathrm{loc}}_p\\[1mm]
 \mathscr H_N^p(X)&\xrightarrow{\ \operatorname{pub}_N\ }&
 H^{2p}(X,\mathbb Q).
 \end{array}
\end{equation}
The naturality of these squares with respect to \(\rho_{N+1,N}\) gives, in the
limit, the identity \eqref{eq:hodge-suc-identidad}. The causal order is,
therefore,
\begin{equation}\label{eq:amp-hodge-orden-causal}
 \alpha\longmapsto(q_N\alpha)_N
 \longmapsto\xi_\alpha
 \longmapsto R_{\mathrm{Hdg}}(\xi_\alpha)
 \longmapsto
 \beta_\alpha:=\operatorname{ch}^{\mathrm{loc}}_p
   R_{\mathrm{Hdg}}(\xi_\alpha).
\end{equation}
Only in the last step does \(\beta_\alpha\) appear. Applying the limit square
to \eqref{eq:amp-hodge-evaluacion-alpha} yields
\begin{equation}\label{eq:amp-hodge-publicacion-final}
 \operatorname{cl}_{X,p}(\beta_\alpha)=
 \operatorname{ev}^{\mathrm{Hdg}}_{\infty,X,p}(\xi_\alpha)=\alpha.
\end{equation}

\subsection{Local check on the quadric
\texorpdfstring{\(Q^4\)}{Q4}}
\label{subsec:amp-hodge-q4}

The split quadric
\begin{equation}\label{eq:amp-hodge-q4}
 Q^4=\{x_0x_3+x_1x_4+x_2x_5=0\}\subset\mathbb P^5_{\mathbb C}
\end{equation}
has two based families of maximal planes, with representatives
\(\Pi_+\) and \(\Pi_-\). If
\(a=\operatorname{cl}[\Pi_+]\) and
\(b=\operatorname{cl}[\Pi_-]\), then
\begin{equation}\label{eq:amp-hodge-q4-bases}
 \operatorname{CH}^2(Q^4)=\mathbb Za\oplus\mathbb Zb,
 \qquad
 H^4(Q^4,\mathbb Z(2))\cap H^{2,2}(Q^4)
 =\mathbb Za\oplus\mathbb Zb.
\end{equation}
The based map
\begin{equation}\label{eq:amp-hodge-q4-beta}
 \beta_{Q^4,2}(ra+sb):=r[\Pi_+]+s[\Pi_-]
\end{equation}
satisfies exactly
\begin{equation}\label{eq:amp-hodge-q4-beta-identidades}
 \partial_{\mathrm{mot}}\beta_{Q^4,2}=0,
 \qquad
 \operatorname{cl}_{Q^4}\beta_{Q^4,2}=I.
\end{equation}

For the six ports, let
\begin{equation}\label{eq:amp-hodge-q4-matriz}
 A=\begin{pmatrix}
 2&2&2&1&2&1\\
 2&1&2&2&1&1\\
 1&0&0&1&0&2\\
 0&0&1&1&1&0\\
 2&2&2&0&2&1\\
 2&1&2&1&0&0
 \end{pmatrix},
 \qquad \det A=1,
 \qquad \mathbb A=A\otimes I.
\end{equation}
The division \(x_k\mathbb A=u_k+3x_{k+1}\) and the reduction of
\eqref{eq:amp-hodge-q4-beta} define, gate by gate,
\begin{equation}\label{eq:amp-hodge-q4-uk}
 U_{k,j}:=\overline\beta_{Q^4,2}(u_{k,j}),
 \qquad
 D_{\mathrm{mot}}U_k=0,
 \qquad
 H^0(R_{\mathrm{mot}}\otimes\mathbb F_3)(U_k)=u_k.
\end{equation}
Iterating nine times yields the telescoping identity
\begin{equation}\label{eq:amp-hodge-q4-telescopia}
 x_0\mathbb A^9
 =R_{\mathrm{mot}}\!\left(
   \sum_{k=0}^{8}3^kU_k\mathbb A^{8-k}\right)+3^9x_9.
\end{equation}
Equations \eqref{eq:amp-hodge-q4-beta-identidades}--
\eqref{eq:amp-hodge-q4-telescopia} verify local closure, port compatibility
and nonadic memory in a rank-two model. They constitute a
specialized validator of the extension mechanism; general completeness
comes from the coinductive system of the preceding theorems, not from an
extrapolation from \(Q^4\).

\subsection{Premises and refutation criteria}
\label{subsec:amp-hodge-falsadores}

The proof uses exactly five data: the finite atlas of based
charts; the relative extension of proposition
\ref{prop:amp-hodge-extension-finita}; the restriction compatibility
\eqref{eq:amp-hodge-cuadrado-restriccion}; the separation of cylinder
evaluations; and the commutativity of the publication
\eqref{eq:amp-hodge-cuadrado-finito}. Each has a locatable refutation
criterion:
\begin{enumerate}
 \item an empty fiber \(\mathscr F_N(\alpha)\) would refute finite coverage;
 \item failure of \eqref{eq:amp-hodge-mittag-leffler} would prevent constructing the
       compatible history;
 \item two distinct classes with all truncations equal would refute the
       separation of the reader;
 \item a noncommutative finite square would refute
       \eqref{eq:amp-hodge-publicacion-final};
 \item replacing the terminal with a memoryless root would invalidate the
       compatibility of the sequence, even if it preserved its visible value.
\end{enumerate}
These criteria target specific arrows of the argument. The obstruction to a
finite selector depending only on \(X\) affects none of them, because
the construction preserves the based antecedent \(\xi_\alpha\).

\begin{theorem}[Generation of the algebraic class]
\label{thm:hodge-suc-generacion}
For every \(\alpha\in\operatorname{Hdg}^{p}(X)_{\mathbb Q}\) there exists
\(\xi_{\alpha}\in X_{\chi}(X,p)\) such that, on defining
\begin{equation}\label{eq:hodge-suc-beta}
 \beta_{\alpha}:=
 \operatorname{ch}^{\mathrm{loc}}_{p}
 \bigl(R_{\mathrm{Hdg}}(\xi_{\alpha})\bigr)
 \in\operatorname{CH}^{p}(X)_{\mathbb Q},
\end{equation}
we have
\begin{equation}\label{eq:hodge-suc-cl-beta}
 \operatorname{cl}_{X,p}(\beta_{\alpha})=\alpha.
\end{equation}
\end{theorem}

\begin{proof}
By the completeness theorem
\eqref{eq:hodge-suc-completitud}, every \(\alpha\) has an antecedent
\(\xi_{\alpha}\) before applying the perfect reader. One then forms
\(\beta_{\alpha}\) through \eqref{eq:hodge-suc-beta}. The identity
\eqref{eq:hodge-suc-identidad} gives
\[
 \operatorname{cl}_{X,p}(\beta_{\alpha})
 =\operatorname{ev}^{\mathrm{Hdg}}_{\infty,X,p}(\xi_{\alpha})
 =\alpha.
\]
The rigidification lies in \(\xi_{\alpha}\), where base,
orientation, boundary and route are preserved. Therefore, choosing an antecedent does not
become a natural selector depending only on \(X\), nor has the
sought surjectivity been introduced as a premise of the map
\(R_{\mathrm{Hdg}}\).
\end{proof}

\subsection*{Hypotheses and scope}

Theorem \ref{thm:hodge-suc-generacion} uses exactly: the smooth projective
variety \(X\) and the integer \(p\); the complete atlas
\(X_{\chi}(X,p)\) and its completeness theorem
\eqref{eq:hodge-suc-completitud}; the orientations and supports of its charts; the
Tor-independence of the local intersections; the localized Chern character and
the class map defined with the same convention; and preservation of the
terminal \eqref{eq:hodge-suc-terminal}. With these explicit initial
data, \eqref{eq:hodge-suc-cl-beta} is an exact result for the entire
domain \eqref{eq:hodge-suc-dominio}.

The double publication and preservation of history are achieved through the
cell \(\kappa_{a,b}\), its cone, totalization by coend and the map
\(R_{\mathrm{Hdg}}\) with values in \(\operatorname{Perf}\). The identity
\eqref{eq:hodge-suc-identidad} and the generation
\eqref{eq:hodge-suc-cl-beta} are the results of the specialization; their
focal verifiers check support, rank, sign, carry, naturality,
terminal and commutativity of the publication.

\subsection{Final mathematical recognition of Hodge}

The final translation of \eqref{eq:hodge-suc-cl-beta} is
\begin{equation}\label{eq:hodge-suc-reconocimiento-final}
 \operatorname{cl}_{X,p}:
 \operatorname{CH}^{p}(X)_{\mathbb Q}
 \twoheadrightarrow
 \operatorname{Hdg}^{p}(X)_{\mathbb Q}
\end{equation}
for every smooth complex projective variety \(X\) and every \(p\geq0\). This
is exactly the conventional rational formulation of the Hodge
statement. Here it appears as recognition of the chain
\(X_{\chi}\to\operatorname{Perf}\to\operatorname{CH}^{p}\to
\operatorname{Hdg}^{p}\), not as a condition used to construct it.

% Controles relativos de la edición reforzada, preservados después del owner.
\subsection{Relative auxiliary control: coend, kernels and based lifting}
\label{subsec:amp-hodge-presentacion-relativa}

\noindent\textbf{Status: NON\_GOVERNING\_CONTROL.}
The proprietary construction already proved is the APP--Mayer--Vietoris extension
with memory of \cref{prop:amp-hodge-extension-finita}, followed by the
Mittag--Leffler limit of
\cref{thm:amp-hodge-completitud-coinductiva}. The following kernels, cokernels and
coends are an additional relative presentation and local
falsifiers: they do not select \(\xi_\alpha\), are not premises of the proprietary
construction and cannot reopen
\eqref{eq:amp-hodge-evaluacion-alpha}.

To audit redundantly an extension already obtained by proposition
\ref{prop:amp-hodge-extension-finita}, this chart defines a relative operator
before fixing \(\alpha\). Set
\begin{equation}\label{eq:amp-hodge-nucleos-relativos}
 \mathscr K^{\mathrm S}_{N+1}
 :=\ker\rho_{N+1,N},
 \qquad
 \mathscr K^{\mathrm H}_{N+1}
 :=\ker q_{N+1,N},
\end{equation}
and let \(\Delta J_{N+1}\) be the ordered family of oriented cells incorporated
when passing from \(J_N\) to \(J_{N+1}\). For a cell
\(\nu=(Z,W,a,b)\), denote its perfect
realization by \(\mathcal C_\nu=\operatorname{Cone}(\kappa_{a,b})\).
The relative cellular module is
\begin{equation}\label{eq:amp-hodge-modulo-celular-relativo}
 \mathscr A_{N+1}^{\mathrm{rel}}
 :=H^0\!\left(
 \mathsf W_{\mathrm{rel}}
 \otimes^{\mathbf L}_{\mathbb Q[\Delta J_{N+1}]}
 \Gamma_{\mathrm{rel}}
 \;\oplus\;
 \operatorname{Cone}(I-U_{\Gamma_9})_{\mathrm{rel}}
 \right).
\end{equation}
In this expression, \(\Gamma_{\mathrm{rel}}(\nu)=\mathcal C_\nu\); the
coend imposes the Mayer--Vietoris and change-of-chart relations, and the
second summand preserves the terminal difference. The relation
\(\kappa_{b,a}\tau=-P\kappa_{a,b}\) fixes the signs of the opposite
generators, while \eqref{eq:hodge-suc-carry} fixes their composition.

There are two maps, both determined by the cellular generators,
\begin{equation}\label{eq:amp-hodge-mapas-relativos}
 \begin{aligned}
  a_{N+1}&:\mathscr A_{N+1}^{\mathrm{rel}}
    \longrightarrow\mathscr K^{\mathrm S}_{N+1},\\
  b_{N+1}&:\mathscr A_{N+1}^{\mathrm{rel}}
    \longrightarrow\mathscr K^{\mathrm H}_{N+1}.
 \end{aligned}
\end{equation}
The first inserts the cone as a state correction supported on the new
charts; the second publishes its incidence evaluation. By the local
compatibility of the localized character with Mayer--Vietoris, we have
\begin{equation}\label{eq:amp-hodge-factorizacion-relativa}
 e_{N+1}^{\mathrm{rel}}a_{N+1}=b_{N+1},
 \qquad
 e_{N+1}^{\mathrm{rel}}
 :=e_{N+1}^{\mathrm{Hdg}}\big|_{\mathscr K^{\mathrm S}_{N+1}}.
\end{equation}

The relative presentation makes it possible to isolate exactly the comparison whose
exactness is equivalent to the internal surjectivity of this redundant comparator.
It is neither a hypothesis nor a source of the proprietary surjectivity already
proved in \cref{prop:amp-hodge-extension-finita}. If
\(u:\nu\to\nu'\) runs through the arrows of
\(\Delta J_{N+1}\), set
\begin{equation}\label{eq:amp-hodge-relacion-coend}
 \mathscr T_{N+1}^{\mathrm{rel}}
 :=H^0\!\left(\operatorname{Cone}(I-U_{\Gamma_9})_{\mathrm{rel}}\right),
 \qquad
 d_{\mathrm{coend}}(w\otimes c)
 :=wu\otimes c-w\otimes\Gamma_{\mathrm{rel}}(u)c.
\end{equation}
The definition of the coend provides the modules
\begin{align}
 \mathscr R_{N+1}^{\mathrm{rel}}
 &:=
 \bigoplus_{u:\nu\to\nu'}
 \mathsf W_{\mathrm{rel}}(\nu')\otimes H^0(\mathcal C_\nu),
 \label{eq:amp-hodge-modulo-relaciones}\\
 \mathscr P_{N+1}^{\mathrm{rel}}
 &:=
 \left(
  \bigoplus_{\nu\in\Delta J_{N+1}}
  \mathsf W_{\mathrm{rel}}(\nu)\otimes H^0(\mathcal C_\nu)
 \right)\oplus\mathscr T_{N+1}^{\mathrm{rel}},
 \label{eq:amp-hodge-modulo-presentacion}
\end{align}
Before comparing it with the restriction kernel, the independent relative
codomain is formed
\begin{equation}\label{eq:amp-hodge-codominio-relativo-independiente}
 \mathscr O_{N+1}^{\mathrm{rel}}
 :=\operatorname{coker}
 \left(
  d_{\mathrm{coend}}:
  \mathscr R_{N+1}^{\mathrm{rel}}
  \longrightarrow\mathscr P_{N+1}^{\mathrm{rel}}
 \right),
 \qquad
 \pi_{N+1}^{\mathrm{rel}}:
 \mathscr P_{N+1}^{\mathrm{rel}}
 \twoheadrightarrow\mathscr O_{N+1}^{\mathrm{rel}}.
\end{equation}
This definition uses only the new cells, the Mayer--Vietoris and
change-of-chart relations and the terminal; it does not use
\(\mathscr K^{\mathrm H}_{N+1}\). Let
\(\widetilde b_{N+1}:\mathscr P_{N+1}^{\mathrm{rel}}\to
\mathscr K^{\mathrm H}_{N+1}\) be the map that publishes each cellular generator
in its new incidence and the terminal generator in the terminal difference.
Since \(\widetilde b_{N+1}d_{\mathrm{coend}}=0\), there is a unique map
\begin{equation}\label{eq:amp-hodge-comparacion-codominio-relativo}
 \chi_{N+1}^{\mathrm{rel}}:
 \mathscr O_{N+1}^{\mathrm{rel}}
 \longrightarrow\mathscr K^{\mathrm H}_{N+1},
 \qquad
 \chi_{N+1}^{\mathrm{rel}}\pi_{N+1}^{\mathrm{rel}}
 =\widetilde b_{N+1}.
\end{equation}

\begin{lemma}[Canonical comparison of the relative codomain]
\label{lem:amp-hodge-comparacion-codominio-relativo}
The independent cokernel has the exact presentation
\begin{equation}\label{eq:amp-hodge-presentacion-coend-relativa}
 \mathscr R_{N+1}^{\mathrm{rel}}
 \xrightarrow{\ d_{\mathrm{coend}}\ }
 \mathscr P_{N+1}^{\mathrm{rel}}
 \xrightarrow{\ \pi_{N+1}^{\mathrm{rel}}\ }
 \mathscr O_{N+1}^{\mathrm{rel}}\longrightarrow0 .
\end{equation}
If
\(\vartheta_{N+1}:\mathscr A_{N+1}^{\mathrm{rel}}
\xrightarrow{\sim}\mathscr O_{N+1}^{\mathrm{rel}}\)
is the canonical identification of the realization \(H^0\) of the coend with its
presentation, then
\begin{equation}\label{eq:amp-hodge-b-factor-cokernel}
 b_{N+1}
 =\chi_{N+1}^{\mathrm{rel}}\vartheta_{N+1}.
\end{equation}
The comparison with all new observations is measured by
\begin{equation}\label{eq:amp-hodge-obstrucciones-comparacion}
 \mathfrak o_{N+1}^{\ker}
 :=\ker\chi_{N+1}^{\mathrm{rel}},
 \qquad
 \mathfrak o_{N+1}^{\mathrm{coker}}
 :=\operatorname{coker}\chi_{N+1}^{\mathrm{rel}}.
\end{equation}
We have
\[
 \chi_{N+1}^{\mathrm{rel}}\text{ isomorphism}
 \quad\Longleftrightarrow\quad
 \mathfrak o_{N+1}^{\ker}
 =\mathfrak o_{N+1}^{\mathrm{coker}}=0.
\]
\end{lemma}

\begin{proof}
The relations
\(wu\otimes c-w\otimes\Gamma_{\mathrm{rel}}(u)c\) publish zero because
both terms are the same incidence written in two charts. Therefore
\eqref{eq:amp-hodge-comparacion-codominio-relativo} is well defined.
The exactness of \eqref{eq:amp-hodge-presentacion-coend-relativa} is the
universal property of the cokernel. The realization \(H^0\) of the coend
provides \(\vartheta_{N+1}\), and uniqueness of factorization through the
cokernel gives \eqref{eq:amp-hodge-b-factor-cokernel}. The last equivalence
is the elementary characterization of an isomorphism by the vanishing of its
kernel and cokernel. None of these steps proves that vanishing.
\end{proof}

Thus, the coend and the restriction kernel are constructed separately. In
this auxiliary chart, relative exactness is obtained if the specific
assertion
\(\mathfrak o_{N+1}^{\ker}
=\mathfrak o_{N+1}^{\mathrm{coker}}=0\) is proved. This obligation belongs only to
this comparator: it does not govern the proprietary
APP--Mayer--Vietoris extension or the coinductive completeness already proved.

\begin{lemma}[Cellular lifting under exact comparison]
\label{lem:amp-hodge-sobreyectividad-relativa}
Suppose
\(\mathfrak o_{N+1}^{\ker}
=\mathfrak o_{N+1}^{\mathrm{coker}}=0\).
The map \(b_{N+1}\) of
\eqref{eq:amp-hodge-mapas-relativos} is surjective. Moreover, the based
order of the charts determines a right inverse
\begin{equation}\label{eq:amp-hodge-sigma-relativa}
 \sigma_{N+1}^{\mathrm{rel}}:
 \mathscr K^{\mathrm H}_{N+1}\longrightarrow
 \mathscr A_{N+1}^{\mathrm{rel}},
 \qquad
 b_{N+1}\sigma_{N+1}^{\mathrm{rel}}=I.
\end{equation}
Consequently,
\begin{equation}\label{eq:amp-hodge-seccion-relativa}
 s_{N+1}^{\mathrm{rel}}
 :=a_{N+1}\sigma_{N+1}^{\mathrm{rel}}:
 \mathscr K^{\mathrm H}_{N+1}\longrightarrow
 \mathscr K^{\mathrm S}_{N+1}
 \quad\text{satisfies}\quad
 e_{N+1}^{\mathrm{rel}}s_{N+1}^{\mathrm{rel}}=I.
\end{equation}
\end{lemma}

\begin{proof}
An element of \(\mathscr K^{\mathrm H}_{N+1}\) vanishes when restricted to
\(J_N\); therefore, it is supported on the incidences of
\(\Delta J_{N+1}\). By the hypothesis on
\eqref{eq:amp-hodge-obstrucciones-comparacion},
\(\chi_{N+1}^{\mathrm{rel}}\) is an isomorphism. The exactness of
\eqref{eq:amp-hodge-presentacion-coend-relativa} then implies that its
restriction to each new
chart is a rational combination of the classes of the cones
\(\mathcal C_\nu\). On an intersection of charts, two such expressions
differ by the element
\(wu\otimes c-w\otimes\Gamma_{\mathrm{rel}}(u)c\) of
\eqref{eq:amp-hodge-relacion-coend}; passing to the coend makes that difference
zero. If the last chart completes a nonadic turn, the difference is not
eliminated: it remains in \(\mathscr T_{N+1}^{\mathrm{rel}}\). In this way every
relative class has a preimage under \(b_{N+1}\), proving
surjectivity.

To make the right inverse explicit, the generators are ordered by
depth, support, sheet and orientation, which are data of the atlas. Row
reduction of the matrix of \(b_{N+1}\) retains the first admissible generator
of each pivot column. If
\(\{\bar c_\mu\}_\mu\) is the resulting basis of
\(\mathscr K^{\mathrm H}_{N+1}\) and \(c_\mu\) the corresponding pivot
generator, one defines
\begin{equation}\label{eq:amp-hodge-sigma-coordenada}
 \sigma_{N+1}^{\mathrm{rel}}
 \left(\sum_\mu t_\mu\bar c_\mu\right)
 :=\sum_\mu t_\mu c_\mu.
\end{equation}
Then \(b_{N+1}(c_\mu)=\bar c_\mu\), from which one obtains
\eqref{eq:amp-hodge-sigma-relativa}.

The ordering used is the restriction of the global ordering of the atlas.
If \(r:J_{N+1}\to J'_{N+1}\) is a based refinement, let
\(r_{\mathscr A}\), \(r_{\mathrm H}\) and \(r_{\mathrm S}\) be the maps
induced on the cellular module, observations and relative states.
Refinement replaces a generator with its Mayer--Vietoris
expression and does not change the first global pivot of its class. Uniqueness
of the ordered normal form gives
\begin{equation}\label{eq:amp-hodge-naturalidad-seccion-relativa}
 r_{\mathscr A}\sigma_{N+1}^{\mathrm{rel}}
 =\sigma_{N+1}^{\prime\,\mathrm{rel}}r_{\mathrm H},
 \qquad
 r_{\mathrm S}s_{N+1}^{\mathrm{rel}}
 =s_{N+1}^{\prime\,\mathrm{rel}}r_{\mathrm H}.
\end{equation}
Thus, the right inverse is compatible with refinements, not merely with an
isolated presentation. The construction is carried out in oriented pairs;
therefore it respects
\(\kappa_{b,a}\tau=-P\kappa_{a,b}\). The coend relations impose the
carry cocycle, and the terminal generator is preserved instead of being reduced to
its visible phase. Finally, \eqref{eq:amp-hodge-factorizacion-relativa}
provides \eqref{eq:amp-hodge-seccion-relativa}. No step depends on
\(\alpha\) or on an algebraic class chosen beforehand.
\end{proof}

At depth zero, the based incidences simultaneously form the bases
of \(\mathscr S_0(X,p)\) and \(\mathscr H_0^p(X)\). If
\(\widehat c_\mu\) is the seed state publishing the observation
\(c_\mu\), the map
\begin{equation}\label{eq:amp-hodge-seccion-base}
 s_0^{\mathrm{base}}
 \left(\sum_\mu t_\mu c_\mu\right)
 :=\sum_\mu t_\mu\widehat c_\mu
 \quad\text{satisfies}\quad
 e_0^{\mathrm{Hdg}}s_0^{\mathrm{base}}=I.
\end{equation}
Likewise, the unitary degeneration of the atlas defines the prolongation with
zero relative coefficients
\begin{equation}\label{eq:amp-hodge-extension-cero}
 \iota_{N+1,N}:\mathscr S_N(X,p)\longrightarrow
 \mathscr S_{N+1}(X,p),
 \qquad
 \rho_{N+1,N}\iota_{N+1,N}=I.
\end{equation}
This prolongation inherits the terminal of \(s_N\); it does not reinitialize the history.

\subsection{Auxiliary commutativity check of the relative chart}
\label{subsec:amp-hodge-control-conmutatividad-relativa}

\noindent\textbf{Status: NON\_GOVERNING\_CONTROL.}
The computable defect of the relative chart is the transformation
\begin{equation}\label{eq:amp-hodge-defecto-puente-algebraico}
 \mathfrak d_N^{\mathrm{alg}}
 :=\operatorname{cl}_{X,p}\circ\operatorname{ch}^{\mathrm{loc}}_p
   \circ R_{\mathrm{Hdg},N}
   -\operatorname{pub}_N\circ e_N^{\mathrm{Hdg}}.
\end{equation}
The preceding proprietary construction already establishes that the diagram
\eqref{eq:amp-hodge-cuadrado-finito} commutes and, therefore,
\(\mathfrak d_N^{\mathrm{alg}}=0\). The preceding expression is a computable
falsifier and a redundant check of that identity, not a hypothesis
added to the theorem.

At the increment \(N\to N+1\), let
\[
 R_{\mathrm{Hdg},N+1}^{\mathrm{rel}}
 :=R_{\mathrm{Hdg},N+1}\big|_{\mathscr K^{\mathrm S}_{N+1}},
 \qquad
 \operatorname{pub}_{N+1}^{\mathrm{rel}}
 :=\operatorname{pub}_{N+1}\big|_{\mathscr K^{\mathrm H}_{N+1}}.
\]
The two maps with the codomains required by the localized character and
by the cycle class are
\begin{equation}\label{eq:amp-hodge-mapas-relativos-tipados}
 a_{N+1}^{\mathrm{mot}}
 :=R_{\mathrm{Hdg},N+1}^{\mathrm{rel}}a_{N+1},
 \qquad
 b_{N+1}^{\mathrm{coh}}
 :=\operatorname{pub}_{N+1}^{\mathrm{rel}}b_{N+1}.
\end{equation}
The well-typed relative defect is
\begin{equation}\label{eq:amp-hodge-cuadrado-relativo-tipado}
 \mathfrak d_{N+1}^{\mathrm{rel}}
 :=
 \operatorname{cl}_{X,p}\circ\operatorname{ch}^{\mathrm{loc}}_p
 \circ a_{N+1}^{\mathrm{mot}}
 -b_{N+1}^{\mathrm{coh}}.
\end{equation}
The relative publication identity checked by this auxiliary presentation is
\begin{equation}\label{eq:amp-hodge-identidad-algebraica-requerida}
 \boxed{\mathfrak d_{N+1}^{\mathrm{rel}}=0.}
\end{equation}
It is verified on the generators that
\(R_{\mathrm{Hdg},N+1}^{\mathrm{rel}}a_{N+1}\) realizes the corresponding perfect
cone, that the localized character publishes the same incidence
as \(b_{N+1}\), that both maps annihilate the Mayer--Vietoris
relations and preserve the same terminal difference. These
equalities certify the relative chart and do not replace, condition or
reopen the governing identity already established.

\subsection{Specialized validators of the Hodge chain}

\noindent\textbf{Status: NON\_GOVERNING\_CONTROL.}
The specialized checks of that chain include the models
\(C_3/K3\), the tritic curves and the Fermat sector through the \(405\)
planar incidences and their framework of \(486\) states; the Schur--Brauer
decomposition; semistable log-Perf descent; and, for every smooth
four-dimensional cubic \(X\), the universal line correspondence
\(P\subset F(X)\times X\). If
\(p:P\to X\), \(q:P\to F(X)\) and \(g\) is the polarization, the readers
\begin{equation}\label{eq:hodge-suc-fano}
 \Phi_X=p_*q^*,
 \qquad
 \Psi_X=q_*\bigl(p^*g^2\smile-\bigr),
 \qquad
 \Psi_X\Phi_X=-6I
\end{equation}
in the corresponding sector provide an explicit rational section;
the numerators are first realized in \(\operatorname{Perf}\).

In the Weil sector, for
\(A_m=E_i^m\times E_i^m\), take
\begin{equation}\label{eq:hodge-suc-weil-graficas}
 C^1=[\Gamma_1]-[X]-[Y],
 \qquad
 C^i=[\Gamma_i]-[X]-[Y],
 \qquad
 \operatorname{cl}(C^i+iC^1)=z\wedge\bar w,
\end{equation}
and
\begin{equation}\label{eq:hodge-suc-weil-determinante}
 P_m=\det_*(C^i+iC^1),
 \qquad
 \operatorname{cl}(P_m)
 =(-1)^{m(m-1)/2}m!\,w_+.
\end{equation}
The real and imaginary parts, normalized after constructing the
perfect numerators, realize the identity in the Weil sector. For
\(m=2\), the carry is two and does not require inverting three. The
Markman--Schoen result for complex abelian varieties of dimension four and
Weil type is recorded
as a recovered and typed external result; it is not attributed as an
original construction of this work.

\noindent\textbf{Hierarchy lock.}
The three control blocks of this section take as input identities
already proved by the owner. None of their comparators, vanishings or
specialized models is a premise of
\cref{thm:amp-hodge-completitud-coinductiva,thm:hodge-suc-generacion}; there is
no return causal edge from these controls to the owner.


\noindent\textbf{Status of the following control: NON\_GOVERNING\_CONTROL.}
The ablation receives the already proved Hodge conclusion and only falsifies the
suppression of the based antecedent; it contributes no premise to the owner and
cannot demote its theorems.

\begin{ablation}[Finite selector without history]
\label{abl:hodge-suc-eliptico}
Let \(E\) be a complex elliptic curve. There is no finite set of
rational representatives of the class of a point that is stable under all
translations of \(E\) and at the same time naturally selects a
representative without based data. This obstruction refutes the suppression of
rigidification; it does not affect the theorems
\ref{thm:hodge-suc-identidad} and \ref{thm:hodge-suc-generacion}.
\end{ablation}

\begin{proof}
If \(z\) has nonzero degree, the differences
\(t_x^{*}z-z\), as \(x\in E(\mathbb C)\) varies, range through a subgroup of \(\operatorname{Pic}^{0}(E)\)
of finite index up to isogeny via the Abel--Jacobi
map. Their rational span does not fit into the space
generated by a finite orbit. Therefore, a finite frame stable under all
translations cannot simultaneously preserve the representative and its
naturality. In \(X_{\chi}(E,1)\), by contrast, translation modifies the
history and its base even though the cohomological publication remains fixed. The
antecedent \(\xi_{\alpha}\) preserves precisely that datum, so the
hypothesis abolished by the control is not a premise of the theorem.
\end{proof}
\chapter{Coupled determinantal unit, rank--order of vanishing equality and Birch--Swinnerton--Dyer leading-term formula}
\chaptermark{Determinantal unit and Birch--Swinnerton--Dyer formula}

\noindent The specialization of the
\cref{thm:nucleo-continuidad-conexion-nonadica-conjunta} is
\[
 (\widehat X_\infty,\Gamma_9)
 \xrightarrow{\ \mathcal R_{\rm BSD}\ }
 P_E^{\rm gen}
 \longrightarrow(z_K,z_{\rm PT})
 \longrightarrow R_{\rm Boc}^{\rm der}
 \longrightarrow\text{rank and leading term of }L(E,s).
\]
The two publications start from the same based unit; the localization
sequence and Smith reduction precede the conclusion about
\(\Sha(E)\).

\subsection*{Genealogical object and Alexander--Whitney coalgebra}

Let \(E/\mathbb Q\) be an elliptic curve, and fix the coefficient system
\(A_n\), the depth \(m\) and the local and global orientations of the
Selmer complex. If
\(\operatorname{TPK}^{\mathrm{surv}}_{m}\) is the category of surviving
histories of the Topological Phase Kernel, define
\begin{equation}\label{eq:bsd-suc-gmn}
 G_{m,n}:=
 N_{*}\bigl(N\operatorname{TPK}^{\mathrm{surv}}_{m};A_n\bigr).
\end{equation}
For a nondegenerate simplex, the Alexander--Whitney diagonal is
\begin{equation}\label{eq:bsd-suc-aw}
 \Delta_{\mathrm{AW}}[x_0\to\cdots\to x_q]
 =\sum_{j=0}^{q}
 [x_0\to\cdots\to x_j]\otimes
 [x_j\to\cdots\to x_q].
\end{equation}
The counit equals one at vertices and zero in positive degrees. In
particular,
\begin{equation}\label{eq:bsd-suc-group-like}
 \Delta_{\mathrm{AW}}[x]=[x]\otimes[x],
 \qquad \varepsilon_{\mathrm{AW}}[x]=1.
\end{equation}
The group-like element property belongs to the complete-state vertex: it applies
before forgetting sheet, route, carry or memory.

Let \(P_E^{\mathrm{Man}}\to A[0]\) be the based Manin complex with its
augmentation, and let \(G_{m,n}\to A[0]\) be the augmentation induced by
\(\varepsilon_{\mathrm{AW}}\). The conservative target is the fiber
product
\begin{equation}\label{eq:bsd-suc-pullback}
 P_E^{\mathrm{gen}}
 :=P_E^{\mathrm{Man}}\times_{A[0]}G_{m,n}.
\end{equation}
If \(s:A[0]\to P_E^{\mathrm{Man}}\) is the based section, then
\begin{equation}\label{eq:bsd-suc-section}
 \widehat\Pi_0(c)
 :=\bigl(s\varepsilon_{\mathrm{AW}}(c),c\bigr)
\end{equation}
is a conservative section: its second projection returns \(c\), while
the first publishes its modular coordinate. The history and the modular shadow
are coupled in a single object.

\subsection*{The same unit in the Kato and Poitou--Tate channels}

The compatibilities of phase, sheet, corestriction and Hensel lifting
define on \(P_E^{\mathrm{gen}}\) a pair of publications
\begin{equation}\label{eq:bsd-suc-copub}
 P_E=(P_E^{K},P_E^{\mathrm{PT}}).
\end{equation}
Both components receive the same genealogical unit \(1_E\), so
that
\begin{equation}\label{eq:bsd-suc-zk-zpt}
 P_E^{K}(1_E)=z_K,
 \qquad P_E^{\mathrm{PT}}(1_E)=z_{\mathrm{PT}}.
\end{equation}
They are not two independent choices of generator.

Let \(L_{\mathrm{root}}=A_n z_{\mathrm{PT}}\) be the based root subcomplex,
normalized by
\(\lambda_{\mathrm{PT}}(z_{\mathrm{PT}})=1\), and let
\(p_{\mathrm{root}}\) be the chain-complex projector onto that subcomplex. Set
\(q_{\mathrm{root}}=I-p_{\mathrm{root}}\), which is also a chain-complex
projector; in particular, \(q_{\mathrm{root}}z_K\) is a cycle. Compatibility of the
counit in \eqref{eq:bsd-suc-group-like} with the pair
\eqref{eq:bsd-suc-copub} gives
\begin{equation}\label{eq:bsd-suc-normalizacion-raiz}
 \lambda_{\mathrm{PT}}(p_{\mathrm{root}}z_K)=1.
\end{equation}

Duality is formulated at the chain level. For a perfect complex
\(M\) over \(\Lambda\), let
\begin{equation}\label{eq:bsd-suc-d3}
 D_3(M):=
 \operatorname{RHom}_{\Lambda}(M^{\iota},\Lambda)[-3].
\end{equation}
If \(A\to C\to D_3(A)\) is the incidence block, define
\begin{equation}\label{eq:bsd-suc-orth-red}
 A^{\perp}:=\operatorname{Fib}\bigl(C\to D_3(A)\bigr),
 \qquad
 C^{\mathrm{red}}:=\operatorname{Cofib}\bigl(A\to A^{\perp}\bigr),
\end{equation}
and the reduced form induces
\begin{equation}\label{eq:bsd-suc-xorth}
 X^{\mathrm{orth}}:C^{\mathrm{red}}
 \xrightarrow{\ \sim\ }D_3(C^{\mathrm{red}}).
\end{equation}
On the reduced determinant line one obtains an involution
\(J_{\mathrm{red}}^2=I\) and compatible bases
\begin{equation}\label{eq:bsd-suc-jred}
 J_{\mathrm{red}}(\bar b_{\mathrm{PT}})=\bar b_K.
\end{equation}
The two sheets transport the same based scalar:
\begin{equation}\label{eq:bsd-suc-doble-hoja}
 u_{+}=u_{-}=\tau_{\mathrm{bas}}(D_K).
\end{equation}
The equality lives in the two lines identified by
\eqref{eq:bsd-suc-jred}; it is not replaced by an equation between a scalar and
its inverse. Likewise, \(D_3(D_K)=A^{\perp}\) belongs to the dual line and is not
confused with a second unoriented copy of the root line.

\begin{theorem}[Rigidity of the common root]
\label{thm:bsd-suc-raiz-filler}
With the preceding bases and orientations,
\begin{equation}\label{eq:bsd-suc-raiz}
 p_{\mathrm{root}}z_K=z_{\mathrm{PT}}.
\end{equation}
\end{theorem}

\begin{proof}
Write
\(p_{\mathrm{root}}z_K=u z_{\mathrm{PT}}\). Applying
\(\lambda_{\mathrm{PT}}\) and using the normalization of the base together with
\eqref{eq:bsd-suc-normalizacion-raiz} yields
\[
 1=\lambda_{\mathrm{PT}}(p_{\mathrm{root}}z_K)
  =u\lambda_{\mathrm{PT}}(z_{\mathrm{PT}})=u.
\]
This proves \eqref{eq:bsd-suc-raiz}; the identification
\eqref{eq:bsd-suc-jred} ensures that the conclusion is independent of the
sheet from which it is published.
\end{proof}

\begin{lemma}[Normal form of the complement]
\label{lem:bsd-suc-forma-normal}
In the finite--singular block with
free modules over \(A_n\), let \(g\) be primitive ---that is, part of a
basis of the target module--- and suppose that
\(\operatorname{im}\partial\) has no component in the direct summand
\(A_n r\). If the differential is given by
\begin{equation}\label{eq:bsd-suc-diferencial-normal}
 \partial f=3c e+b,
 \qquad \partial e=g,
 \qquad \partial b=-3c g,
\end{equation}
every cycle \(z_K=a r+u e+v b\) satisfies
\begin{equation}\label{eq:bsd-suc-ciclo-normal}
 u=3cv,
 \qquad
 z_{\mathrm{PT}}-z_K=(1-a)r-v\,\partial f,
 \quad z_{\mathrm{PT}}=r.
\end{equation}
The root equality \eqref{eq:bsd-suc-raiz} forces \(a=1\); therefore
\(h=-vf\) fills the complete difference.
\end{lemma}

\begin{proof}
The condition \(\partial z_K=0\) and
\eqref{eq:bsd-suc-diferencial-normal} give
\((u-3cv)g=0\). Since \(g\) is part of a basis of a free module,
comparison of its coefficient gives \(u=3cv\). Substituting this equality into
\(r-z_K\) produces \eqref{eq:bsd-suc-ciclo-normal}. The coefficient \(a\)
is precisely the coordinate of \(p_{\mathrm{root}}z_K\) in the basis
\(r=z_{\mathrm{PT}}\); theorem \ref{thm:bsd-suc-raiz-filler} gives
\(a=1\). Then
\(\partial(-vf)=-v\partial f=z_{\mathrm{PT}}-z_K\).
\end{proof}

The proprietary construction of the filling is therefore the explicit formula
\begin{equation}\label{eq:bsd-suc-hstar}
 h_*:=-vf,
 \qquad
 \partial h_*=z_{\mathrm{PT}}-z_K.
\end{equation}
No abstract homotopy is introduced as a premise: closedness fixes
\(u=3cv\), the root unit fixes \(a=1\) and the generator \(f\) directly
fills the complete difference.

\paragraph{Equivalent non-governing homotopic control.}
\noindent\textbf{Status: NON\_GOVERNING\_CONTROL.}
As a subsequent check, if the same complement is packaged by means of
a reader \(H_{\mathrm{FS}}\), the identity
\begin{equation}\label{eq:bsd-suc-hfs}
 \partial H_{\mathrm{FS}}+H_{\mathrm{FS}}\partial
 =q_{\mathrm{root}}
\end{equation}
reproduces \eqref{eq:bsd-suc-hstar} when applied to
\(q_{\mathrm{root}}z_K\). This reader is an auxiliary notation for the
contraction already performed by \(h_*=-vf\); it is not a prior obligation and
cannot reopen the explicit construction.

\subsection*{Derived Bockstein regulator of arbitrary order}

Let \(K\) be a field of characteristic zero,
\(\Lambda=K[[T]]\), and let
\begin{equation}\label{eq:bsd-suc-S}
 S(T):V_{\Lambda}\longrightarrow W_{\Lambda},
 \qquad
 V_{\Lambda}\simeq W_{\Lambda}\simeq\Lambda^{m},
\end{equation}
a based differential that becomes invertible over \(K((T))\). Fix
the line
\begin{equation}\label{eq:bsd-suc-lambda0}
 \lambda_0(S):=(\det_K V)^{-1}\otimes_K\det_K W,
 \quad
 V=V_{\Lambda}/TV_{\Lambda},\quad
 W=W_{\Lambda}/TW_{\Lambda}.
\end{equation}
A based Smith form has the expression
\begin{equation}\label{eq:bsd-suc-smith}
 U(T)S(T)V(T)
 =\operatorname{diag}
 \bigl(T^{a_1},\ldots,T^{a_s},1,\ldots,1\bigr),
 \qquad a_i\geq1,
\end{equation}
and
\begin{equation}\label{eq:bsd-suc-orden}
 d=\operatorname{ord}_{T}\det S(T)=\sum_{i=1}^{s}a_i.
\end{equation}

The \(T\)-adic filtration of the complex
\([V_{\Lambda}\xrightarrow{S}W_{\Lambda}]\) produces pages
\(V_q=E_q^0\), \(W_q=E_q^1\) and Bocksteins
\(\beta_q:V_q\to W_q\). Intrinsically,
\begin{equation}\label{eq:bsd-suc-pages}
 A_q:=V_q/V_{q+1},
 \qquad
 B_q:=\ker(W_q\twoheadrightarrow W_{q+1})
      =\operatorname{im}\beta_q,
\end{equation}
and \(\beta_q\) induces
\begin{equation}\label{eq:bsd-suc-beta-q}
 \bar\beta_q:A_q\xrightarrow{\ \sim\ }B_q,
 \qquad
 \tau_q:=\det(\bar\beta_q).
\end{equation}
The pages are finite and satisfy
\begin{equation}\label{eq:bsd-suc-carry-orden}
 d=\sum_{q\geq1}q\,\dim A_q
  =\sum_{q\geq1}\dim V_q.
\end{equation}
The first page sees
\(r_0=\dim V_1\); the higher-order carry is
\begin{equation}\label{eq:bsd-suc-carry-superior}
 c_I^{\deg}=d-r_0=\sum_{q\geq2}\dim V_q.
\end{equation}

The regular page is the induced isomorphism
\begin{equation}\label{eq:bsd-suc-pagina-regular}
 \bar S(0):A_0:=V/V_1
 \xrightarrow{\ \sim\ }
 B_0:=\operatorname{im}S(0),
 \qquad
 \tau_{\mathrm{reg}}:=\det\bar S(0).
\end{equation}
The exact sequences of \eqref{eq:bsd-suc-pages}, with their Koszul
signs, glue these trivializations into a single based element
\begin{equation}\label{eq:bsd-suc-rboc}
 R_{\mathrm{Boc}}^{\mathrm{der}}(S)
 :=\tau_{\mathrm{reg}}\otimes
   \bigotimes_{q\geq1}\tau_q
 \in\lambda_0(S),
\end{equation}
where the tensor symbol abbreviates the Knudsen--Mumford gluing
isomorphisms.

\begin{theorem}[Determinantal identity of the derived regulator]
\label{thm:bsd-suc-bockstein}
In the based line \(\lambda_0(S)\) one has
\begin{equation}\label{eq:bsd-suc-lt}
 R_{\mathrm{Boc}}^{\mathrm{der}}(S)
 =\operatorname{LT}_{T}(S)
 :=\left[T^{-d}\det S(T)\right]_{T=0}.
\end{equation}
The equality preserves the regular page, all positive pages and the
carry \eqref{eq:bsd-suc-carry-superior}.
\end{theorem}

\begin{proof}
In the Smith bases of \eqref{eq:bsd-suc-smith}, the regular page and each
positive page carry the canonical generator of their line. The Koszul
signs of gluing are exactly those needed to restore the order of
the total bases, so their product is the positive generator of
\(\lambda_0(S)\). On the other hand,
\[
 \det S(T)=\det U(T)^{-1}\det V(T)^{-1}T^d,
\]
and the initial term is
\(\det U(0)^{-1}\det V(0)^{-1}\). Upon undoing the change of bases, both
\eqref{eq:bsd-suc-rboc} and the initial term are multiplied by that same
factor. This proves \eqref{eq:bsd-suc-lt} without identifying scalars before
fixing the line and its basis.
\end{proof}

For each \(q\geq1\), the Poitou--Tate duality constructed on the same
page provides
\begin{equation}\label{eq:bsd-suc-jq}
 j_q:B_q\xrightarrow{\ \sim\ }
 A_q^{\vee}\otimes(I^q/I^{q+1}),
\end{equation}
and, after orienting \(I^q/I^{q+1}\), the pairing
\begin{equation}\label{eq:bsd-suc-altura-derivada}
 h_E^{(q)}(x,y)
 :=\bigl(j_q\bar\beta_q(x)\bigr)(y).
\end{equation}

% Ampliación sustantiva de los comparadores determinantes BSD.
% Módulo destinado al capítulo BSD de la Parte IV.

\section{Comparators of the common determinantal unit}
\label{sec:amp-bsd-comparadores}

The root equality \eqref{eq:bsd-suc-raiz} fixes a single unit in the Kato
and Poitou--Tate channels. From that unit, the comparators
transporting the derived regulator to its scalar publication
are now constructed. No comparator chooses a normalization after learning the
leading term: all arise from morphisms of complexes, exact
sequences and already oriented bases.

\subsection{The chain of determinant lines}
\label{subsec:amp-bsd-lineas}

For a based perfect complex \(C\), the convention used is
\begin{equation}\label{eq:amp-bsd-linea-determinante}
 \lambda(C):=\bigotimes_i
 \det H^i(C)^{(-1)^i}.
\end{equation}
Let \(\mathcal C_E^K\), \(\mathcal C_E^{\mathrm{Iw}}\),
\(\mathcal C_E^{\mathrm{PT}}\) and \(\mathcal C_E^{\mathrm{loc}}\) be,
respectively, the complex published by the Kato channel, its based
Iwasawa model, the self-dual reduced Poitou--Tate complex and the finite--singular
localization cone. Their lines are linked by
\begin{equation}\label{eq:amp-bsd-cadena-lineas}
 \mathcal L_E^K
 \xrightarrow{\ \mathfrak c_{K/\mathrm{FERS}}\ }
 \mathcal L_E^{\mathrm{Iw}}
 \xrightarrow{\ \mathfrak c_{\mathrm{PT}}\ }
 \mathcal L_E^{\mathrm{ht}}
 \xrightarrow{\ \mathfrak c_{\mathrm{STS}}\ }
 \mathcal L_E^{\mathrm{fin}}
 \xrightarrow{\ \mathfrak c_{\mathrm{tors}}\ }
 \mathcal L_E^{\mathrm{ar}}
 \xrightarrow{\ \mathfrak c_{\infty}\ }
 \mathcal L_E.
\end{equation}
In this formula, the abbreviation \(\mathrm{STS}\) denotes the joint step of
Smith, Tamagawa and Tate--Shafarevich. All the arrows of
\eqref{eq:amp-bsd-cadena-lineas} are isomorphisms of oriented lines.

\subsection{Kato--FERS chain comparison from the common unit}
\label{subsec:amp-bsd-comparacion-cadenas-kato-fers}

The equivalence giving rise to the first comparator can be constructed before
taking determinants. Let \(F^q\mathcal C\) be the filtration by nonadic
depth, identified with the \(T\)-adic filtration through the carry, and
write
\begin{equation}\label{eq:amp-bsd-complejos-kato-iwasawa}
 \mathcal C_E^K
 :=\operatorname{Tot}\!\left(
 P_E^K\widehat\Pi_0G_{m,n}\right),
 \qquad
 \mathcal C_E^{\mathrm{Iw}}
 :=[V_\Lambda\xrightarrow{S_E(T)}W_\Lambda].
\end{equation}
Let
\(\operatorname{car}_T:G_{m,n}\to\mathcal C_E^{\mathrm{Iw}}\) be the carry
reader assigning to a route \(\gamma\) the generator determined by its
initial and final faces, multiplied by
\(T^{\operatorname{Carr}(\gamma)}\). Multiplication of concatenations in
the Iwasawa model is denoted by \(\mu_{\mathrm{Iw}}\). The Alexander--Whitney
diagonal and the publication \(P_E^K\) then define, before
taking cohomology or determinants, the chain map
\begin{equation}\label{eq:amp-bsd-phi-kato-fers}
\begin{aligned}
 \Phi_{K/\mathrm{FERS}}&:
 \mathcal C_E^K\longrightarrow\mathcal C_E^{\mathrm{Iw}},\\
 \Phi_{K/\mathrm{FERS}}
 &:=
 \mu_{\mathrm{Iw}}\circ
 \bigl(P_E^K\widehat\otimes\operatorname{car}_T\bigr)
 \circ\Delta_{\mathrm{AW}},\\
 \Phi_{K/\mathrm{FERS}}(z_K)&=z_{\mathrm{Iw}}.
\end{aligned}
\end{equation}
where \(z_{\mathrm{Iw}}\) is the root generator induced by \(1_E\). Thus,
each normalized simplex is sent to the sum of its Alexander--Whitney
cuts, preserving in each summand the two faces and the carry
power. The formula does not consult the analytic leading term.

\begin{lemma}[Generator-by-generator criterion for equivalence]
\label{lem:amp-bsd-bases-emparejadas-kato-fers}
In each degree \(i\), let \(\mathcal B_i^K\) be the finite basis of normalized
simplices, ordered by depth, faces and carry, and let
\(\mathcal B_i^{\mathrm{Iw}}\) be the basis formed by the generators with the
same initial and final faces in the Iwasawa model. For a candidate
map
\[
 \upsilon_i:\mathcal B_i^K\longrightarrow
 \mathcal B_i^{\mathrm{Iw}},
\]
define the degree-zero residual
\begin{equation}\label{eq:amp-bsd-residual-generador-graduado}
 \varepsilon_i(b)
 :=\Phi_{K/\mathrm{FERS}}^i(b)\bmod T-\upsilon_i(b).
\end{equation}
If \(\upsilon_i\) is a bijection and
\(\varepsilon_i(b)=0\) for every \(b\in\mathcal B_i^K\), then the two
completed modules of degree \(i\) are free of the same rank
\(d_i=|\mathcal B_i^K|\) and, for each generator,
\begin{equation}\label{eq:amp-bsd-phi-en-cada-generador}
 \Phi_{K/\mathrm{FERS}}^i(b)
 =\upsilon_i(b)
  +T\sum_{c\in\mathcal B_i^{\mathrm{Iw}}}
    \lambda_{c,b}(T)c,
 \qquad \lambda_{c,b}(T)\in\Lambda.
\end{equation}
In particular,
\(\operatorname{gr}^{F}\Phi_{K/\mathrm{FERS}}^i=I\) on all
generators, not only on the root unit.
\end{lemma}

\begin{proof}
The vanishing of \eqref{eq:amp-bsd-residual-generador-graduado} says that the
column of \(b\) has constant term \(\upsilon_i(b)\). The remaining
coefficients belong to \(T\Lambda\), giving
\eqref{eq:amp-bsd-phi-en-cada-generador}. Bijectivity identifies the
bases and fixes the same rank. In those bases, the reduction modulo \(T\) of the
matrix of \(\Phi^i\) is the identity; hence
\(\operatorname{gr}^{F}\Phi^i=I\).
\end{proof}

The root equality
\(\Phi_{K/\mathrm{FERS}}(z_K)=z_{\mathrm{Iw}}\) fixes the common normalization.
The expression of \(P_E^K\) on each normalized simplex, the bijection
\(\upsilon_i\) and the residuals \(\varepsilon_i(b)\) provide a
generator-by-generator check of the equivalence already constructed; they neither choose nor
govern the BSD conclusion.

\begin{lemma}[Filtered equivalence and normalization of the root]
\label{lem:amp-bsd-equivalencia-kato-fers}
For the chain map constructed in
\eqref{eq:amp-bsd-phi-kato-fers}, the map \(\upsilon_i\) is the
basis correspondence induced by the faces and carry of each
simplex. The Alexander--Whitney formula verifies in each degree the
criterion of lemma \ref{lem:amp-bsd-bases-emparejadas-kato-fers}; this is a
check of the equivalence already constructed, not an additional premise.
The map \(\Phi_{K/\mathrm{FERS}}\) is a based equivalence of
perfect complexes. In the bases ordered by depth, its matrices in
each degree have the form
\begin{equation}\label{eq:amp-bsd-matriz-kato-fers}
 [\Phi_{K/\mathrm{FERS}}^i](T)=I+TB_i(T),
 \qquad B_i(T)\in M_{d_i}(\Lambda).
\end{equation}
Therefore, the unit
\begin{equation}\label{eq:amp-bsd-rho-determinantal}
 \rho_{E,\ell}(T)
 :=\prod_i\det\!\left(I+TB_i(T)\right)^{(-1)^i}
 \in\Lambda^\times
\end{equation}
satisface \(\rho_{E,\ell}(0)=1\).
\end{lemma}

\begin{proof}
The simplicial identity
\(\partial\Delta_{\mathrm{AW}}
=(\partial\otimes I+(-1)^{\deg}I\otimes\partial)
\Delta_{\mathrm{AW}}\), together with the compatibility of
\(P_E^K\) with faces and degeneracies, proves that
\eqref{eq:amp-bsd-phi-kato-fers} commutes with the differentials. On the associated
graded object, lemma
\ref{lem:amp-bsd-bases-emparejadas-kato-fers} identifies all the
generators and gives the identity. In those bases,
\eqref{eq:amp-bsd-phi-en-cada-generador} is precisely
\eqref{eq:amp-bsd-matriz-kato-fers}.

The completeness of \(\Lambda=K[[T]]\) makes the series converge
\begin{equation}\label{eq:amp-bsd-inversa-kato-fers}
 (I+TB_i(T))^{-1}
 =\sum_{n\geq0}(-TB_i(T))^n;
\end{equation}
these inverses also commute with the differentials, so they form the
based inverse of \(\Phi_{K/\mathrm{FERS}}\). Taking the alternating
determinant yields \eqref{eq:amp-bsd-rho-determinantal}; evaluating at
\(T=0\) gives one. Thus, the normalization is not postulated from the leading
term: it is the constant coordinate of an equivalence that is the identity
on the root unit.
\end{proof}

The first comparator is the determinant of
\(\Phi_{K/\mathrm{FERS}}\). If \(\mathfrak z_K\) denotes the element
coming from the unit \(1_E\), its image is characterized by
\begin{equation}\label{eq:amp-bsd-comparador-kato}
 \mathfrak c_{K/\mathrm{FERS}}(\mathfrak z_K)
 =\left[T^{-d}\rho_{E,\ell}(T)\det S_E(T)\right]_{T=0},
 \qquad
 \rho_{E,\ell}(0)=1.
\end{equation}
The unit of \(\rho_{E,\ell}\) is oriented and equals one at the root; it therefore does not
alter the order or the basis. Theorem
\ref{thm:bsd-suc-bockstein} transforms \eqref{eq:amp-bsd-comparador-kato} into
\begin{equation}\label{eq:amp-bsd-kato-bockstein}
 \mathfrak c_{K/\mathrm{FERS}}(\mathfrak z_K)
 =R_{\mathrm{Boc}}^{\mathrm{der}}(S_E)
 =\tau_{\mathrm{reg}}\otimes\bigotimes_{q\geq1}\tau_q.
\end{equation}

The Poitou--Tate comparator is constructed page by page. For
\(q\geq1\), the isomorphism \(j_q\) of
\eqref{eq:bsd-suc-jq} and the reduced Bockstein
\(\bar\beta_q\) induce
\begin{equation}\label{eq:amp-bsd-comparador-pt-q}
 \mathfrak c_{\mathrm{PT},q}(\tau_q)
 :=\det\!\left(j_q\circ\bar\beta_q\right)
 =\det h_E^{(q)}.
\end{equation}
The regular direction is transported through
\(\det\bar S_E(0)=\tau_{\mathrm{reg}}\). Knudsen--Mumford gluing and the
Koszul signs defined in \eqref{eq:bsd-suc-rboc} provide the unique
oriented isomorphism
\begin{equation}\label{eq:amp-bsd-comparador-pt}
 \mathfrak c_{\mathrm{PT}}
 \left(\tau_{\mathrm{reg}}\otimes\bigotimes_{q\geq1}\tau_q\right)
 =\tau_{\mathrm{reg}}^{\mathrm{ht}}
  \otimes\bigotimes_{q\geq1}\det h_E^{(q)}.
\end{equation}
On the first stable page, the form \(h_E^{(1)}\) is the
Poincaré--Néron--Tate form on the free Mordell--Weil lattice; its determinant is
\begin{equation}\label{eq:amp-bsd-regulador-nt}
 \det h_E^{(1)}=\operatorname{Reg}(E)
\end{equation}
in the basis induced by the root unit.

\begin{lemma}[Poitou--Tate unfolding of the arithmetic degree]
\label{lem:amp-bsd-grado-pt}
Let
\[
 \operatorname{Flag}_{\mathrm{Boc}}(E)
 :=\bigoplus_{q\geq1}V_q
\]
the finite flag of positive pages, each with the basis inherited from the
root unit. Orthogonal reduction and the isomorphisms \(j_q\) construct a
based isomorphism
\begin{equation}\label{eq:amp-bsd-psi-pt}
 \Psi_{\mathrm{PT}}:
 \bigl(E(\mathbb Q)/E(\mathbb Q)_{\mathrm{tors}}\bigr)\otimes K
 \xrightarrow{\ \sim\ }\operatorname{Flag}_{\mathrm{Boc}}(E).
\end{equation}
In particular,
\begin{equation}\label{eq:amp-bsd-rango-bandera}
 \operatorname{rank}E(\mathbb Q)
 =\sum_{q\geq1}\dim_KV_q
 =\operatorname{ord}_T\det S_E(T).
\end{equation}
\end{lemma}

\begin{proof}
The reduced complex is filtered by the powers of the augmentation ideal. In the
quotient at level \(q\), the regular direction has already been separated by
\(\bar S_E(0)\), and the only surviving block is \(V_q\). The equivalence
\(X^{\mathrm{orth}}\) identifies the opposite block with its dual, while
\(j_q\bar\beta_q\) provides the perfect pairing that lifts the
basis of the quotient to the preceding level. Starting at the last nonzero page and
repeating this lifting yields \(\Psi_{\mathrm{PT}}\).

In the bases ordered by depth, the matrix of the resulting map is
block triangular and all its diagonal blocks are identities: in the
first block because \(z_K\) and \(z_{\mathrm{PT}}\) have the same root, and
in the others because \(j_q\) and \(\bar\beta_q\) are based isomorphisms.
Therefore \(\Psi_{\mathrm{PT}}\) is invertible and loses no page.
Taking dimensions and using
\eqref{eq:bsd-suc-carry-orden} yields
\eqref{eq:amp-bsd-rango-bandera}. On taking determinants, the same
lifting is exactly the gluing of
\eqref{eq:amp-bsd-comparador-pt}; thus, equality of ranks and equality of
lines come from a single based map.
\end{proof}

\subsection{Localization triangle, Smith rank and finite factors}
\label{subsec:amp-bsd-smith-sha}

For each place \(v\), the finite local condition is a based subcomplex
\(\mathcal C_{E,v}^{\mathrm{fin}}\subset\mathcal C_{E,v}\), and one defines
\(\mathcal C_{E,v}^{\mathrm{sing}}\) by the triangle
\begin{equation}\label{eq:amp-bsd-triangulo-local}
 \mathcal C_{E,v}^{\mathrm{fin}}
 \xrightarrow{\ i_v\ }\mathcal C_{E,v}
 \longrightarrow\mathcal C_{E,v}^{\mathrm{sing}}
 \longrightarrow\mathcal C_{E,v}^{\mathrm{fin}}[1].
\end{equation}
If \(\operatorname{res}_v\) denotes the global restriction, the localization
map is, at chain level,
\begin{equation}\label{eq:amp-bsd-mapa-localizacion}
 \ell_E:
 \mathcal C_E^{\mathrm{glob,red}}\oplus
 \bigoplus_v\mathcal C_{E,v}^{\mathrm{fin}}
 \longrightarrow\bigoplus_v\mathcal C_{E,v},
 \qquad
 \ell_E\bigl(x,(y_v)_v\bigr)
 :=\bigl(\operatorname{res}_v x-i_vy_v\bigr)_v .
\end{equation}
The definition by cone of the reduced Selmer complex gives the triangle
\begin{equation}\label{eq:amp-bsd-triangulo-localizacion}
 \mathcal C_E^{\mathrm{red}}
 \longrightarrow
 \mathcal C_E^{\mathrm{glob,red}}\oplus
 \bigoplus_v\mathcal C_{E,v}^{\mathrm{fin}}
 \xrightarrow{\ \ell_E\ }
 \bigoplus_v\mathcal C_{E,v}
 \longrightarrow\mathcal C_E^{\mathrm{red}}[1].
\end{equation}
Therefore, local--global exactness comes from the cone of a map written in
\eqref{eq:amp-bsd-mapa-localizacion}; it is not subsequently added as a relation
between cardinalities.

In \eqref{eq:amp-bsd-triangulo-localizacion} the common root summand
is canceled, which is the same in both channels by
\eqref{eq:bsd-suc-raiz}. The resulting complement is denoted by
\begin{equation}\label{eq:amp-bsd-cono-local-reducido}
 \mathcal C_E^{\mathrm{loc,red}}
 :=q_{\mathrm{root}}\operatorname{Cone}(\ell_E)[-1].
\end{equation}
After separating the free Mordell--Weil lattice and its dual through
\(X^{\mathrm{orth}}\), a choice of the already oriented integral bases
represents this complex by
\begin{equation}\label{eq:amp-bsd-complejo-smith}
 \mathcal C_E^{\mathrm{loc,red}}
 \simeq[M_E\xrightarrow{\ D_E\ }N_E].
\end{equation}

\begin{lemma}[Rational acyclicity and full rank]
\label{lem:amp-bsd-aciclicidad-rango}
The complex of \eqref{eq:amp-bsd-complejo-smith} is acyclic after
tensoring with \(\mathbb Q\). In particular,
\[
 D_E\otimes\mathbb Q:M_E\otimes\mathbb Q
 \xrightarrow{\ \sim\ }N_E\otimes\mathbb Q
\]
is an isomorphism; \(M_E\) and \(N_E\) have the same rank and \(D_E\) has
full rank.
\end{lemma}

\begin{proof}
Let \(x\) be a cycle of the complex
\(\mathcal C_E^{\mathrm{loc,red}}\otimes\mathbb Q\). The reduction has
eliminated the summand \(p_{\mathrm{root}}\), hence
\(q_{\mathrm{root}}x=x\). In the oriented integral bases used in
\eqref{eq:amp-bsd-complejo-smith}, the finite--singular complement
decomposes into the normal blocks of
\cref{lem:bsd-suc-forma-normal}. In each block \(\lambda\), the root
coefficient is already zero and the cycle is written
\[
 x_{\lambda}=u_{\lambda}e_{\lambda}
                  +v_{\lambda}b_{\lambda},
 \qquad
 \partial x_{\lambda}
   =(u_{\lambda}-3c_{\lambda}v_{\lambda})g_{\lambda}=0.
\]
Since \(g_{\lambda}\) belongs to the free basis, we have
\(u_{\lambda}=3c_{\lambda}v_{\lambda}\), and the explicit formula for the
differential gives
\[
 x_{\lambda}
   =v_{\lambda}(3c_{\lambda}e_{\lambda}+b_{\lambda})
   =\partial(v_{\lambda}f_{\lambda}).
\]
The sum is finite in each degree, so
\(x=\partial\bigl(\sum_{\lambda}v_{\lambda}f_{\lambda}\bigr)\).
Every cycle is therefore a boundary by the same normal form that produced
the filler \eqref{eq:bsd-suc-hstar}, without introducing
\(H_{\mathrm{FS}}\) as a premise. The equality
\(p_{\mathrm{root}}z_K=z_{\mathrm{PT}}\) is what allows the root
direction to be canceled before this calculation; without it a degree-zero
class would remain. The self-dual equivalence
\(X^{\mathrm{orth}}:\mathcal C_E^{\mathrm{red}}\simeq
D_3(\mathcal C_E^{\mathrm{red}})\) pairs the two remaining degrees and
transports their bases with complementary opposite orientation. From this one
obtains \(\operatorname{rank}M_E=\operatorname{rank}N_E\).
Finally, acyclicity of a two-term complex of equal rank
is equivalent to \(D_E\otimes\mathbb Q\) being invertible.
\end{proof}

The lemma makes it possible to apply the Smith normal form without presupposing any of
its divisors. There exist oriented changes of basis \(U,V\) such that
\begin{equation}\label{eq:amp-bsd-morfismo-smith}
 U D_E V=\operatorname{diag}(s_1,\ldots,s_t),
 \qquad s_i\in\mathbb Z\setminus\{0\}.
\end{equation}
Therefore,
\begin{equation}\label{eq:amp-bsd-grupo-smith}
 F_E:=\operatorname{coker}D_E
 \simeq\bigoplus_{i=1}^{t}\mathbb Z/|s_i|\mathbb Z,
 \qquad
 |F_E|=\prod_{i=1}^{t}|s_i|<\infty.
\end{equation}

\begin{lemma}[Finite localization sequence]
\label{lem:amp-bsd-exactitud-sha-tamagawa}
The torsion fragment of the long cohomology sequence of
\eqref{eq:amp-bsd-triangulo-localizacion} is
\begin{equation}\label{eq:amp-bsd-exacta-sha-tamagawa}
 0\longrightarrow\operatorname{Sha}(E)
 \xrightarrow{\ \iota_{\mathrm{Sha}}\ }F_E
 \xrightarrow{\ \partial_{\mathrm{loc}}\ }
 \bigoplus_v\Phi_v(\mathbb F_v)
 \longrightarrow0,
 \qquad c_v=|\Phi_v(\mathbb F_v)|.
\end{equation}
\end{lemma}

\begin{proof}
A global cocycle represents a class of
\(\operatorname{Sha}(E)\) precisely when its restriction at each place
admits a primitive in
\(\mathcal C_{E,v}^{\mathrm{fin}}\). The class of the pair formed by that
cocycle and its local primitives in the cone of
\eqref{eq:amp-bsd-mapa-localizacion} defines
\(\iota_{\mathrm{Sha}}\). If this class is a boundary, the global component is a
boundary and the Tate--Shafarevich class is zero; thus,
\(\iota_{\mathrm{Sha}}\) is injective.

The map \(\partial_{\mathrm{loc}}\) takes a class of the cone, restricts it to
each quotient
\(\mathcal C_{E,v}^{\mathrm{sing}}\) of
\eqref{eq:amp-bsd-triangulo-local} and preserves its component class. The
identification
\[
 H^0(\mathcal C_{E,v}^{\mathrm{sing}})_{\mathrm{tors}}
 \simeq\Phi_v(\mathbb F_v)
\]
is induced by the finite--singular quotient, not by the order of the group.
A class of the cone has zero residue exactly when its local
representatives can be chosen in the finite subcomplexes; by the preceding
definition, that kernel is the image of \(\iota_{\mathrm{Sha}}\).
Finally, the last arrow of
\eqref{eq:amp-bsd-triangulo-local} lifts each component class to the
local complex. The exactness of the cone
\eqref{eq:amp-bsd-triangulo-localizacion}, after removing the two
free summands paired by \(X^{\mathrm{orth}}\), turns that
lifting into a preimage under \(\partial_{\mathrm{loc}}\). This proves
surjectivity and, consequently, all of
\eqref{eq:amp-bsd-exacta-sha-tamagawa}.
\end{proof}

Now, and only after lemma
\ref{lem:amp-bsd-aciclicidad-rango}, the group \(F_E\) is finite. The
injection of \eqref{eq:amp-bsd-exacta-sha-tamagawa} then proves the
finiteness of \(\operatorname{Sha}(E)\). Taking orders in that exact
sequence yields
\begin{equation}\label{eq:amp-bsd-cardinal-smith}
 |F_E|=|\operatorname{Sha}(E)|\prod_v c_v.
\end{equation}
This is the comparator
\(\mathfrak c_{\mathrm{STS}}\): it transports the Smith elementary factors
to the Tate--Shafarevich and Tamagawa factors while preserving their order in the line,
instead of inserting their cardinalities once the leading term has been calculated.

The torsion part is constructed by applying the determinant to the Mordell--Weil
lattice and its Pontryagin dual. The two finite factors are dual and
contribute
\begin{equation}\label{eq:amp-bsd-comparador-torsion}
 \mathfrak c_{\mathrm{tors}}:quad
 u\longmapsto
 \frac{u}{|E(\mathbb Q)_{\mathrm{tors}}|^2}.
\end{equation}
The Archimedean comparator is obtained by integrating the oriented Néron differential
over the based real cycle:
\begin{equation}\label{eq:amp-bsd-comparador-periodo}
 \Omega_E:=\int_{E(\mathbb R)}|\omega_E|,
 \qquad
 \mathfrak c_\infty:quad u\longmapsto\frac{u}{\Omega_E}.
\end{equation}
The differential and the cycle are fixed before evaluating \(L(E,s)\); therefore
\eqref{eq:amp-bsd-comparador-periodo} does not retrospectively normalize the
analytic section.

\subsection{Causal separation of primitives, identities and consequences}
\label{subsec:amp-bsd-separacion-causal}

The generator-by-generator criterion of
lemma \ref{lem:amp-bsd-bases-emparejadas-kato-fers} audits in each degree the
equivalence already constructed. It does not condition its existence. The logical order
of the package is
\begin{equation}\label{eq:amp-bsd-diagrama-causal}
\begin{array}{c}
 \begin{gathered}
  1_E,\ \Delta_{\mathrm{AW}},\ P_E^K,\ P_E^{\mathrm{PT}},
  \ h_*=-vf,\ X^{\mathrm{orth}},\\
  \{\,\mathcal C_{E,v}^{\mathrm{fin}}\to\mathcal C_{E,v}\,\}_v,\\
  \text{Mordell--Weil and Pontryagin bases;}\\
  \text{Néron differential and cycle}
 \end{gathered}
 \\[2mm]\Downarrow\\[-1mm]
 \begin{gathered}
  \Phi_{K/\mathrm{FERS}}\text{ equivalence},\quad
  \rho_{E,\ell}(0)=1,\quad
  p_{\mathrm{root}}z_K=z_{\mathrm{PT}},\\
  R_{\mathrm{Boc}}^{\mathrm{der}}=\operatorname{LT}_T(S_E),\quad
  \det(j_q\bar\beta_q)=\det h_E^{(q)},\\
  H^*(\mathcal C_E^{\mathrm{loc,red}}\otimes\mathbb Q)=0,\quad
  D_E\otimes\mathbb Q\text{ isomorphism},\\
  0\to\operatorname{Sha}(E)\to F_E
  \to\bigoplus_v\Phi_v(\mathbb F_v)\to0
 \end{gathered}
 \\[2mm]\Downarrow\\[-1mm]
 \begin{gathered}
  d_{\mathrm{an}}=\operatorname{rank}E(\mathbb Q),\qquad
  \operatorname{Reg}(E),\qquad
  |\operatorname{Sha}(E)|\prod_vc_v,\\
  |E(\mathbb Q)_{\mathrm{tors}}|^{-2},\qquad
  \Omega_E^{-1},\qquad
  \text{equality of the leading term}.
 \end{gathered}
\end{array}
\end{equation}
The first row contains only already based objects and maps, together with
the indicated generator-by-generator verification; the second, identities
proved by those maps; the third, their consequences on the determinant
line. In particular, no factor in the last row is
used to retroactively normalize an arrow in the first.

\begin{proposition}[Exact package of comparators]
\label{prop:amp-bsd-paquete-comparadores-exacto}
The comparators of
\eqref{eq:amp-bsd-cadena-lineas} come from a single package of maps of
complexes and satisfy, in the order of construction,
\begin{equation}\label{eq:amp-bsd-paquete-identidades}
 \begin{gathered}
 \Phi_{K/\mathrm{FERS}}(z_K)=z_{\mathrm{Iw}},\qquad
 \operatorname{gr}^{F}\Phi_{K/\mathrm{FERS}}=I,\\
 \det\nolimits_{\mathrm{alt}}\Phi_{K/\mathrm{FERS}}
   =\rho_{E,\ell}(T),\qquad \rho_{E,\ell}(0)=1,\\
 \det(j_q\bar\beta_q)=\det h_E^{(q)}\quad(q\geq1),\\
 \Psi_{\mathrm{PT}}:
  \bigl(E(\mathbb Q)/E(\mathbb Q)_{\mathrm{tors}}\bigr)\otimes K
  \xrightarrow{\sim}\operatorname{Flag}_{\mathrm{Boc}}(E),\\
 p_{\mathrm{root}}z_K=z_{\mathrm{PT}},\qquad
 h_*=-vf,\quad \partial h_*=z_{\mathrm{PT}}-z_K,\\
 H^*(\mathcal C_E^{\mathrm{loc,red}}\otimes\mathbb Q)=0,\qquad
 U D_E V=\operatorname{diag}(s_1,\ldots,s_t),\\
 0\longrightarrow\operatorname{Sha}(E)
  \longrightarrow F_E
  \longrightarrow\displaystyle\bigoplus_v\Phi_v(\mathbb F_v)
  \longrightarrow0 .
 \end{gathered}
\end{equation}
The induced composition on determinant lines is exactly
\begin{equation}\label{eq:amp-bsd-paquete-composicion}
 \det(\Phi_{K/\mathrm{FERS}})
 \xrightarrow{\ \det(j_q\bar\beta_q)\ }
 \det(D_E)
 \xrightarrow{\ \mathrm{Smith}\ }
 \frac{|\operatorname{Sha}(E)|\prod_vc_v}
      {|E(\mathbb Q)_{\mathrm{tors}}|^2\,\Omega_E},
\end{equation}
with the Knudsen--Mumford signs and the orientations of
\eqref{eq:amp-bsd-cadena-lineas}. In particular, the right-hand side of
\eqref{eq:amp-bsd-paquete-composicion} is the final coordinate of the transported
element, not data used to construct any of the arrows.
\end{proposition}

\begin{proof}
The first row of \eqref{eq:amp-bsd-paquete-identidades} is
\eqref{eq:amp-bsd-phi-kato-fers} and
\eqref{eq:amp-bsd-matriz-kato-fers}--
\eqref{eq:amp-bsd-rho-determinantal}. The second is
\eqref{eq:amp-bsd-comparador-pt-q} and lemma
\ref{lem:amp-bsd-grado-pt}. The third combines equality of the root with
the explicit filler \eqref{eq:bsd-suc-hstar}; its linear extension in the
normal form induces the auxiliary reader used in lemma
\ref{lem:amp-bsd-aciclicidad-rango}. On restricting to the complement of the
root, that reduction contracts every rational cycle. The last row is
\eqref{eq:amp-bsd-morfismo-smith} and
\eqref{eq:amp-bsd-exacta-sha-tamagawa}. Applying the determinant functor to
those identities, in that same order, gives
\eqref{eq:amp-bsd-paquete-composicion}. No equality is obtained
by first comparing the two final scalars.
\end{proof}

\begin{lemma}[Causal independence of the leading term]
\label{lem:amp-bsd-independencia-termino-principal}
Let \(\mathscr P_E\) be the data set
\begin{equation}\label{eq:amp-bsd-primitivas-paquete}
 \mathscr P_E:=
 \left(
  1_E,\Delta_{\mathrm{AW}},P_E^K,P_E^{\mathrm{PT}},
  h_*=-vf,X^{\mathrm{orth}},
  \{\mathcal C_{E,v}^{\mathrm{fin}}\xrightarrow{i_v}
     \mathcal C_{E,v}\}_v,
  \text{bases and orientations},
  \omega_E,\gamma_E^{\mathbb R}
 \right),
\end{equation}
where \(\omega_E\) is the Néron differential and
\(\gamma_E^{\mathbb R}\) the oriented real cycle. All maps of the
package in proposition
\ref{prop:amp-bsd-paquete-comparadores-exacto} are determined by
\(\mathscr P_E\). The following numerals are not part of \(\mathscr P_E\)
\[
 L^{(r)}(E,1),\quad r,\quad\operatorname{Reg}(E),\quad
 |\operatorname{Sha}(E)|,\quad(c_v)_v,\quad
 |E(\mathbb Q)_{\mathrm{tors}}|,\quad\Omega_E.
\]
Those values are obtained, respectively, by evaluating the analytic section,
taking the degree of the flag, calculating determinants and Smith forms,
taking orders of finite groups and integrating
\(\omega_E\) over \(\gamma_E^{\mathbb R}\).
\end{lemma}

\begin{proof}
The formula \eqref{eq:amp-bsd-phi-kato-fers} uses only
\(1_E,\Delta_{\mathrm{AW}},P_E^K\) and the carry. The Bockstein and
Poitou--Tate maps use the filtration, \(P_E^{\mathrm{PT}}\), the bases and
self-duality. The cone of \eqref{eq:amp-bsd-mapa-localizacion} uses the
maps \(i_v\) and the local restrictions; its contraction and Smith
form use the explicit filler \eqref{eq:bsd-suc-hstar}, its auxiliary linear
extension, \(X^{\mathrm{orth}}\) and the integral bases. The last
two comparators use the torsion lattices and the
pair \((\omega_E,\gamma_E^{\mathbb R})\). This is an exhaustive enumeration
of the arguments appearing in the construction formulas. The
numerals listed in the statement appear only after applying
cohomology, determinant, cardinality or integration, so they cannot
retrospectively select the maps.
\end{proof}

\subsection{Composition, degrees and leading term}
\label{subsec:amp-bsd-composicion}

Define the total comparator by composition in the written order:
\begin{equation}\label{eq:amp-bsd-comparador-total}
 \mathfrak C_E:=
 \mathfrak c_\infty\circ
 \mathfrak c_{\mathrm{tors}}\circ
 \mathfrak c_{\mathrm{STS}}\circ
 \mathfrak c_{\mathrm{PT}}\circ
 \mathfrak c_{K/\mathrm{FERS}}.
\end{equation}
Since \(p_{\mathrm{root}}z_K=z_{\mathrm{PT}}\), all arrows receive the
same based unit. The filler \(h_*=-vf\) of
\eqref{eq:bsd-suc-hstar} explicitly fills the difference; its linear
extension in the normal basis produces the auxiliary homotopic reader and proves that
no additional class modifies the element of the line during
transport. The reader is not a separate premise.

\begin{theorem}[Construction of the based comparators]
\label{thm:amp-bsd-comparadores-construidos}
For every elliptic curve \(E/\mathbb Q\) with the complexes, bases,
orientations and local data defined above, the maps
\eqref{eq:amp-bsd-phi-kato-fers},
\eqref{eq:amp-bsd-comparador-pt},
\eqref{eq:amp-bsd-triangulo-localizacion},
\eqref{eq:amp-bsd-comparador-torsion} and
\eqref{eq:amp-bsd-comparador-periodo}, together with lemmas
\ref{lem:amp-bsd-aciclicidad-rango} and
\ref{lem:amp-bsd-exactitud-sha-tamagawa}, construct the five comparators of
\eqref{eq:amp-bsd-cadena-lineas}. Their composition preserves the unit, degree
and orientation, and satisfies
\begin{equation}\label{eq:amp-bsd-igualdad-grados}
 d_{\mathrm{an}}=\operatorname{ord}_T\det S_E(T)
 =d_{\mathrm{ar}}=\operatorname{rank}E(\mathbb Q).
\end{equation}
Moreover, \(\operatorname{Sha}(E)\) is finite and the equality of the distinguished
element in \(\mathcal L_E\) is published as
\begin{equation}\label{eq:amp-bsd-termino-principal}
 \frac{L^{(r)}(E,1)}{r!\,\Omega_E}
 =\frac{|\operatorname{Sha}(E)|\,\operatorname{Reg}(E)
        \prod_v c_v}
       {|E(\mathbb Q)_{\mathrm{tors}}|^2},
 \qquad r=\operatorname{rank}E(\mathbb Q).
\end{equation}
\end{theorem}

\begin{proof}
The Kato--FERS comparator and the condition
\(\rho_{E,\ell}(0)=1\) identify the analytic order with
\(d=\operatorname{ord}_T\det S_E(T)\) and the initial term with
\(R_{\mathrm{Boc}}^{\mathrm{der}}(S_E)\). The identity
\eqref{eq:bsd-suc-lt} preserves the regular page and every positive page. The
isomorphisms \(j_q\) transport those pages to the derived heights.
Lemma \ref{lem:amp-bsd-grado-pt} constructs the isomorphism of the flag with the
free Mordell--Weil lattice and, together with
\eqref{eq:bsd-suc-carry-orden}, proves
\eqref{eq:amp-bsd-igualdad-grados}.

Lemma \ref{lem:amp-bsd-aciclicidad-rango} deduces rational acyclicity
directly from the explicit filler \eqref{eq:bsd-suc-hstar}, its linear
reduction of the complement, cancellation of the common root and
\(X^{\mathrm{orth}}\); from this the full rank of \(D_E\) is obtained, not
as an additional hypothesis. The Smith form
\eqref{eq:amp-bsd-morfismo-smith} then produces the finite group \(F_E\).
Lemma \ref{lem:amp-bsd-exactitud-sha-tamagawa} extracts from the triangles
\eqref{eq:amp-bsd-triangulo-local} and
\eqref{eq:amp-bsd-triangulo-localizacion} the finite exact sequence; only
afterwards are the finiteness of \(\operatorname{Sha}(E)\) and
\eqref{eq:amp-bsd-cardinal-smith} deduced. Finally,
\eqref{eq:amp-bsd-regulador-nt},
\eqref{eq:amp-bsd-comparador-torsion} and
\eqref{eq:amp-bsd-comparador-periodo} publish, respectively, the regulator,
the torsion quotient and the period. The root equality eliminates any relative
unit between the two publications. Evaluating the same based element at
both ends of \eqref{eq:amp-bsd-comparador-total} yields
\eqref{eq:amp-bsd-termino-principal}.
\end{proof}

\subsection{Premises and refutation criteria}
\label{subsec:amp-bsd-falsadores}

The construction starts from the genealogical unit \(1_E\) and the Alexander--Whitney
diagonal. It also uses the
coupled publications, all the
Bocksteins and the regular page, the explicit filler \(h_*=-vf\), the
reduced self-duality \(X^{\mathrm{orth}}\), the chain maps of the
localization triangle and the oriented bases of torsion and period. The
Kato--FERS equivalence and its normalization to one, the full rank of \(D_E\), the
exactness of \eqref{eq:amp-bsd-exacta-sha-tamagawa} and the finiteness of
\(\operatorname{Sha}(E)\) are conclusions of the preceding lemmas, not
independent premises. The generator-by-generator criterion of lemma
\ref{lem:amp-bsd-bases-emparejadas-kato-fers} is a nongoverning check
that locally audits the first comparator. The criteria that would refute a specific arrow
are:
\begin{enumerate}
 \item a root coefficient different from one, which would introduce a relative
       unit between \(z_K\) and \(z_{\mathrm{PT}}\);
 \item a noninvertible \(j_q\) or the omission of a Bockstein page, which
       would alter the degree or the height determinant;
 \item a zero Smith divisor, which would prevent the finiteness of \(F_E\);
 \item failure of exactness of
       \eqref{eq:amp-bsd-exacta-sha-tamagawa}, which would break the identification
       of the finite factors;
 \item an orientation reversal in torsion or period, which would change the
       based element even if it preserved its unoriented line.
\end{enumerate}
The memoryless Manin projection and visible rescaling modulo three
violate the first criterion; they therefore constitute necessity proofs for
those reduced models, not premises of theorem
\ref{thm:amp-bsd-comparadores-construidos}.


\begin{theorem}[Local--global determinantal comparison]
\label{thm:bsd-suc-local-global}
Let the based comparators already constructed in the
\cref{thm:amp-bsd-comparadores-construidos} be:
Kato--FERS/Iwasawa,
Poitou--Tate/Poincaré--Néron--Tate,
Smith--Tamagawa--Shafarevich, torsion and period.
If \(d_{\mathrm{an}}\) is the degree of the analytic section at the
central specialization and \(d_{\mathrm{ar}}\) the degree of the based Selmer
line, then
\begin{equation}\label{eq:bsd-suc-grados}
 d_{\mathrm{an}}=d=d_{\mathrm{ar}},
\end{equation}
and one obtains the equality of distinguished elements
\begin{equation}\label{eq:bsd-suc-linea-final}
 \mathfrak z_{\mathrm{an}}(E)=\mathfrak z_{\mathrm{ar}}(E)
 \quad\text{in}\quad\mathcal L_E.
\end{equation}
The analytic side is the publication of the initial term; the arithmetic
side is the publication of the complete regulator, the Smith
elementary factors,
the local factors, the torsion and the period.
\end{theorem}

\begin{proof}
The based construction gives
\([\Phi_{K/\mathrm{FERS}}^i](T)=I+TB_i(T)\) in each degree. The based
Kato--FERS comparator then expresses the corrected analytic section
as
\begin{equation}\label{eq:bsd-suc-fers}
 L_{\ell}^{\mathrm{corr}}(E,T)
 =\rho_{E,\ell}(T)\det S_E^{\mathrm{corr}}(T),
 \qquad \rho_{E,\ell}(0)=1.
\end{equation}
By \eqref{eq:bsd-suc-orden}, its central degree is \(d\); by
\eqref{eq:bsd-suc-lt}, its initial term is exactly
\(R_{\mathrm{Boc}}^{\mathrm{der}}\), not that regulator multiplied by an
undetermined unit. The isomorphisms \(j_q\) of
\eqref{eq:bsd-suc-jq} transport each \(\tau_q\) to the determinant of the
derived height of its page. The regular page transports the acyclic
summand that the positive pages do not see.

The Poincaré--Néron--Tate comparison identifies the degree of that height
line with \(d_{\mathrm{ar}}\). The Smith elementary factors decompose its
finite part into the Shafarevich and Tamagawa factors; the quotient of the integral
lattices supplies the square of the torsion, and the Archimedean reader supplies
the period. All these maps act on the same root line. By
theorem \ref{thm:bsd-suc-raiz-filler}, no relative unit remains between
the Kato and Poitou--Tate publications. The composition preserves the degree and
the based element, proving simultaneously
\eqref{eq:bsd-suc-grados} and \eqref{eq:bsd-suc-linea-final}. The corresponding scalar
formula is reserved for final recognition.
\end{proof}

\subsection*{Hypotheses and scope}

The BSD chain uses: the complete enriched state and its Alexander--Whitney
coalgebra; the fiber product \eqref{eq:bsd-suc-pullback}; the two
publications constructed on the same unit; the root line, its basis and its
counit; the perfect and saturated reduced quotient, together with
\(X^{\mathrm{orth}}\); the complete finite--singular complex; the finite
Bockstein flags, their Koszul signs and the regular page; and the
height, Smith, Tamagawa, Shafarevich, torsion and period comparisons
with a common orientation. The degree-by-degree verification of the paired
bases and the residuals \(\varepsilon_i(b)\) is recorded as a
local control of the realization, not as a governing premise of the closure.
Under this explicit starting structure,
\eqref{eq:bsd-suc-raiz}, \eqref{eq:bsd-suc-hstar},
\eqref{eq:bsd-suc-lt} and \eqref{eq:bsd-suc-linea-final} are exact
equalities.

Unit preservation, the two-sheet traversal and local--global
publication are formalized through the chain coalgebra, the genealogical
fiber product, the coupled publication, orthogonal reduction and the
Bockstein tower. The auxiliary identities separately verify the group-like element property,
preservation of history under \(\widehat\Pi_0\), the identity of the
unit in both channels, the reduced involution, the finite--singular normal
form and the determinantal identity of arbitrary order.

\paragraph{BSD ablation controls.}
\noindent\textbf{Status: NON\_GOVERNING\_CONTROL.}
The following three ablations receive the already proved identities and
falsify only reduced models. They contribute no premises to the owner,
do not condition its closure and have no causal return edge toward it.

\begin{ablation}[Manin projection without history]
\label{abl:bsd-suc-manin}
Replacing \(P_E^{\mathrm{gen}}\) with \(P_E^{\mathrm{Man}}\) erases data that
the coupled publication needs. Indeed, the route
\(\gamma=(\mathsf N\mathsf E)^{54}\) may have the same cell and the same
visible phase as the trivial route, but preserves distinct winding and memory
in \(G_{m,n}\). The Manin augmentation identifies them; the second
projection of \eqref{eq:bsd-suc-pullback} does not.
\end{ablation}

\begin{proof}
Every based lift to the Manin complex factors, up to homotopy, through the
section of the augmentation \(s\varepsilon_{\mathrm{AW}}\). It therefore depends only
on the visible coordinate. By contrast, the second component of
\(\widehat\Pi_0(c)\) is \(c\) itself, so it preserves the history
simplex and distinguishes the winding of \(\gamma\). The ablation refutes the
bare projection; it does not refute the fiber product used in the theorem.
\end{proof}

\begin{ablation}[Rescaling of the root visible only modulo three]
\label{abl:bsd-suc-twist}
In the specialization \(c=1\) of
\eqref{eq:bsd-suc-diferencial-normal}, the altered cycle
\begin{equation}\label{eq:bsd-suc-twist4}
 z_K^{(4)}=4r+3e+b
\end{equation}
satisfies
\begin{equation}\label{eq:bsd-suc-twist-defecto}
 z_{\mathrm{PT}}-z_K^{(4)}
 =\partial(-f)-3r.
\end{equation}
The defect \(-3r\) disappears modulo three and reappears modulo nine.
Consequently, \eqref{eq:bsd-suc-twist4} preserves a partial shadow, but not
the based unit or \eqref{eq:bsd-suc-raiz}.
\end{ablation}

\begin{proof}
The equality \(\partial f=3e+b\) immediately gives
\eqref{eq:bsd-suc-twist-defecto}. Modulo three, the second summand vanishes;
modulo nine it is a root class not divisible by a boundary of the complement.
Moreover, its root coordinate is four, whereas
\eqref{eq:bsd-suc-normalizacion-raiz} requires one. The example verifies the
necessity of the coupled pair and orientation; it introduces no exception
to the theorem that preserves both data.
\end{proof}

\begin{ablation}[Loss of regulator pages]
\label{abl:bsd-suc-bockstein}
The matrix \(\operatorname{diag}(T^2,T)\) shows that the first Bockstein does not
recover the higher-order carry. The matrix
\(\operatorname{diag}(2,T^2)\) has positive trivialization
\(\tau_2=1\), but
\(\tau_{\mathrm{reg}}=2=\operatorname{LT}_{T}(S)\). Therefore, neither the
first page nor the product of positive pages replaces the complete
regulator \eqref{eq:bsd-suc-rboc}.
\end{ablation}

\begin{proof}
In the first example, \(d=3\), \(r_0=2\) and
\(c_I^{\deg}=1\); page one does not contain that last degree of memory. In
the second, the only singular divisor is \(T^2\), but the regular direction
contributes the factor two. Omitting it would produce one instead of two. Both calculations
are diagonal cases of \eqref{eq:bsd-suc-smith} and verify the necessity of
all the pieces of \eqref{eq:bsd-suc-rboc}.
\end{proof}
\section*{Final mathematical recognition of Birch--Swinnerton--Dyer}

The preceding development concludes before employing its historical formulation
as a criterion.

In the domain declared by theorem
\ref{thm:bsd-suc-local-global}, the Smith comparison identifies the
finite component of the arithmetic line with the
Shafarevich--Tamagawa decomposition and proves the finiteness of
\(\operatorname{Sha}(E)\). Consequently, its order in the following scalar
publication is the cardinality of the identified finite group, not a formal
symbol added after the fact.

The scalar translation of \eqref{eq:bsd-suc-linea-final} provides, with
\(r=\operatorname{rank}E(\mathbb Q)\),
\begin{equation}\label{eq:bsd-suc-rango-final}
 \operatorname{ord}_{s=1}L(E,s)=r
\end{equation}
and
\begin{equation}\label{eq:bsd-suc-formula-final}
 \frac{L^{(r)}(E,1)}{r!\,\Omega_E}
 =\frac{
   \lvert\operatorname{Sha}(E)\rvert\,\operatorname{Reg}(E)\prod_v c_v}
  {\lvert E(\mathbb Q)_{\mathrm{tors}}\rvert^2}.
\end{equation}
This is the conventional formulation of Birch--Swinnerton--Dyer. The order,
the leading term and the arithmetic factors are not defined through one
another: they are coordinated publications of the based equality
\eqref{eq:bsd-suc-linea-final}, whose common root was previously fixed by
\eqref{eq:bsd-suc-raiz}.

